% Sensibilisation sur les déséquilibres énergétiques au sein des organismes et leurs conséquences — SVT 1ère D — Cours signé OnBuch+
% Programme officiel MINESEC. Compiler avec : tectonic ND05-desequilibres-energetiques-organismes-d.tex
\newcommand{\DOCMATIERE}{SVT}
\newcommand{\DOCNIVEAU}{1ère D}
\newcommand{\DOCMODULE}{Module 1 : Le Monde Vivant}
\newcommand{\DOCLECON}{ND05}
\newcommand{\DOCTITRE}{Sensibilisation sur les déséquilibres énergétiques au sein des organismes et leurs conséquences}
% OnBuch+ — préambule « cours signés » ENRICHI (compilé avec tectonic / XeLaTeX)
% Système visuel « Pop » d'OnBuch+ : fond crème (jamais blanc), bordures encre
% épaisses, ombres « dures » décalées, titres Archivo Black en capitales.
% Version MATHÉMATIQUES : graphes (pgfplots/tikz), boîtes pédagogiques
% (définition, propriété, méthode, attention, le savais-tu, exercice résolu,
% exercice, TP, bilan), page de garde et signature OnBuch+ ancrée en bas.
%
% Macros attendues dans main.tex AVANT \input{preamble} :
%   \DOCMATIERE  \DOCNIVEAU  \DOCMODULE  \DOCLECON (numéro)  \DOCTITRE
\documentclass[11pt]{article}

\usepackage{fontspec}
\usepackage[french]{babel}
\usepackage[a4paper,margin=2cm,top=3.4cm,bottom=2.8cm,headheight=1.4cm,footskip=1.3cm]{geometry}
\usepackage{amsmath,amssymb,amsfonts}
\usepackage{mathtools} % \xmapsto, \coloneqq... souvent utilisés par les modèles
\usepackage{cancel} % simplifications barrées (bilans de forces)
% \abs{z} (module, valeur absolue) et \norm{u} (norme), utilisés par les modèles
\providecommand{\abs}{}\let\abs\relax\DeclarePairedDelimiter{\abs}{\lvert}{\rvert}
\providecommand{\norm}{}\let\norm\relax\DeclarePairedDelimiter{\norm}{\lVert}{\rVert}
\usepackage[table]{xcolor} % [table] : fournit aussi \rowcolors (alternance de lignes)
\usepackage{pagecolor}
\usepackage{tikz}
\usetikzlibrary{arrows.meta,positioning,calc,shapes.geometric,shapes.misc,shapes.arrows,decorations.pathmorphing,decorations.pathreplacing,decorations.markings,patterns,shadows,mindmap,trees,fit,backgrounds,matrix,intersections,angles,quotes}
\usepackage{pgfplots}
\usepackage[version=4]{mhchem} % équations nucléaires \ce{^{235}_{92}U}
\usepackage[european,siunitx]{circuitikz} % schémas électriques
\pgfplotsset{compat=1.18}
\usepgfplotslibrary{fillbetween,groupplots} % aires entre courbes (calcul intégral)
\usepackage{enumitem}
\usepackage{fancyhdr}
% [most] charge toutes les bibliothèques tcolorbox (listings, minted, raster,
% documentation...) et gonfle inutilement la mémoire TeX sur un document long
% (~300 boîtes) : on ne charge que ce qui est réellement utilisé.
\usepackage[skins,breakable]{tcolorbox}
\usepackage{graphicx}
\usepackage{colortbl}
\usepackage{booktabs}
\usepackage{tabularx}
\usepackage{diagbox} % case d'en-tête diagonale des tableaux à double entrée
% Colonnes extensibles centrée / gauche / droite (C, L, R), souvent utilisées par les modèles
\newcolumntype{C}{>{\centering\arraybackslash}X}
\newcolumntype{L}{>{\raggedright\arraybackslash}X}
\newcolumntype{R}{>{\raggedleft\arraybackslash}X}
\usepackage{array}
\usepackage{multicol}
\usepackage{siunitx}
% parse-numbers=false : accepte aussi \SI{2,5 \times 10^{-3}}{...}
\sisetup{locale=FR,output-decimal-marker={,},group-minimum-digits=4,parse-numbers=false}
\DeclareSIUnit\ohm{\ensuremath{\symup{\Omega}}} % U+2126 absent de la police
\DeclareSIUnit\division{div} % réglages d'oscilloscope (ms/div, V/div)
\DeclareSIUnit\tr{tr} % tours (vitesse de rotation, ex. \SI{1500}{\tr\per\minute})
\DeclareSIUnit\ch{ch} % cheval-vapeur (puissance moteur, unité usuelle hors SI)
\DeclareSIUnit\franccfa{FCFA} % franc CFA (coûts, contexte camerounais)
\DeclareSIUnit\dioptre{\ensuremath{\delta}} % vergence d'une lentille (dioptrie)
\usepackage{qrcode}
% \degree : commande fréquemment utilisée par les modèles pour un angle ou une
% température en dehors de \SI{}{} (non fournie par les paquets chargés).
\providecommand{\degree}{^{\circ}}
% \qed (fin de démonstration), utilisé par les modèles sans amsthm
\providecommand{\qed}{\hfill$\square$}
% \barre : double barre des valeurs interdites dans un tableau de variations (tkz-tab)
\providecommand{\barre}{\ensuremath{\|}}
% environnement proof (amsthm non chargé) utilisé par les modèles
\ifdefined\proof\else\newenvironment{proof}[1][Démonstration]{\par\noindent\textit{#1.}\ }{\qed\par}\fi
\usepackage{lastpage}
\usepackage{hyperref}
\hypersetup{hidelinks}

% ============================================================
% Palette « Pop » OnBuch+ — mêmes hex que la classe P (plus_ui.dart)
% ============================================================
\definecolor{poporange}{HTML}{F59321}
\definecolor{popdark}{HTML}{E07A0C}
\definecolor{popcream}{HTML}{FFF6E8}
\definecolor{popink}{HTML}{1C1714}
\definecolor{poppurple}{HTML}{7A5AE0}
\definecolor{popgreen}{HTML}{2E9E6A}
\definecolor{poppink}{HTML}{E0457B}
\definecolor{popblue}{HTML}{2F7DE1}
\definecolor{popgold}{HTML}{C98A12}
\definecolor{popmuted}{HTML}{8A7F76}
\definecolor{popline}{HTML}{EADFD2}
\definecolor{popgray}{HTML}{8A7F76}
\colorlet{popgrey}{popgray}
% Alias tolérants (le modèle nomme parfois les couleurs différemment)
\colorlet{popwhite}{white}
\colorlet{popblack}{popink}
\colorlet{popred}{poppink}
\colorlet{popyellow}{popgold}
% Teintes claires pour fonds de tableaux / zones
\colorlet{popblueL}{popblue!10!white}
\colorlet{poporangeL}{poporange!14!white}
\colorlet{popgreenL}{popgreen!12!white}
\colorlet{poppurpleL}{poppurple!12!white}
\colorlet{poppinkL}{poppink!12!white}
\colorlet{popgoldL}{popgold!14!white}

\pagecolor{popcream}

% ============================================================
% Typographie
% ============================================================
\defaultfontfeatures{Ligatures=TeX}
\setmainfont{PlusJakartaSans-Medium.ttf}[Path=../../../fonts/,BoldFont=PlusJakartaSans-Bold.ttf]
% Maths en sans-serif (Fira Math) avec chiffres et lettres droites de Plus
% Jakarta, pour que les formules se fondent dans le texte.
\usepackage[math-style=ISO,bold-style=ISO]{unicode-math}
\setmathfont{FiraMath-Regular.otf}[Path=../../../fonts/]
\setmathfont{PlusJakartaSans-Medium.ttf}[Path=../../../fonts/,range={up/num,up/{Latin,latin},\mathup}]
\usepackage{icomma} % virgule décimale sans espace en mode math
% Caractères mathématiques Unicode absents des polices de texte (Plus Jakarta,
% Archivo Black) : rendus par leur équivalent mathématique, en texte comme en
% formule — sinon ils s'affichent en carré vide (ex. « Arithmétique dans ℤ »).
\usepackage{newunicodechar}
\newunicodechar{ℝ}{\ensuremath{\mathbb{R}}}
\newunicodechar{ℤ}{\ensuremath{\mathbb{Z}}}
\newunicodechar{ℂ}{\ensuremath{\mathbb{C}}}
\newunicodechar{ℕ}{\ensuremath{\mathbb{N}}}
\newunicodechar{ℚ}{\ensuremath{\mathbb{Q}}}
\newunicodechar{𝔻}{\ensuremath{\mathbb{D}}}
\newunicodechar{α}{\ensuremath{\alpha}}
\newunicodechar{β}{\ensuremath{\beta}}
\newunicodechar{γ}{\ensuremath{\gamma}}
\newunicodechar{δ}{\ensuremath{\delta}}
\newunicodechar{ε}{\ensuremath{\varepsilon}}
\newunicodechar{θ}{\ensuremath{\theta}}
\newunicodechar{λ}{\ensuremath{\lambda}}
\newunicodechar{μ}{\ensuremath{\mu}}
\newunicodechar{σ}{\ensuremath{\sigma}}
\newunicodechar{τ}{\ensuremath{\tau}}
\newunicodechar{φ}{\ensuremath{\varphi}}
\newunicodechar{ψ}{\ensuremath{\psi}}
\newunicodechar{ω}{\ensuremath{\omega}}
\newunicodechar{Σ}{\ensuremath{\Sigma}}
% majuscules produites par \MakeUppercase dans les titres (α-aminés -> Α)
\newunicodechar{Α}{\ensuremath{\alpha}}
\newunicodechar{Β}{\ensuremath{\beta}}
\newunicodechar{Γ}{\ensuremath{\Gamma}}
\newunicodechar{Δ}{\ensuremath{\Delta}}
\newunicodechar{Ε}{\ensuremath{\varepsilon}}
\newunicodechar{Θ}{\ensuremath{\Theta}}
\newunicodechar{Λ}{\ensuremath{\Lambda}}
\newunicodechar{Μ}{\ensuremath{\mu}}
\newunicodechar{Τ}{\ensuremath{\tau}}
\newunicodechar{Φ}{\ensuremath{\Phi}}
\newunicodechar{Ψ}{\ensuremath{\Psi}}
\newunicodechar{Ω}{\ensuremath{\Omega}}
\newunicodechar{↦}{\ensuremath{\mapsto}}
\newunicodechar{∈}{\ensuremath{\in}}
\newunicodechar{∉}{\ensuremath{\notin}}
\newunicodechar{∧}{\ensuremath{\wedge}}
\newunicodechar{⇔}{\ensuremath{\Leftrightarrow}}
\newunicodechar{⟺}{\ensuremath{\Longleftrightarrow}}
\newunicodechar{⇒}{\ensuremath{\Rightarrow}}
\newunicodechar{⟹}{\ensuremath{\Longrightarrow}}
\newunicodechar{∘}{\ensuremath{\circ}}
\newunicodechar{∪}{\ensuremath{\cup}}
\newunicodechar{∩}{\ensuremath{\cap}}
\newunicodechar{≡}{\ensuremath{\equiv}}
\newunicodechar{⊂}{\ensuremath{\subset}}
\newunicodechar{⊥}{\ensuremath{\perp}}
\newunicodechar{∃}{\ensuremath{\exists}}
\newunicodechar{∀}{\ensuremath{\forall}}
\newunicodechar{∛}{\ensuremath{\sqrt[3]{\,}}}
\newunicodechar{∜}{\ensuremath{\sqrt[4]{\,}}}
\newunicodechar{⃗}{\ensuremath{{}^{\to}}}
\newunicodechar{ⁿ}{\ifmmode^{n}\else\textsuperscript{\ensuremath{n}}\fi}
\newunicodechar{ˣ}{\ifmmode^{x}\else\textsuperscript{\ensuremath{x}}\fi}
\newunicodechar{ᵇ}{\ifmmode^{b}\else\textsuperscript{\ensuremath{b}}\fi}
\newunicodechar{ʳ}{\ifmmode^{r}\else\textsuperscript{\ensuremath{r}}\fi}
\newunicodechar{ᵘ}{\ifmmode^{u}\else\textsuperscript{\ensuremath{u}}\fi}
\newunicodechar{ʸ}{\ifmmode^{y}\else\textsuperscript{\ensuremath{y}}\fi}
\newunicodechar{ᵃ}{\ifmmode^{a}\else\textsuperscript{\ensuremath{a}}\fi}
\newunicodechar{ᵉ}{\ifmmode^{e}\else\textsuperscript{\ensuremath{e}}\fi}
\newunicodechar{ᵏ}{\ifmmode^{k}\else\textsuperscript{\ensuremath{k}}\fi}
\newunicodechar{ᵖ}{\ifmmode^{p}\else\textsuperscript{\ensuremath{p}}\fi}
\newunicodechar{ᴺ}{\ifmmode^{N}\else\textsuperscript{\ensuremath{N}}\fi}
\newunicodechar{ᶜ}{\ifmmode^{c}\else\textsuperscript{\ensuremath{c}}\fi}
\newunicodechar{ʰ}{\ifmmode^{h}\else\textsuperscript{\ensuremath{h}}\fi}
\newunicodechar{ᵠ}{\ifmmode^{q}\else\textsuperscript{\ensuremath{q}}\fi}
\newunicodechar{ₙ}{\ifmmode_{n}\else\textsubscript{\ensuremath{n}}\fi}
\newunicodechar{ₐ}{\ifmmode_{a}\else\textsubscript{\ensuremath{a}}\fi}
\newunicodechar{ᵢ}{\ifmmode_{i}\else\textsubscript{\ensuremath{i}}\fi}
\newunicodechar{ᵣ}{\ifmmode_{r}\else\textsubscript{\ensuremath{r}}\fi}
\newunicodechar{ₖ}{\ifmmode_{k}\else\textsubscript{\ensuremath{k}}\fi}
\newunicodechar{ᵦ}{\ifmmode_{\beta}\else\textsubscript{\ensuremath{\beta}}\fi}
\newunicodechar{ₓ}{\ifmmode_{x}\else\textsubscript{\ensuremath{x}}\fi}
\newfontfamily\popdisplay{ArchivoBlack-Regular.ttf}[Path=../../../fonts/]
\newfontfamily\poplogo{SpaceGrotesk-Bold.ttf}[Path=../../../fonts/]
\newfontfamily\popsemi{PlusJakartaSans-SemiBold.ttf}[Path=../../../fonts/]
\newfontfamily\popxbold{PlusJakartaSans-ExtraBold.ttf}[Path=../../../fonts/]
\newfontfamily\popxboldls{PlusJakartaSans-ExtraBold.ttf}[Path=../../../fonts/,LetterSpace=3.0]

\setlist[enumerate]{leftmargin=1.7em,itemsep=5pt,topsep=4pt}
\setlist[itemize]{leftmargin=1.7em,itemsep=4pt,topsep=4pt,label={\color{poporange}\textbullet}}
\setlength{\parindent}{0pt}
\setlength{\parskip}{5pt}
\renewcommand{\arraystretch}{1.25}

% pgfplots : style Pop par défaut
\pgfplotsset{
  every axis/.append style={axis line style={popink,line width=1pt},tick style={popink},
    grid style={popline,line width=0.5pt},label style={font=\small},tick label style={font=\footnotesize}},
  popcurve/.style={poporange,line width=1.8pt,no markers,smooth},
}
% Même style disponible dans un \draw TikZ simple (hors pgfplots)
\tikzset{popcurve/.style={poporange,line width=1.8pt}}
\tikzset{popgrid/.style={popline,very thin}}
\colorlet{popcurve}{poporange} % le nom du style est parfois utilisé comme couleur

% ============================================================
% Pill / squiggle
% ============================================================
\newcommand{\poppill}[3]{%
  \tikz[baseline=-0.55ex]\node[draw=popink,line width=1.2pt,rounded corners=2.6pt,%
    fill=#2,inner xsep=6.5pt,inner ysep=3pt]{%
    \popxboldls\fontsize{7}{7}\selectfont\color{#3}\MakeUppercase{#1}};%
}
\newcommand{\popsquiggle}[1]{%
  \tikz[baseline=-0.4ex]{
    \draw[#1,line width=2.6pt,line cap=round]
      plot [smooth] coordinates {(0,0.09) (0.34,0.01) (0.68,0.21) (1.02,0.01) (1.36,0.10)};
  }%
}
% Mot-clé surligné (marqueur orange sous le texte)
\newcommand{\cle}[1]{{\popxbold\color{popdark}#1}}

% ============================================================
% En-tête / pied
% ============================================================
\pagestyle{fancy}
\fancyhf{}
\renewcommand{\headrulewidth}{0pt}
\renewcommand{\footrulewidth}{0pt}
\lhead{\raisebox{-0.15ex}{\poplogo\fontsize{13}{13}\selectfont\color{popink}OnBuch\color{poporange}+}}
\rhead{\poppill{\DOCMATIERE\ \textbullet\ \DOCNIVEAU\ \textbullet\ Leçon \DOCLECON}{white}{popink}}
\lfoot{\popxbold\fontsize{7.6}{7.6}\selectfont\color{popmuted}\MakeUppercase{Cours signé OnBuch+}\ \textbullet\ {\popsemi\DOCTITRE}}
\rfoot{\tikz[baseline=-0.6ex]\node[rounded corners=8.5pt,fill=popink,minimum height=17pt,minimum width=17pt,inner xsep=5pt,inner sep=0]{\popxbold\fontsize{8}{8}\selectfont\color{popcream}\thepage\,/\,\pageref*{LastPage}};}

% ============================================================
% Titres
% ============================================================
\newcommand{\chaptitre}[2]{%
  \begin{tcolorbox}[enhanced,colback=poporange,colframe=popink,boxrule=2.6pt,arc=11pt,
      drop shadow=popink,left=15pt,right=15pt,top=12pt,bottom=11pt,before skip=3pt,after skip=18pt]
    \poppill{#2}{popink}{popcream}\par\vspace{9pt}
    {\raggedright\hyphenpenalty=10000\exhyphenpenalty=10000\popdisplay\fontsize{22}{24}\selectfont\color{popcream}\MakeUppercase{#1}\par}
  \end{tcolorbox}
}
% Section numérotée : pastille orange avec le numéro + titre Archivo
\newcounter{popsec}
\newcommand{\coursec}[1]{%
  \stepcounter{popsec}\setcounter{popsub}{0}%
  \par\vspace{16pt}\noindent
  \begin{minipage}{\linewidth}
  \tikz[baseline=-0.7ex]\node[draw=popink,line width=1.6pt,fill=poporange,rounded corners=4pt,minimum size=22pt,inner sep=0]{\popdisplay\fontsize{12}{12}\selectfont\color{popcream}\thepopsec};%
  \hspace{8pt}{\popdisplay\fontsize{15}{17}\selectfont\color{popink}\MakeUppercase{#1}}\par
  \vspace{4pt}\noindent{\color{poporange}\rule{36pt}{3.6pt}}
  \end{minipage}\par\nopagebreak\vspace{8pt}
}
\newcounter{popsub}
\newcommand{\courssub}[1]{%
  \stepcounter{popsub}%
  \par\vspace{9pt}\noindent{\popxbold\fontsize{11.5}{13}\selectfont\color{popink}{\color{poporange}\thepopsec.\thepopsub}\hspace{6pt}#1}\par\nopagebreak\vspace{4pt}
}

% ============================================================
% Boîtes pédagogiques — même recette : fond blanc, bordure épaisse colorée,
% ombre dure encre, badge plein. Titre optionnel [#1].
% ============================================================
\tcbset{popbase/.style={breakable,enhanced,colback=white,boxrule=2.3pt,arc=9pt,
  shadow={3pt}{-3pt}{0pt}{popink},left=14pt,right=14pt,top=11pt,bottom=11pt,before skip=11pt,after skip=15pt,
  fontupper=\popsemi\fontsize{10.6}{14.5}\selectfont\color{popink},
  fontlower=\popsemi\fontsize{10.6}{14.5}\selectfont\color{popink}}}
% Coche dessinée (le glyphe ✓ n'existe pas dans les polices du document)
\newcommand{\popcheck}[1]{\tikz[baseline=-0.5ex]\draw[#1,line width=1.6pt,line cap=round,line join=round](-0.13,0.0)--(-0.04,-0.09)--(0.13,0.1);}
\providecommand{\coche}{\popcheck{popgreen}}
\providecommand{\resultat}[1]{\par\smallskip\noindent\fcolorbox{popgreen}{popgreenL}{\parbox{\dimexpr\linewidth-2\fboxsep-2\fboxrule\relax}{\textbf{Résultat.} #1}}\par\smallskip}
\newcommand{\popboxhead}[3]{\poppill{#1}{#2}{white}\ifx\relax#3\relax\else\hspace{6pt}{\popxbold\fontsize{10.6}{12}\selectfont\color{popink}#3}\fi\par\vspace{7pt}}

\newtcolorbox{aretenir}[1][]{popbase,colframe=popgreen,before upper={\popboxhead{À retenir}{popgreen}{#1}}}
\newtcolorbox{exemplebox}[1][]{popbase,colframe=poporange,before upper={\popboxhead{Exemple}{poporange}{#1}}}
\newtcolorbox{definition}[1][]{popbase,colframe=popblue,before upper={\popboxhead{Définition}{popblue}{#1}}}
\newtcolorbox{propriete}[1][]{popbase,colframe=poppurple,before upper={\popboxhead{Propriété}{poppurple}{#1}}}
\newtcolorbox{methode}[1][]{popbase,colframe=popgold,before upper={\popboxhead{Méthode}{popgold}{#1}}}
\newtcolorbox{attention}[1][]{popbase,colframe=poppink,before upper={\popboxhead{Attention}{poppink}{#1}}}
\newtcolorbox{experience}[1][]{popbase,colframe=popblue,colback=popblueL,before upper={\popboxhead{Expérience}{popblue}{#1}}}
% « Le savais-tu ? » : carte violette pleine (contexte, culture, Cameroun)
\newtcolorbox{savaistu}[1][]{popbase,colback=poppurple,colframe=popink,
  fontupper=\popsemi\fontsize{10.4}{14.2}\selectfont\color{white},
  before upper={\poppill{Le savais-tu\,?}{white}{poppurple}\ifx\relax#1\relax\else\hspace{6pt}{\popxbold\color{white}#1}\fi\par\vspace{7pt}}}
% Exercice résolu : énoncé en haut, \tcblower puis la solution sur fond crème
\newcounter{popexo}\newcounter{popexr}
\newtcolorbox{exoresolu}[1][]{popbase,colframe=poporange,colbacklower=popcream,
  segmentation style={popink,line width=1.2pt,dashed},
  before upper={\stepcounter{popexr}\popboxhead{Exercice résolu \thepopexr}{poporange}{#1}},
  before lower={\poppill{Solution}{popink}{popcream}\par\vspace{6pt}}}
% Exercice (non résolu en place) — la correction est dans la partie Corrigés
\newtcolorbox{exercice}[1][]{popbase,colframe=popink,
  before upper={\stepcounter{popexo}\popboxhead{Exercice \thepopexo}{popink}{#1}}}
% Corrigé
\newtcolorbox{corrige}[1][]{popbase,colframe=popgreen,colback=popgreenL,
  before upper={\popboxhead{Corrigé}{popgreen}{#1}}}
% Objectifs / prérequis / bilan
\newtcolorbox{objectifs}{popbase,colframe=popink,colback=poporangeL,
  before upper={\popboxhead{Objectifs de la leçon}{popink}{}}}
\newtcolorbox{prerequis}{popbase,colframe=popink,colback=popblueL,
  before upper={\popboxhead{Prérequis}{popblue}{}}}
\newtcolorbox{bilan}{popbase,colframe=popink,colback=white,boxrule=2.8pt,
  before upper={\popboxhead{Fiche bilan}{poporange}{L'essentiel à connaître pour l'examen}}}
% Figure encadrée avec légende
\newtcolorbox{popfigure}[1][]{enhanced,breakable=false,colback=white,colframe=popline,boxrule=1.4pt,arc=8pt,
  left=10pt,right=10pt,top=10pt,bottom=8pt,before skip=10pt,after skip=12pt,halign=center,
  lower separated=false,halign lower=center,
  fontlower=\popsemi\fontsize{9}{11.5}\selectfont\color{popmuted},
  #1}
\newcommand{\legende}[1]{\par\vspace{6pt}{\popsemi\fontsize{9}{11.5}\selectfont\color{popmuted}#1}}

% ============================================================
% Signature OnBuch+
% ============================================================
\newcommand{\poplogoinline}[1]{{\poplogo\fontsize{#1}{#1}\selectfont\color{popink}OnBuch\color{poporange}+}}
\newcommand{\signatureonbuch}{%
  \begin{tcolorbox}[enhanced,colback=popink,colframe=popink,boxrule=2.2pt,arc=10pt,
      drop shadow=poporange,left=14pt,right=14pt,top=10pt,bottom=10pt,before skip=0pt,after skip=0pt]
    \tikz[baseline=-0.7ex]\node[circle,fill=popgreen,draw=popcream,line width=1.3pt,minimum size=22pt,inner sep=0]{\popcheck{white}};%
    \hspace{9pt}\begin{minipage}[c]{0.62\linewidth}
      {\popxbold\fontsize{12}{13}\selectfont\color{popcream}Signé\ \textbullet\ {\poplogo OnBuch\color{poporange}+}}\par\vspace{2pt}
      {\raggedright\popsemi\fontsize{8}{10.5}\selectfont\color{popline}\MakeUppercase{Équipe pédagogique OnBuch+}\\\MakeUppercase{Conforme au programme officiel MINESEC}\par}
    \end{minipage}\hfill
    {\poplogo\fontsize{22}{22}\selectfont\color{popcream}OnBuch\color{poporange}+}
  \end{tcolorbox}%
}

% ============================================================
% Page de garde
% #1 titre · #2 matière · #3 niveau · #4 module · #5 n° leçon · #6 durée ·
% #7 description / objectifs (texte ou itemize)
% ============================================================
\newcommand{\pagegarde}[7]{%
  \thispagestyle{empty}
  \newgeometry{a4paper,margin=1.7cm,top=1.6cm,bottom=1.4cm}%
  \begin{tikzpicture}[remember picture,overlay]
    \fill[poporange!18] ([xshift=-3.2cm,yshift=-3.6cm]current page.north east) circle (4.6cm);
    \fill[poppurple!14] ([xshift=2.4cm,yshift=7.6cm]current page.south west) circle (3.8cm);
  \end{tikzpicture}%
  {\poplogo\fontsize{30}{30}\selectfont\color{popink}OnBuch\color{poporange}+}\hfill
  \raisebox{0.6ex}{\poppill{Cours signé \textbullet\ Premium}{popink}{popcream}}\par
  \vspace{1pt}\popsquiggle{poppurple}\par
  \vspace{8pt}
  \begin{tcolorbox}[enhanced,colback=poporange,colframe=popink,boxrule=2.8pt,arc=13pt,
      drop shadow=popink,left=18pt,right=18pt,top=12pt,bottom=13pt,after skip=10pt]
    \poppill{#2 \textbullet\ #3}{popink}{popcream}\hspace{5pt}\poppill{#4}{white}{popink}\par\vspace{8pt}
    {\popdisplay\fontsize{14}{14}\selectfont\color{popink}LEÇON #5}\par\vspace{4pt}
    {\raggedright\hyphenpenalty=10000\exhyphenpenalty=10000\popdisplay\fontsize{25}{27}\selectfont\color{popcream}\MakeUppercase{#1}\par}
    \vspace{8pt}
    {\popxbold\fontsize{9.5}{9.5}\selectfont\color{popink}\MakeUppercase{Durée conseillée : #6}}
  \end{tcolorbox}
  \begin{tcolorbox}[enhanced,colback=white,colframe=popink,boxrule=2.2pt,arc=10pt,
      drop shadow=popink,left=16pt,right=16pt,top=10pt,bottom=10pt,after skip=8pt]
    \poppill{Dans cette leçon}{poppurple}{white}\par\vspace{5pt}
    \popsemi\fontsize{9.3}{12.6}\selectfont\color{popink}#7
  \end{tcolorbox}
  \noindent\begin{minipage}[t]{0.487\linewidth}
    \begin{tcolorbox}[enhanced,colback=popgreen,colframe=popink,boxrule=2.2pt,arc=10pt,
        drop shadow=popink,left=11pt,right=11pt,top=10pt,bottom=10pt,height=3.1cm,valign=center]
      \tcbox[colback=white,colframe=popink,boxrule=1.3pt,arc=5pt,boxsep=0pt,nobeforeafter,
        left=3pt,right=3pt,top=3pt,bottom=3pt,valign=center]{\qrcode[height=1.9cm]{https://play.google.com/store/apps/details?id=com.luvvix.onbuch&hl=fr}}%
      \hspace{8pt}\begin{minipage}[c]{0.5\linewidth}
        {\popdisplay\fontsize{11}{12.5}\selectfont\color{white}\MakeUppercase{Télécharge OnBuch}}\par\vspace{4pt}
        {\raggedright\popsemi\fontsize{8.4}{11}\selectfont\color{white}Gratuit \textbullet\ Google Play\\Tuteur IA, annales, concours, résultats\par}
      \end{minipage}
    \end{tcolorbox}
  \end{minipage}\hfill
  \begin{minipage}[t]{0.487\linewidth}
    \begin{tcolorbox}[enhanced,colback=poppurple,colframe=popink,boxrule=2.2pt,arc=10pt,
        drop shadow=popink,left=11pt,right=11pt,top=10pt,bottom=10pt,height=3.1cm,valign=center]
      \tikz[baseline=-0.7ex]\node[circle,fill=white,draw=popink,line width=1.3pt,minimum size=1.6cm,inner sep=0]{\popdisplay\fontsize{24}{24}\selectfont\color{poppurple}+};%
      \hspace{8pt}\begin{minipage}[c]{0.6\linewidth}
        {\popdisplay\fontsize{11}{12.5}\selectfont\color{white}\MakeUppercase{Passe à OnBuch+}}\par\vspace{4pt}
        {\raggedright\popsemi\fontsize{8.4}{11}\selectfont\color{white}Tous les cours signés, Tuteur IA illimité, fiches PDF — onglet OnBuch+ de l'app\par}
      \end{minipage}
    \end{tcolorbox}
  \end{minipage}
  \vspace{8pt}
  \popcontenu
  \vfill
  \signatureonbuch
  \restoregeometry
  \newpage
}

% Rangée « Ce cours contient » (page de garde)
\newcommand{\poptile}[3]{% #1 couleur #2 gros texte #3 légende
  \begin{tcolorbox}[enhanced,colback=#1,colframe=popink,boxrule=2pt,arc=9pt,shadow={2.5pt}{-2.5pt}{0pt}{popink},
      left=6pt,right=6pt,top=9pt,bottom=9pt,halign=center,valign=center,height=1.75cm,width=\linewidth,nobeforeafter]
    {\popdisplay\fontsize{12.5}{13}\selectfont\color{white}\MakeUppercase{#2}}\par\vspace{3pt}
    {\centering\popsemi\fontsize{7.8}{9.5}\selectfont\color{white}#3\par}
  \end{tcolorbox}}
\newcommand{\popcontenu}{%
  \par\noindent{\popxboldls\fontsize{7.5}{7.5}\selectfont\color{popmuted}CE COURS SIGNÉ CONTIENT}\par\vspace{7pt}
  \noindent\begin{minipage}[t]{0.235\linewidth}\poptile{popblue}{Cours}{complet, illustré, pas à pas}\end{minipage}\hfill
  \begin{minipage}[t]{0.235\linewidth}\poptile{poporange}{Méthodes}{et exercices résolus}\end{minipage}\hfill
  \begin{minipage}[t]{0.235\linewidth}\poptile{poppink}{Exercices}{type examen + corrigés}\end{minipage}\hfill
  \begin{minipage}[t]{0.235\linewidth}\poptile{popgreen}{Fiche bilan}{l'essentiel à retenir}\end{minipage}\par}

% Sommaire
\newtcolorbox{sommaire}{enhanced,colback=white,colframe=popink,boxrule=2.2pt,arc=10pt,
  drop shadow=popink,left=18pt,right=18pt,top=13pt,bottom=13pt,before skip=6pt,after skip=20pt,
  before upper={\poppill{Sommaire}{poppurple}{white}\par\vspace{9pt}\popsemi\fontsize{10.6}{14.5}\selectfont\color{popink}}}

% Signature de fin de cours — poussée tout en bas de la dernière page
\newcommand{\findecours}{%
  \par\vfill\nopagebreak
  \begin{center}\popsquiggle{poporange}\end{center}\vspace{4pt}
  \signatureonbuch
}
\AtBeginDocument{\def\llbracket{\mathopen{[\mkern-3.5mu[}}\def\rrbracket{\mathclose{]\mkern-3.5mu]}}} % intervalles d'entiers (glyphe ⟦ absent de la police)
% Glyphes absents de Fira Math : redéfinis à partir de glyphes disponibles
\AtBeginDocument{%
  \def\setminus{\mathbin{\backslash}}\let\smallsetminus\setminus
  \def\vdots{\mathord{\vbox{\baselineskip3.2pt\lineskiplimit0pt\kern4pt\hbox{.}\hbox{.}\hbox{.}}}}%
  \def\ddots{\mathinner{\mkern1mu\raise6.5pt\hbox{.}\mkern2mu\raise3.6pt\hbox{.}\mkern2mu\raise0.7pt\hbox{.}\mkern1mu}}%
  \def\triangle{\mathord{\scalebox{0.9}{$\symup{\Delta}$}}}%
  \def\nabla{\mathord{\raisebox{\depth}{\scalebox{1}[-1]{$\symup{\Delta}$}}}}%
  \def\checkmark{\popcheck{popdark}}%
  \def\star{\mathbin{\ast}}\def\bigstar{\mathord{\ast}}%
  \def\diamond{\mathbin{\scalebox{0.6}{$\Diamond$}}}%
  \def\bigcup{\mathop{\mathchoice{\raisebox{-0.3em}{\scalebox{1.7}{$\cup$}}}{\scalebox{1.25}{$\cup$}}{\cup}{\cup}}\displaylimits}%
  \def\bigcap{\mathop{\mathchoice{\raisebox{-0.3em}{\scalebox{1.7}{$\cap$}}}{\scalebox{1.25}{$\cap$}}{\cap}{\cap}}\displaylimits}%
  \def\lesseqqgtr{\mathrel{\lessgtr}}\def\bigoplus{\mathop{\oplus}\displaylimits}%
}
\providecommand{\ssi}{si et seulement si} % abréviation fréquente
\AtBeginDocument{\providecommand{\wideparen}{\overparen}} % arc de cercle

%% FIGFIX-BEGIN v3 — correctifs globaux des figures (docs/onbuch-generation/tools/figfix)
\usepackage{adjustbox}\usepackage{etoolbox}
% toute figure TikZ/pgfplots plus large que la ligne est réduite à la largeur disponible
\BeforeBeginEnvironment{tikzpicture}{\begin{adjustbox}{max width=\linewidth}}
\AfterEndEnvironment{tikzpicture}{\end{adjustbox}}
% légendes pgfplots lisibles par-dessus les courbes
\pgfplotsset{every axis/.append style={legend style={fill=white,fill opacity=0.92,text opacity=1,draw=black!15,inner sep=3pt}}}
% étiquettes d'arêtes (syntaxe edge["..."]) avec fond blanc
\tikzset{every edge quotes/.append style={fill=white,inner sep=1.2pt}}
% bulles de cartes mentales : texte à la ligne dans la bulle, bulle un peu plus grande
\tikzset{every concept/.append style={text width=2.0cm,align=center,minimum size=2.3cm,inner sep=2pt}}
%% FIGFIX-END
\begin{document}

\pagegarde{Sensibilisation sur les déséquilibres énergétiques au sein des organismes et leurs conséquences}{SVT}{1ère D}{Module 1 : Le Monde Vivant}{ND05}{22 h}{Cette leçon aborde les sources de l'énergie cellulaire, la construction d'un bilan énergétique, les conséquences physiologiques des déséquilibres et les adaptations métaboliques à l'effort et au jeûne, en vue de proposer des mesures correctives adaptées au contexte camerounais.
\vspace{4pt}
\begin{multicols}{2}\small
\begin{itemize}
\item À la fin de cette leçon, je sais décrire les voies de production d'ATP (glycolyse, respiration aérobie, fermentations) et comparer leurs rendements.
\item Je sais calculer les apports énergétiques quotidiens à partir de la composition d'un repas typique camerounais.
\item Je sais estimer les dépenses énergétiques (métabolisme de base, activité physique, thermogenèse) et établir un bilan énergétique quantitatif.
\item Je sais interpréter des courbes de consommation d'O₂, de production de CO₂ ou de variation pondérale pour détecter un déséquilibre.
\item Je sais relier un déséquilibre énergétique positif ou négatif à ses conséquences physiologiques (stockage lipidique, résistance à l'insuline, catabolisme protéique, troubles immunitaires).
\item Je sais expliquer les adaptations métaboliques lors d'un effort physique (utilisation des substrats, rôle hormonal) et lors d'une privation alimentaire (glycogénolyse, néoglucogenèse, cétogenèse).
\end{itemize}\end{multicols}}

\begin{sommaire}
\begin{multicols}{2}
\begin{enumerate}
\item 1. Sources et voies de production de l'énergie cellulaire
\item 2. Bilan énergétique de l'organisme : apports et dépenses
\item 3. Conséquences physiologiques des déséquilibres énergétiques
\item 4. Adaptations métaboliques à l'effort et à la privation alimentaire
\item 5. Mesures correctives, prévention et éducation nutritionnelle
\item Activité d'intégration
\item Méthodes et exercices résolus
\item Exercices
\item Corrigés des exercices
\item Fiche bilan
\end{enumerate}
\end{multicols}
\end{sommaire}

\begin{objectifs}
\begin{itemize}
\item À la fin de cette leçon, je sais décrire les voies de production d'ATP (glycolyse, respiration aérobie, fermentations) et comparer leurs rendements.
\item Je sais calculer les apports énergétiques quotidiens à partir de la composition d'un repas typique camerounais.
\item Je sais estimer les dépenses énergétiques (métabolisme de base, activité physique, thermogenèse) et établir un bilan énergétique quantitatif.
\item Je sais interpréter des courbes de consommation d'O₂, de production de CO₂ ou de variation pondérale pour détecter un déséquilibre.
\item Je sais relier un déséquilibre énergétique positif ou négatif à ses conséquences physiologiques (stockage lipidique, résistance à l'insuline, catabolisme protéique, troubles immunitaires).
\item Je sais expliquer les adaptations métaboliques lors d'un effort physique (utilisation des substrats, rôle hormonal) et lors d'une privation alimentaire (glycogénolyse, néoglucogenèse, cétogenèse).
\item Je sais proposer des mesures correctives nutritionnelles et d'activité physique adaptées à un cas concret.
\item Je sais communiquer mes résultats sous forme de tableaux, graphiques et courts rapports argumentés.
\end{itemize}
\end{objectifs}

\begin{prerequis}
\begin{itemize}
\item Notion de macronutriments (glucides, lipides, protéines) et de leur valeur énergétique (kJ/g).
\item Équation globale de la respiration cellulaire et de la fermentation alcoolique/lactique (Seconde).
\item Concept de métabolisme de base et de facteur d'activité physique (Première).
\item Lecture de graphiques (courbes, histogrammes) et calculs de pourcentages.
\item Notions de base sur les hormones régulatrices (insuline, glucagon, adrénaline).
\end{itemize}
\end{prerequis}

\newpage

\coursec{Sources et voies de production de l'énergie cellulaire}

\textbf{Situation d'accroche.} Dans le quartier de Nkolbisson, à Yaoundé, Marie, 16 ans, mange chaque midi un généreux plat de ndolé accompagné de plantains frits, arrosé d'une boisson sucrée. Après les cours, elle s'installe devant la télévision pendant trois heures. Six mois plus tard, son médecin constate une prise de poids de 4 kg et Marie se plaint d'une fatigue inhabituelle dès qu'elle monte les escaliers. D'où vient cette fatigue ? Pourquoi le corps stocke-t-il ? Tout commence dans les cellules : chaque bouchée de plantain contient du glucose, et chaque cellule doit transformer ce glucose en une énergie utilisable, l'\cle{ATP}. Comprendre comment la cellule produit cette énergie — avec ou sans oxygène — est la première étape pour expliquer le déséquilibre de Marie.

\textbf{Problématique.} Comment la cellule produit-elle de l'énergie à partir des nutriments ? Quelles sont les voies possibles, leurs conditions et leurs rendements respectifs ?

\courssub{La glycolyse : la voie commune à toutes les cellules}

\textbf{Observation.} Une levure de boulanger placée dans du jus de fruit sucré, même enfermée dans un bocal sans air, continue à vivre et produit des bulles de gaz. Un muscle qui contracte très fort, même quand l'apport en dioxygène ne suffit plus, continue un certain temps à fonctionner. Ces cellules tirent donc de l'énergie du glucose \emph{sans oxygène}. L'étape commune à toutes ces situations est la \cle{glycolyse}.

\begin{definition}[Glycolyse]
La \cle{glycolyse} est une suite de dix réactions enzymatiques se déroulant dans le \textbf{cytoplasme} (cytosol) de toute cellule, au cours desquelles une molécule de \textbf{glucose} (6 atomes de carbone) est scindée en deux molécules de \textbf{pyruvate} (3 atomes de carbone). Elle ne consomme \textbf{pas de dioxygène}.
\end{definition}

\begin{propriete}[Bilan de la glycolyse]
Pour une molécule de glucose dégradée, la glycolyse produit :
\begin{itemize}
\item \textbf{2 molécules de pyruvate} ;
\item un gain net de \textbf{2 ATP} (4 ATP produits moins 2 ATP investis au départ) ;
\item \textbf{2 NADH, H$^+$} (coenzymes réduits, transporteurs d'électrons).
\end{itemize}
\[\ce{C6H12O6 + 2 ADP + 2 Pi + 2 NAD+ -> 2 C3H4O3 + 2 ATP + 2 NADH,H+ + 2 H2O}\]
\end{propriete}

\begin{attention}[Piège fréquent]
Le bilan net de la glycolyse est de \textbf{2 ATP}, pas 4. On oublie souvent que 2 ATP sont \emph{consommés} dans les premières étapes (phosphorylation du glucose et du fructose-6-phosphate). Au Probatoire, écris toujours « gain net : 2 ATP ».
\end{attention}

Le pyruvate formé se trouve alors à un \textbf{carrefour métabolique} : son destin dépend de la présence ou de l'absence de dioxygène.

\courssub{La respiration cellulaire aérobie : le rendement maximal}

\textbf{Observation.} Un athlète au repos respire calmement : ses cellules musculaires reçoivent assez de \ce{O2}. Si l'on mesure sa consommation de \ce{O2} et son rejet de \ce{CO2}, on constate qu'ils sont proportionnels à l'oxydation du glucose. C'est la signature de la \cle{respiration cellulaire}.

\begin{definition}[Respiration cellulaire aérobie]
La \cle{respiration cellulaire} est l'ensemble des réactions d'\textbf{oxydation} qui, en \textbf{présence de dioxygène}, dégradent complètement le pyruvate (issu du glucose) jusqu'à \ce{CO2} et \ce{H2O}, avec production abondante d'\textbf{ATP}. Elle se déroule dans la \textbf{mitochondrie} (matrice et crêtes mitochondriales).
\end{definition}

Elle comprend trois étapes :

\textbf{1. La décarboxylation oxydative du pyruvate.} Dans la matrice mitochondriale, chaque pyruvate (3 C) perd un carbone sous forme de \ce{CO2} et est transformé en \textbf{acétyl-CoA} (2 C), avec formation d'un NADH, H$^+$. Pour un glucose : 2 acétyl-CoA, 2 \ce{CO2}, 2 NADH.

\textbf{2. Le cycle de Krebs} (cycle de l'acide citrique). Dans la matrice, l'acétyl-CoA entre dans un cycle de réactions qui oxyde complètement ses 2 carbones en \ce{CO2}. Par tour de cycle : 3 NADH, 1 \ce{FADH2}, 1 ATP (via GTP) et 2 \ce{CO2}. Comme un glucose fournit 2 acétyl-CoA, le cycle tourne \textbf{deux fois} : au total 6 NADH, 2 \ce{FADH2}, 2 ATP, 4 \ce{CO2}.

\textbf{3. La chaîne de transport d'électrons (CTE) et la phosphorylation oxydative.} Sur les crêtes mitochondriales (membrane interne), les NADH et \ce{FADH2} cèdent leurs électrons à une chaîne de transporteurs. L'énergie libérée pompe des protons \ce{H+} vers l'espace intermembranaire ; leur retour dans la matrice à travers l'\textbf{ATP synthase} entraîne la synthèse massive d'ATP. Le dioxygène est l'\textbf{accepteur final d'électrons} : il est réduit en eau. C'est lui qui « aspire » les électrons et permet à toute la chaîne de fonctionner.

\begin{propriete}[Rendement maximal théorique de la respiration (admis)]
L'oxydation complète d'une molécule de glucose par la respiration aérobie produit au maximum \textbf{environ 38 ATP} :
\begin{itemize}
\item 2 ATP nets (glycolyse) + 2 ATP (cycle de Krebs) = 4 ATP directs (phosphorylation au niveau du substrat) ;
\item environ 34 ATP issus de la chaîne de transport d'électrons (revalorisation des 10 NADH et 2 \ce{FADH2}).
\end{itemize}
Équation globale : \[\ce{C6H12O6 + 6 O2 -> 6 CO2 + 6 H2O} + \text{énergie } (\approx 38\ \text{ATP})\]
Ce résultat est admis ; la méthode de calcul à partir des NADH/\ce{FADH2} est, elle, exigible.
\end{propriete}

\begin{methode}[Calculer le rendement ATP à partir des NADH et \ce{FADH2}]
Retiens les équivalences actuelles (rapports P/O) : \textbf{1 NADH $\approx$ 2,5 ATP} et \textbf{1 \ce{FADH2} $\approx$ 1,5 ATP}.
\begin{enumerate}
\item Compte les NADH produits : 2 (glycolyse) + 2 (décarboxylation) + 6 (Krebs) = \textbf{10 NADH}.
\item Compte les \ce{FADH2} : 2 (Krebs).
\item Calcule : $10 \times 2{,}5 + 2 \times 1{,}5 = 25 + 3 = 28$ ATP par la CTE (les NADH cytosoliques de la glycolyse valent parfois moins selon la navette utilisée — navette malate-aspartate ou glycérol-3-phosphate —, d'où des totaux de 30 à 32 ATP selon les cellules).
\item Ajoute les 4 ATP directs (glycolyse + Krebs). Total : environ 32 à 38 ATP. La valeur théorique maximale retenue au lycée est \textbf{38 ATP}.
\end{enumerate}
\end{methode}

\begin{exemplebox}[Valeur énergétique du glucose]
1 mole de glucose (180 g) oxydée libère environ \SI{2870}{kJ} ; la synthèse de 38 mol d'ATP en récupère environ $38 \times 30{,}5 \approx \SI{1159}{kJ}$, soit un rendement proche de 40 \% (le reste est dissipé en chaleur, ce qui maintient notre température à 37 \textdegree C). En valeur physiologique, \textbf{1 g de glucose $\approx$ 17 kJ} (environ 4 kcal).
\end{exemplebox}

\courssub{Les fermentations : produire de l'ATP sans oxygène}

\textbf{Observation.} Lors d'un 100 m, les muscles de la cuisse travaillent si intensément que l'apport en \ce{O2} ne suit plus : une douleur de « brûlure » apparaît. De même, la pâte à beignets (puff-puff) laissée au repos « lève » et dégage une odeur d'alcool. Dans les deux cas, la glycolyse continue seule : c'est la \cle{fermentation}.

\begin{definition}[Fermentation]
La \cle{fermentation} est une voie de production d'énergie \textbf{sans dioxygène} (anaérobie) et \textbf{sans chaîne de transport d'électrons}, qui se déroule dans le cytoplasme. Son rôle essentiel est de \textbf{régénérer le \ce{NAD+}} à partir du NADH, permettant à la glycolyse de continuer. Elle ne produit \textbf{aucun ATP supplémentaire} : le rendement reste celui de la glycolyse, soit \textbf{2 ATP par glucose}.
\end{definition}

\textbf{Fermentation lactique} (cellules musculaires en anaérobiose, certaines bactéries lactiques) : le pyruvate est réduit en \textbf{lactate} par le NADH, qui redevient \ce{NAD+}.
\[\ce{C3H4O3 + NADH,H+ -> C3H6O3 + NAD+}\]
L'accumulation de lactate est associée à la fatigue musculaire ; au repos, le lactate est reconverti (foie, cycle de Cori) en pyruvate puis en glucose (néoglucogenèse).

\textbf{Fermentation alcoolique} (levures) : le pyruvate est d'abord décarboxylé en éthanal (avec dégagement de \ce{CO2}, qui fait lever la pâte), puis réduit en \textbf{éthanol}.
\[\ce{C3H4O3 -> C2H4O + CO2} \qquad \ce{C2H4O + NADH,H+ -> C2H6O + NAD+}\]

\begin{attention}[Fermentation $\neq$ respiration anaérobie]
Ne confonds jamais :
\begin{itemize}
\item \textbf{Fermentation} : pas de chaîne de transport d'électrons, \textbf{aucun} accepteur d'électrons externe, aucun \ce{O2} consommé, rendement 2 ATP.
\item \textbf{Respiration anaérobie} (certains micro-organismes) : une CTE fonctionne, mais l'accepteur final d'électrons n'est \textbf{pas} le \ce{O2} (c'est par exemple le nitrate \ce{NO3-} ou le sulfate \ce{SO4^2-}).
\end{itemize}
Au lycée, quand on dit « sans oxygène » pour une cellule musculaire ou une levure, il s'agit de \textbf{fermentation}.
\end{attention}

\begin{popfigure}
\begin{center}
\begin{tikzpicture}[font=\small, node distance=0.55cm and 1.1cm,
  mol/.style={draw=popdark, fill=popcream, rounded corners=2pt, minimum height=0.8cm, inner sep=4pt},
  etape/.style={font=\footnotesize\itshape, text=popmuted},
  fleche/.style={-{Stealth[length=2.5mm]}, thick, popdark}]
\node[mol, fill=popgoldL] (glu) {Glucose (6 C)};
\node[mol, right=2.2cm of glu] (pyr) {2 Pyruvate (3 C)};
\draw[fleche] (glu) -- (pyr) node[midway, above, etape, align=center]{Glycolyse (cytoplasme)\\ 2 ATP + 2 NADH};
\node[mol, above right=0.9cm and 1.6cm of pyr, fill=popgreenL] (ace) {2 Acétyl-CoA (2 C)};
\draw[fleche, popgreen] (pyr) -- (ace) node[midway, above, etape, align=center]{Aérobie : présence de \ce{O2}\\ 2 NADH + 2 \ce{CO2}};
\node[mol, right=1.4cm of ace, fill=popgreenL] (krebs) {Cycle de Krebs};
\draw[fleche, popgreen] (ace) -- (krebs) node[midway, above, etape, align=center]{2 ATP + 6 NADH\\ + 2 \ce{FADH2} + 4 \ce{CO2}};
\node[mol, right=1.3cm of krebs, fill=popblueL] (cte) {CTE + \ce{O2}};
\draw[fleche, popblue] (krebs) -- (cte) node[midway, above, etape, align=center]{$\approx$ 34 ATP\\ \ce{O2 -> H2O}};
\node[mol, below right=0.9cm and 1.6cm of pyr, fill=poppinkL] (lac) {2 Lactate (3 C)};
\draw[fleche, poppink] (pyr) -- (lac) node[midway, below, etape, align=center]{Anaérobie (muscle)\\ fermentation lactique\\ 0 ATP, \ce{NAD+} régénéré};
\node[mol, below=1.5cm of lac, fill=poporangeL] (eth) {2 Éthanol + 2 \ce{CO2}};
\draw[fleche, poporange] (pyr) |- (eth) node[pos=0.75, below, etape, align=center]{Anaérobie (levure)\\ fermentation alcoolique\\ 0 ATP, \ce{NAD+} régénéré};
\end{tikzpicture}
\end{center}
\legende{Les destins du pyruvate. La glycolyse (voie commune, cytoplasme) produit 2 ATP nets. En présence de \ce{O2} (vert, bleu), le pyruvate entre dans la mitochondrie : décarboxylation, cycle de Krebs puis chaîne de transport d'électrons (rendement total $\approx$ 38 ATP). En l'absence de \ce{O2} (rose, orange), les fermentations régénèrent le \ce{NAD+} sans produire d'ATP supplémentaire (rendement total : 2 ATP).}
\end{popfigure}

\courssub{Comparaison des rendements énergétiques}

\begin{popfigure}
\begin{center}
\begin{tikzpicture}
\begin{axis}[
  ybar, bar width=0.9cm,
  width=0.85\linewidth, height=6.2cm,
  ymin=0, ymax=44,
  ylabel={ATP produits par glucose},
  symbolic x coords={Glycolyse seule, Fermentation, Respiration aérobie},
  xtick=data, x tick label style={font=\small},
  ytick={0,2,10,20,30,38},
  nodes near coords, every node near coord/.append style={font=\small\bfseries},
  axis lines=left, enlarge x limits=0.3,
  grid=major, grid style={popline!40},
]
\addplot[fill=popgold, draw=popdark] coordinates {(Glycolyse seule,2) (Fermentation,2) (Respiration aérobie,38)};
\end{axis}
\end{tikzpicture}
\end{center}
\legende{Rendement énergétique comparé des trois voies. La fermentation n'ajoute aucun ATP à la glycolyse (2 ATP), tandis que la respiration aérobie multiplie le rendement par 19 ($\approx$ 38 ATP), grâce à la chaîne de transport d'électrons.}
\end{popfigure}

\begin{center}
\begin{tabularx}{\linewidth}{|l|X|X|X|}
\hline
\rowcolor{popblueL} \textbf{Critère} & \textbf{Respiration aérobie} & \textbf{Fermentation lactique} & \textbf{Fermentation alcoolique} \\
\hline
Localisation & Cytoplasme + mitochondrie & Cytoplasme & Cytoplasme \\
\hline
\ce{O2} requis & Oui (accepteur final) & Non & Non \\
\hline
Produits finaux & \ce{CO2} + \ce{H2O} & Lactate & Éthanol + \ce{CO2} \\
\hline
ATP par glucose & $\approx$ 38 & 2 & 2 \\
\hline
Cellules concernées & Toutes cellules aérobies & Muscle en effort intense, bactéries lactiques & Levures \\
\hline
\end{tabularx}
\end{center}

\courssub{Les substrats alternatifs : acides gras et acides aminés}

\textbf{Observation.} Après plusieurs heures de jeûne, ou lors d'un régime, l'organisme continue à produire de l'ATP alors que les réserves de glucose s'épuisent. D'où vient l'énergie ? Des \textbf{graisses} et, en dernier recours, des \textbf{protéines}.

\begin{definition}[Substrats énergétiques alternatifs]
Outre le glucose, la cellule peut produire de l'\textbf{acétyl-CoA} — la molécule d'entrée du cycle de Krebs — à partir :
\begin{itemize}
\item des \textbf{acides gras}, par \cle{$\beta$-oxydation} : dans la matrice mitochondriale, la chaîne carbonée des acides gras est raccourcie de 2 carbones à chaque cycle, libérant des acétyl-CoA, des NADH et des \ce{FADH2}. Le rendement énergétique des lipides est très élevé : \textbf{1 g de lipide $\approx$ 38 kJ} (environ 9 kcal), soit plus du double des glucides ;
\item des \textbf{acides aminés}, par \cle{déamination} : le groupement aminé (\ce{-NH2}) est retiré (il sera éliminé sous forme d'\textbf{urée} par le foie et les reins), et le squelette carboné restant entre dans la glycolyse ou le cycle de Krebs.
\end{itemize}
\end{definition}

\begin{exoresolu}[Rendement de la respiration à partir des coenzymes]
\textbf{Énoncé.} L'oxydation complète d'une molécule de glucose produit 10 NADH et 2 \ce{FADH2} réoxydés par la chaîne de transport d'électrons, plus 4 ATP directs. En admettant que 1 NADH $\approx$ 2,5 ATP et 1 \ce{FADH2} $\approx$ 1,5 ATP, calcule le nombre total théorique d'ATP, puis compare à une fermentation lactique.
\tcblower
\textbf{Solution.}
\begin{align*}
\text{ATP (CTE)} &= 10 \times 2{,}5 + 2 \times 1{,}5 = 25 + 3 = 28 \text{ ATP}\\
\text{ATP totaux} &= 28 + 4 = 32 \text{ ATP}
\end{align*}
(les NADH cytosoliques de la glycolyse valent parfois davantage selon la navette employée, ce qui conduit au maximum théorique de \textbf{38 ATP} retenu au programme). En fermentation lactique : seule la glycolyse fonctionne, soit \textbf{2 ATP}. La respiration est donc \textbf{16 à 19 fois plus rentable}. C'est pourquoi un effort prolongé ne peut être alimenté que par la voie aérobie.
\end{exoresolu}

\begin{savaistu}[Un cœur qui brûle des graisses]
Au repos, le \textbf{muscle cardiaque} — qui ne s'arrête jamais — tire 60 à 70 \% de son énergie de la \textbf{$\beta$-oxydation des acides gras}, et non du glucose. Riche en mitochondries et irrigué en permanence, il fonctionne presque exclusivement en aérobiose. C'est pourquoi une obstruction des coronaires (infarctus), qui prive le cœur d'\ce{O2}, est si rapidement dramatique : sans \ce{O2}, la CTE s'arrête et le rendement s'effondre de 38 à 2 ATP par glucose.
\end{savaistu}

\begin{aretenir}[L'essentiel de la section]
\begin{itemize}
\item La \textbf{glycolyse} (cytoplasme, sans \ce{O2}) transforme 1 glucose en 2 pyruvate : gain net \textbf{2 ATP} + 2 NADH.
\item En présence de \ce{O2}, la \textbf{respiration} (décarboxylation + cycle de Krebs + chaîne de transport d'électrons, dans la mitochondrie) porte le rendement à \textbf{$\approx$ 38 ATP} ; le \ce{O2} est l'accepteur final d'électrons.
\item En l'absence de \ce{O2}, les \textbf{fermentations} (lactique chez le muscle, alcoolique chez la levure) régénèrent le \ce{NAD+} \textbf{sans produire d'ATP supplémentaire}.
\item Les \textbf{acides gras} ($\beta$-oxydation) et les \textbf{acides aminés} (déamination) alimentent le cycle de Krebs via l'acétyl-CoA : ce sont les carburants de secours de l'organisme.
\end{itemize}
\end{aretenir}

\begin{exercice}[Pour vérifier ta compréhension]
\begin{enumerate}
\item Pourquoi la fermentation est-elle indispensable à la poursuite de la glycolyse en anaérobiose, alors qu'elle ne produit aucun ATP ?
\item Un coureur de marathon puise son énergie principalement dans ses réserves lipidiques après la 25\textsuperscript{e} kilomètre. Justifie ce choix métabolique par deux arguments.
\item Calcule l'énergie (en kJ) libérée par l'oxydation complète de 10 g de glucose, puis la fraction stockée sous forme d'ATP si le rendement est de 40 \%.
\end{enumerate}
\end{exercice}

\begin{corrige}[Exercice n]
\begin{enumerate}
\item La glycolyse consomme du \ce{NAD+} (réduit en NADH). Sans fermentation, tout le \ce{NAD+} serait bloqué sous forme de NADH et la glycolyse s'arrêterait. En réduisant le pyruvate (en lactate ou en éthanol), la fermentation \textbf{régénère le \ce{NAD+}}, permettant à la glycolyse — seule source d'ATP en anaérobiose — de continuer.
\item (a) Les lipides ont une densité énergétique élevée ($\approx$ 38 kJ/g, plus du double des glucides) : ils constituent une réserve massive et légère. (b) Les réserves de glycogène (glucose stocké) sont limitées et s'épuisent après environ 1 h 30 à 2 h d'effort ; la $\beta$-oxydation des acides gras prend alors le relais, en aérobiose, pour alimenter le cycle de Krebs.
\item Énergie libérée : $10 \times 17 = \SI{170}{kJ}$. Fraction stockée en ATP : $170 \times 0{,}40 = \SI{68}{kJ}$ ; le reste ($\SI{102}{kJ}$) est dissipé en chaleur.
\end{enumerate}
\end{corrige}

\coursec{Bilan énergétique de l'organisme : apports et dépenses}

Dans la section précédente, nous avons vu comment nos cellules transforment l'énergie chimique des nutriments (glucides, lipides, protéines) en ATP utilisable, principalement par la respiration cellulaire et, dans certaines conditions, par la fermentation. Nous avons identifié les \cle{substrats énergétiques} et les \cle{voies métaboliques} (glycolyse, cycle de Krebs, chaîne respiratoire). Mais l'organisme n'est pas une simple cellule en culture : c'est un système ouvert, en échange permanent avec son environnement. Pour maintenir son \cle{homéostasie} (température, glycémie, masse corporelle), il doit équilibrer son \textbf{budget énergétique}. Comment savoir si nous mangeons « trop », « pas assez » ou « juste ce qu'il faut » ? C'est l'objet de cette section : apprendre à \textbf{quantifier les apports et les dépenses} pour établir un \textbf{bilan énergétique} fiable.

\courssub{Valeurs énergétiques de référence des macronutriments}

\begin{experience}[Mesure de l'énergie libérée par la combustion de nutriments]
\textbf{Problème :} Comment comparer l'énergie fournie par 1 g de riz, 1 g d'huile ou 1 g de viande ?
\textbf{Matériel :} Calorimètre à bombe (enceinte étanche en acier immergée dans un volume d'eau connu), échantillons secs de glucose, d'huile d'arachide, de blanc d'œuf (protéines), oxygène sous pression, système d'allumage électrique, thermomètre de précision.
\textbf{Protocole :}
\begin{enumerate}
    \item Peser exactement 1 g de l'échantillon sec, le placer dans le creuset du calorimètre.
    \item Remplir la bombe d'\ce{O2} à haute pression (environ 30 atm) pour assurer une combustion complète.
    \item Enregistrer la température initiale de l'eau ($T_i$).
    \item Enflammer l'échantillon par l'allumage électrique.
    \item Mesurer la température finale maximale de l'eau ($T_f$).
    \item Calculer l'énergie libérée : $Q = m_{\text{eau}} \times c_{\text{eau}} \times (T_f - T_i)$ (avec $c_{\text{eau}} \approx 4,18$ kJ$\cdot$kg$^{-1}\cdot$K$^{-1}$).
\end{enumerate}
\textbf{Observations (valeurs moyennes expérimentales, \emph{valeurs de combustion}) :}
\begin{itemize}
    \item Glucides (glucose) : $\approx$ 17,5 kJ/g.
    \item Lipides (tripalmitine) : $\approx$ 39,5 kJ/g.
    \item Protéines (albumine) : $\approx$ 23,5 kJ/g.
\end{itemize}
\textbf{Interprétation :} Les lipides libèrent plus du double d'énergie par gramme que les glucides ou les protéines. Cela s'explique par leur structure chimique : les acides gras sont très \cle{réduits} (beaucoup de liaisons C--H, peu d'atomes d'oxygène), donc leur oxydation complète libère beaucoup d'électrons pour la chaîne respiratoire.
\textbf{Correction physiologique (valeurs d'Atwater) :} Dans l'organisme, l'oxydation des protéines est incomplète (production d'urée, excrétion d'azote) et la digestibilité n'est pas de 100\%. On utilise donc les \textbf{coefficients d'Atwater} (valeurs physiologiques) pour les calculs nutritionnels :
\end{experience}

\begin{propriete}[Valeurs énergétiques physiologiques de référence (coefficients d'Atwater)]
Pour établir un bilan nutritionnel, on retient les valeurs suivantes (arrondies au kJ près) :
\begin{center}
\begin{tabularx}{\linewidth}{|l|X|c|}\hline
\rowcolor{popblueL}
\textbf{Macronutriment} & \textbf{Exemples sources camerounaises} & \textbf{Valeur énergétique (kJ/g)} \\ \hline
\cle{Glucides} & Riz, maïs, manioc, plantain, igname, fruits, sucre & 17 \\ \hline
\cle{Lipides} & Huile rouge, huile d'arachide, beurre, viande grasse, poisson gras, jaune d'œuf & 37 \\ \hline
\cle{Protéines} & Poisson sec, viande, œufs, légumineuses (haricots, niébé), lait & 17 \\ \hline
\cle{Alcool (éthanol)} & Billet, vin de palme, bière, whisky & 29 \\ \hline
\end{tabularx}
\end{center}
\textbf{Note :} Les fibres alimentaires (légumes verts, enveloppes des céréales complètes) apportent $\approx$ 8 kJ/g par fermentation colique, mais on les néglige souvent dans les bilans simples.
\end{propriete}

\begin{attention}[Piège fréquent : Confondre valeur de combustion et valeur physiologique]
Ne pas utiliser 39,5 kJ/g pour les lipides ni 23,5 kJ/g pour les protéines dans un bilan alimentaire. Au Probatoire, les énoncés précisent « valeurs d'Atwater » ou donnent les coefficients à utiliser. Par défaut, appliquez : \textbf{G = 17, L = 37, P = 17} (kJ/g).
\end{attention}

\courssub{Estimation des apports énergétiques quotidiens : le journal alimentaire}

Pour connaître les \textbf{apports} ($A$), on note tout ce qui est ingéré (solides et liquides) sur 24 h ou plusieurs jours (idéalement 3 jours dont un week-end). On convertit les portions en grammes de macronutriments via des tables de composition (ex. table Ciqual, ou données locales INRAB/Ministère de la Santé).

\begin{exemplebox}[Journal alimentaire type : journée d'un lycéen de 17 ans à Yaoundé]
\begin{center}
\begin{tabular}{|l|c|c|c|c|c|}\hline
\rowcolor{popblueL}
\textbf{Aliment / Boisson} & \textbf{Quantité} & \textbf{G (g)} & \textbf{L (g)} & \textbf{P (g)} & \textbf{Energie (kJ)} \\ \hline
\textbf{Petit-déjeuner} & & & & & \\
Pain blanc & 100 g & 55 & 2 & 8 & $55\times17 + 2\times37 + 8\times17 = 1145$ \\
Beurre & 10 g & 0 & 8 & 0 & $8\times37 = 296$ \\
Café au lait (lait entier 100 ml + sucre 10 g) & -- & 12 & 4 & 3 & $12\times17 + 4\times37 + 3\times17 = 403$ \\ \hline
\textbf{Déjeuner (cantine)} & & & & & \\
Riz blanc cuit & 250 g & 70 & 1 & 6 & $70\times17 + 1\times37 + 6\times17 = 1329$ \\
Sauce ndolé (feuilles, viande, huile, crevettes) & 300 g & 10 & 25 & 20 & $10\times17 + 25\times37 + 20\times17 = 1435$ \\
Mangue & 150 g & 20 & 0 & 1 & $20\times17 = 340$ \\ \hline
\textbf{Collation} & & & & & \\
Beurre d'arachide (sur pain) & 20 g & 4 & 10 & 5 & $4\times17 + 10\times37 + 5\times17 = 523$ \\
Jus de fruit industriel & 250 ml & 25 & 0 & 0 & $25\times17 = 425$ \\ \hline
\textbf{Dîner} & & & & & \\
Pâte de maïs (couscous) & 200 g & 50 & 1 & 5 & $50\times17 + 1\times37 + 5\times17 = 972$ \\
Sauce pistache (égusi, viande, huile) & 250 g & 8 & 20 & 15 & $8\times17 + 20\times37 + 15\times17 = 1131$ \\ \hline
\rowcolor{popcream}
\textbf{TOTAL JOUR} & & \textbf{254 g} & \textbf{71 g} & \textbf{63 g} & \textbf{7 999 kJ} $\approx$ \textbf{8 000 kJ} \\ \hline
\end{tabular}
\end{center}
\textbf{Analyse :} Les lipides ne représentent que 71 g sur 388 g de macronutriments (18\% en masse) mais apportent $71\times37 = 2627$ kJ, soit \textbf{33\% de l'énergie totale}. C'est typique d'une alimentation riche en sauces à l'huile (ndolé, pistache, arachide). Les glucides apportent 54\% de l'énergie, les protéines 13\%.
\end{exemplebox}

\begin{methode}[Établir les apports énergétiques à partir d'un journal alimentaire]
\begin{enumerate}
    \item \textbf{Lister} chaque aliment/boisson avec sa quantité (pesée ou estimée : 1 cuillère à soupe d'huile $\approx$ 10 g, 1 boule de pâte $\approx$ 150--200 g, 1 assiette de riz $\approx$ 250 g).
    \item \textbf{Rechercher} la composition en G, L, P (tables de composition, étiquettes, données locales).
    \item \textbf{Calculer} l'énergie par aliment : $E = 17\times G + 37\times L + 17\times P$ (kJ).
    \item \textbf{Sommer} sur la journée : $A_{\text{total}} = \sum E_i$.
    \item \textbf{Moyenner} sur 3 jours si possible pour lisser les variations.
\end{enumerate}
\end{methode}

\courssub{Dépenses énergétiques : les trois composantes}

Les \textbf{dépenses} ($D$) se décomposent en trois postes majeurs :
\[
D = MB + TPP + AP
\]
où $MB$ = Métabolisme de Base, $TPP$ = Thermogenèse Post-Prandiale, $AP$ = Activité Physique.

\paragraph{2.3.1 Métabolisme de Base (MB) ou Métabolisme de Repos (BMR/RMR).}
C'est l'énergie minimale pour maintenir les fonctions vitales (cœur, respiration, rénaux, activité cérébrale, turnover protéique, thermorégulation) chez un sujet \textbf{à jeun depuis 12 h, au repos complet, allongé, à la thermoneutralité (22--24 °C), éveillé}. Il représente \textbf{60 à 70\%} des dépenses totales.

\begin{propriete}[Formules prédictives du MB (kJ/jour)]
Les formules les plus utilisées en pratique clinique et sportive sont celles de \textbf{Mifflin--St Jeor (1990)}, plus précises que les anciennes formules de Harris--Benedict (1919) sur les populations actuelles.
\begin{align*}
\textbf{Homme :}\quad MB_{\text{kcal}} &= 10,0 \times m + 6,25 \times t - 5,0 \times a + 5 \\
\textbf{Femme :}\quad MB_{\text{kcal}} &= 10,0 \times m + 6,25 \times t - 5,0 \times a - 161
\end{align*}
avec $m$ = masse en kg, $t$ = taille en cm, $a$ = âge en ans. Le résultat est en \textbf{kcal/j} ; multiplier par \textbf{4,184} pour obtenir des \textbf{kJ/j}.
\textbf{Exemple :} Garçon 17 ans, 65 kg, 1,75 m.
$MB_{\text{kcal}} = 10\times65 + 6,25\times175 - 5\times17 + 5 = 650 + 1093,75 - 85 + 5 = 1663,75$ kcal/j.
$MB_{\text{kJ}} = 1663,75 \times 4,184 \approx \textbf{6 960 kJ/j}$.
\end{propriete}

\begin{attention}[Conditions d'application des formules]
Ces formules s'appliquent à des \textbf{adultes ou adolescents en fin de croissance}, en bonne santé, non obèses (IMC < 30). Pour un enfant en croissance, un sportif de haut niveau, une femme enceinte/allaitante, ou un sujet obèse, elles surestiment ou sous-estiment le MB réel. La \textbf{calorimétrie indirecte} (mesure de $\dot{V}\ce{O2}$ et $\dot{V}\ce{CO2}$ au repos) reste la référence (gold standard).
\end{attention}

\paragraph{2.3.2 Thermogenèse Post-Prandiale (TPP) ou Effet Thermique des Aliments.}
C'est l'énergie dépensée pour \textbf{digérer, absorber, transporter, métaboliser et stocker} les nutriments. Elle représente en moyenne \textbf{10\% des apports énergétiques totaux} ($A$).
\[
TPP \approx 0,10 \times A
\]
Elle varie selon la composition du repas : Protéines (20--30\% de leur énergie) $>$ Glucides (5--10\%) $>$ Lipides (0--3\%). Un repas riche en protéines « coûte » plus cher à traiter.

\paragraph{2.3.3 Dépense liée à l'Activité Physique (AP).}
C'est la composante la plus \textbf{variable} (15 à 30\% du total, jusqu'à 50\% chez un manœuvre ou un sportif). On l'estime par un \textbf{coefficient d'activité physique} (NAP / PAL -- Physical Activity Level) multiplicateur du MB :
\[
AP = (NAP - 1) \times MB \quad \text{donc} \quad D_{\text{total}} = NAP \times MB + TPP
\]
(Si on inclut la TPP dans le NAP, on simplifie souvent en $D_{\text{total}} \approx NAP \times MB$, la TPP étant « absorbée » dans le coefficient moyen).

\begin{center}
\begin{tabularx}{\linewidth}{|l|X|c|}\hline
\rowcolor{popblueL}
\textbf{Niveau d'activité (NAP / PAL)} & \textbf{Description (contexte camerounais)} & \textbf{Coefficient} \\ \hline
Sédentaire & Travail de bureau, élève/étudiant assis, transports motorisés, peu de sport & 1,40 -- 1,55 \\ \hline
Activité légère & Enseignant, commerce debout, ménagère active, marche 30--60 min/j & 1,55 -- 1,70 \\ \hline
Activité modérée & Ouvrier agricole (houe, machette), maçon, livreur à pied/vélo, sport 1 h/j & 1,70 -- 1,90 \\ \hline
Activité intense & Bûcheron, carrier, paysan sur grands champs sans machines, athlète entraînement 2--3 h/j & 1,90 -- 2,20 \\ \hline
Activité très intense & Sportif de haut niveau (double séance), militaire en opération, porteur de charges lourdes & 2,20 -- 2,50 \\ \hline
\end{tabularx}
\end{center}

\begin{exemplebox}[Calcul des dépenses pour notre lycéen de 17 ans (65 kg, 1,75 m)]
\begin{itemize}
    \item $MB \approx 6\,960$ kJ/j (calculé ci-dessus).
    \item Niveau : \textbf{Sédentaire} (cours assis, taxi/moto pour aller à l'école, devoirs, écrans, pas de sport régulier) $\rightarrow$ NAP = 1,50.
    \item $D_{\text{activite}} = NAP \times MB = 1,50 \times 6\,960 = 10\,440$ kJ/j.
    \item $TPP \approx 0,10 \times A$. Si $A = 8\,000$ kJ (journal ci-dessus), $TPP \approx 800$ kJ.
    \item \textbf{Total Dépenses $D$} $\approx 10\,440 + 800 = \textbf{11 240 kJ/j}$.
    \item \textbf{Bilan $\Delta E = A - D = 8\,000 - 11\,240 = -3\,240$ kJ/j} (déficit).
\end{itemize}
\textbf{Interprétation :} Ce lycéen est en \textbf{déficit énergétique chronique}. Il puise dans ses réserves. S'il maintient ce régime, il va maigrir.
\end{exemplebox}

\courssub{Équation du bilan énergétique et interprétation}

\begin{definition}[Bilan énergétique de l'organisme]
Le \textbf{bilan énergétique} ($\Delta E$) sur une période donnée (généralement la journée) est la différence entre l'énergie totale apportée par l'alimentation ($A$) et l'énergie totale dépensée ($D$) :
\[
\boxed{\Delta E = A - D}
\]
\begin{itemize}
    \item $\Delta E = 0$ : \textbf{Équilibre énergétique} (maintien de la masse corporelle, homéostasie).
    \item $\Delta E > 0$ : \textbf{Surplus (bilan positif)}. L'excédent est \textbf{stocké} : principalement sous forme de \cle{triglycérides} dans le \cle{tissu adipeux} (adipocytes), secondairement sous forme de \cle{glycogène} (foie, muscle) et de \cle{protéines} (muscle, si stimulus + apport protéique).
    \item $\Delta E < 0$ : \textbf{Déficit (bilan négatif)}. L'organisme \textbf{mobilise ses réserves} : lipolyse (graisses), glycogénolyse puis néoglucogenèse (protéines musculaires), oxydation des acides gras et corps cétoniques.
\end{itemize}
\end{definition}

\begin{propriete}[Équivalents énergétiques des variations de masse corporelle]
Pour relier un bilan énergétique cumulé à une variation de poids, on utilise les densités énergétiques moyennes des tissus gagnés/perdus (composition mixte eau/protéines/lipides) :
\begin{itemize}
    \item \textbf{1 kg de tissu adipeux (graisse corporelle)} $\approx$ \textbf{37 000 kJ} (87\% lipides, 13\% eau/protéines).
    \item \textbf{1 kg de tissu maigre (muscle + eau + glycogène)} $\approx$ \textbf{18 000 kJ} (73\% eau, 22\% protéines, 5\% glycogène/lipides).
    \item \textbf{Perte de poids « mixte » typique (régime modéré + activité)} : $\approx$ 25 000 à 30 000 kJ/kg (car on perd de l'eau, du glycogène, de la graisse et un peu de muscle).
\end{itemize}
\textbf{Conséquence pratique :} Un déficit constant de 500 kJ/j (environ 120 kcal/j) entraîne une perte de $\approx$ 1 kg de graisse en $37\,000 / 500 = 74$ jours ($\sim$ 2,5 mois). La perte de poids rapide au début d'un régime est surtout de l'eau (glycogène + 3 g eau/g glycogène).
\end{propriete}

\begin{exoresolu}[Bilan et prévision de variation de masse]
\textbf{Énoncé :} Une jeune fille de 16 ans, 55 kg, 1,60 m, activité modérée (NAP = 1,75). Son journal alimentaire moyen donne $A = 9\,500$ kJ/j. Calculer son MB, ses dépenses $D$, son bilan $\Delta E$ et estimer sa prise de poids sur 30 jours en supposant que le surplus est stocké uniquement en tissu adipeux.
\textbf{Solution détaillée :}
\begin{enumerate}
    \item \textbf{MB (Mifflin-St Jeor, femme) :}
    $MB_{\text{kcal}} = 10\times55 + 6,25\times160 - 5\times16 - 161 = 550 + 1000 - 80 - 161 = 1309$ kcal/j.
    $MB_{\text{kJ}} = 1309 \times 4,184 \approx \textbf{5 477 kJ/j}$.
    \item \textbf{Dépenses totales $D$ :}
    $D_{\text{activite}} = NAP \times MB = 1,75 \times 5\,477 = 9\,585$ kJ/j.
    $TPP = 0,10 \times A = 0,10 \times 9\,500 = 950$ kJ/j.
    $D = 9\,585 + 950 = \textbf{10 535 kJ/j}$.
    \item \textbf{Bilan $\Delta E$ :}
    $\Delta E = A - D = 9\,500 - 10\,535 = \textbf{-1 035 kJ/j}$.
    \item \textbf{Interprétation :} Contrairement à l'intitulé suggérant un surplus, cette jeune fille est en \textbf{déficit} de $\approx$ 1 000 kJ/j.
    \item \textbf{Variation de masse sur 30 j :}
    Déficit cumulé = $30 \times 1\,035 = 31\,050$ kJ.
    Perte de tissu adipeux $\approx 31\,050 / 37\,000 \approx \textbf{0,84 kg}$.
    Perte totale (mixte) $\approx 31\,050 / 28\,000 \approx \textbf{1,1 kg}$.
\end{enumerate}
\textbf{Conclusion :} Elle va maigrir. Pour stabiliser son poids, elle doit augmenter ses apports d'environ 1 000 kJ/j (ex. + 30 g huile + 50 g riz + 1 fruit) ou réduire son activité.
\end{exoresolu}

\courssub{Représentations graphiques du bilan énergétique}

Deux représentations sont essentielles pour visualiser et communiquer le bilan.

\paragraph{2.5.1 Histogramme : Apports par macronutriment vs Dépenses par composante.}
Ce graphique permet de voir \textbf{d'où vient l'énergie} et \textbf{où elle va}.

\begin{popfigure}
\begin{tikzpicture}
\begin{axis}[
    ybar,
    bar width=18pt,
    width=\linewidth,
    height=10cm,
    enlarge x limits=0.25,
    ylabel={Énergie (kJ/jour)},
    ymin=0, ymax=12000,
    xtick=data,
    xticklabels={Apports Glucides, Apports Lipides, Apports Protéines, MB, Activité Physique, Thermogenèse (TPP)},
    xticklabel style={rotate=30, anchor=east, font=\small},
    ymajorgrids=true,
    grid style=dashed,
    legend style={at={(0.5,-0.25)}, anchor=north, legend columns=2, font=\small},
    nodes near coords={\pgfmathprintnumber\pgfplotspointmeta},
    nodes near coords align={vertical},
    every node near coord/.append style={font=\tiny, rotate=90, anchor=west, yshift=2pt}
]
% Données de l'exemple lycéen (A=8000, D=11240)
% Apports : G=254*17=4318, L=71*37=2627, P=63*17=1071
% Dépenses : MB=6960, AP=10440-6960=3480, TPP=800
\addplot[popblue, fill=popblueL] coordinates {
    (0,4318) (1,2627) (2,1071) (3,0) (4,0) (5,0)
};
\addplot[poporange, fill=poporangeL] coordinates {
    (0,0) (1,0) (2,0) (3,6960) (4,3480) (5,800)
};
\legend{Apports (kJ), Dépenses (kJ)}
\end{axis}
\end{tikzpicture}
\legende{Histogramme comparant les apports énergétiques par macronutriment (bleu) et les dépenses par composante (orange) pour le lycéen de l'exemple (A $\approx$ 8 000 kJ, D $\approx$ 11 240 kJ). Le déficit est visible : la somme des barres bleues est inférieure à la somme des barres orange.}
\end{popfigure}

\paragraph{2.5.2 Courbe de variation de la masse corporelle sur plusieurs semaines.}
On modélise l'évolution du poids $M(t)$ (en kg) sur $N$ semaines pour trois scénarios théoriques à partir d'un poids initial $M_0 = 65$ kg, $MB_0 \approx 7\,000$ kJ/j, $NAP = 1,55$ ($D_{\text{maintien}} \approx 10\,850$ kJ/j + TPP).
\begin{itemize}
    \item \textbf{Scénario 1 (Équilibre) :} $A = D_{\text{maintien}} \approx 11\,000$ kJ/j $\rightarrow \Delta E = 0$.
    \item \textbf{Scénario 2 (Surplus +500 kJ/j) :} $A = 11\,500$ kJ/j $\rightarrow \Delta E = +500$ kJ/j. Stockage adipeux $\approx 500/37\,000 = 0,0135$ kg/j $\approx 0,095$ kg/sem.
    \item \textbf{Scénario 3 (Déficit -500 kJ/j) :} $A = 10\,500$ kJ/j $\rightarrow \Delta E = -500$ kJ/j. Perte mixte $\approx 500/28\,000 = 0,018$ kg/j $\approx 0,125$ kg/sem (plus rapide au début car perte d'eau/glycogène).
\end{itemize}

\begin{popfigure}
\begin{tikzpicture}
\begin{axis}[
    width=\linewidth,
    height=10cm,
    xlabel={Temps (semaines)},
    ylabel={Masse corporelle (kg)},
    xmin=0, xmax=8,
    ymin=62, ymax=68,
    xtick={0,1,2,3,4,5,6,7,8},
    ytick={62,63,64,65,66,67,68},
    grid=major,
    grid style={dashed, popmuted},
    legend style={at={(0.02,0.98)}, anchor=north west, font=\small, fill=popcream, draw=popline},
    title style={font=\normalsize},
    title={Modélisation de l'évolution pondérale sur 8 semaines (M$_0$ = 65 kg)}
]
% Scénario 1 : Equilibre (ligne horizontale)
\addplot[popgreen, thick, mark=*] coordinates {
    (0,65) (1,65) (2,65) (3,65) (4,65) (5,65) (6,65) (7,65) (8,65)
};
\addlegendentry{Équilibre ($\Delta E = 0$)}

% Scénario 2 : Surplus +500 kJ/j -> +0.095 kg/sem (linéaire simplifié)
\addplot[poporange, thick, mark=square*] coordinates {
    (0,65) (1,65.095) (2,65.19) (3,65.285) (4,65.38) (5,65.475) (6,65.57) (7,65.665) (8,65.76)
};
\addlegendentry{Surplus +500 kJ/j (+0,76 kg)}

% Scénario 3 : Déficit -500 kJ/j -> -0.125 kg/sem (courbe légèrement concave au début)
% On simule une perte plus rapide les 2 premières semaines (eau/glycogène) puis linéaire
\addplot[popblue, thick, mark=triangle*] coordinates {
    (0,65) (1,64.7) (2,64.45) (3,64.32) (4,64.20) (5,64.07) (6,63.95) (7,63.82) (8,63.70)
};
\addlegendentry{Déficit -500 kJ/j (-1,3 kg)}
\end{axis}
\end{tikzpicture}
\legende{Évolution simulée de la masse corporelle sur 8 semaines pour un sujet de 65 kg (MB $\approx$ 7 000 kJ/j, NAP 1,55). L'équilibre maintient le poids. Un surplus constant de 500 kJ/j entraîne une prise de $\sim$ 0,76 kg (principalement graisse). Un déficit de 500 kJ/j entraîne une perte de $\sim$ 1,3 kg (eau + graisse + peu de muscle), plus marquée les premières semaines.}
\end{popfigure}

\begin{aretenir}[Lecture d'une courbe de variation pondérale]
\begin{itemize}
    \item \textbf{Pente} $\approx$ vitesse de changement des réserves ($\Delta M / \Delta t$).
    \item \textbf{Plateau} : nouvel équilibre (apports = nouvelles dépenses adaptées au nouveau poids) ou fin de la restriction.
    \item \textbf{Variations brusques (1--2 kg en quelques jours)} : quasi exclusivement \textbf{eau} (glycogène, sel, hydratation, cycle menstruel), pas de graisse.
    \item \textbf{Seuil d'alerte :} Perte $> 1$ kg/semaine prolongée $\rightarrow$ risque de fonte musculaire, calculs biliaires, troubles électrolytiques.
\end{itemize}
\end{aretenir}

\courssub{Mesures correctives, prévention et éducation nutritionnelle (introduction)}

L'établissement du bilan n'est pas une fin en soi ; il guide l'action. La section 5 développera les stratégies, mais retenons déjà les principes de base :

\begin{methode}[Corriger un déséquilibre énergétique identifié]
\begin{enumerate}
    \item \textbf{Si $\Delta E > 0$ (surplus / prise de poids non désirée) :}
    \begin{itemize}
        \item \textbf{Réduire les apports} : privilégier la densité nutritionnelle (légumes, protéines maigres, céréales complètes) vs densité énergétique (huiles, sucres ajoutés, fritures, boissons sucrées, alcool). Cibler -500 kJ/j pour une perte de $\sim$ 0,5 kg/sem.
        \item \textbf{Augmenter les dépenses} : augmenter le NAP (marche rapide 30 min/j = +0,15 à 0,20 sur NAP $\approx$ +1 000 à 1 400 kJ/j), musculation (augmente le MB à long terme).
        \item \textbf{Combiner les deux} (plus durable, moins de frustration, préserve la masse musculaire).
    \end{itemize}
    \item \textbf{Si $\Delta E < 0$ (déficit / maigreur, dénutrition) :}
    \begin{itemize}
        \item \textbf{Augmenter les apports} : fractionner les repas (3 repas + 2--3 collations), enrichir en énergie sans augmenter le volume (huile, beurre d'arachide, lait entier, poudre de lait, œufs, purées de légumineuses). Cibler +500 à +1 000 kJ/j.
        \item \textbf{Adapter l'activité} : réduire l'intensité si catabolisme excessif, privilégier la musculation douce pour stimuler l'appétit et l'anabolisme.
        \item \textbf{Rechercher la cause} : pathologie (TB, VIH, diabète, hyperthyroïdie, cancer), troubles du comportement alimentaire, insécurité alimentaire, malabsorption.
    \end{itemize}
    \item \textbf{Suivi :} Pesée hebdomadaire (même conditions : matin, à jeun, nu), réévaluation du journal alimentaire et du NAP tous les 15 jours, ajustement progressif.
\end{enumerate}
\end{methode}

\begin{savaistu}[Le « piège » de la thermogenèse adaptative]
Quand on maintient un \textbf{déficit énergétique prolongé}, l'organisme défend son « set point » (poids de référence) en \textbf{réduisant ses dépenses} : le MB chute de 10--15\% au-delà de la perte de masse maigre (thermogénèse adaptative), la spontanéité motrice (NEAT -- \emph{Non-Exercise Activity Thermogenesis}) diminue inconsciemment (on bouge moins, on a froid, on est moins fidgety). C'est pourquoi les régimes très restrictifs « bloquent » et pourquoi la reprise de poids est fréquente (effet yoyo). \textbf{Solution :} déficit modéré, activité physique régulière (maintient le MB et le NEAT), apports protéiques suffisants ($> 1,2$ g/kg/j) pour préserver la masse musculaire.
\end{savaistu}

\begin{exercice}[Application : Bilan énergétique d'une vendeuse de marché]
\textbf{Contexte :} Mme M., 42 ans, 68 kg, 1,58 m. Elle tient un stand de légumes au marché Mfoundi (debout, marche, porte des charges légères, 10 h/j). Elle dort 7 h. Son journal alimentaire type :
\begin{itemize}
    \item Matin : Bouillie de maïs (300 g) + lait concentré (30 g) + sucre (20 g).
    \item Midi : Riz (250 g) + sauce gombo/viande (200 g, 15 g huile) + banane plantain frite (100 g, 10 g huile absorbée).
    \item Soir : Pâte de manioc (200 g) + sauce feuilles (150 g, 10 g huile) + poisson sec (50 g).
    \item Collations : 2 mangues (300 g), 1 sachet d'eau (0,5 L), 1 jus bissap sucré (250 ml, 25 g sucre).
\end{itemize}
\textbf{Données composition (approx. pour 100 g ou portion indiquée) :} Bouillie 100g : 20G/2L/2P ; Lait concentré 100g : 55G/8L/8P ; Riz cuit 100g : 28G/0,3L/2,5P ; Gombo/viande 100g : 4G/7L/6P ; Plantain frit 100g : 30G/10L/1,5P (inclut l'huile absorbée) ; Pâte manioc 100g : 25G/0,5L/1,5P ; Sauce feuilles 100g : 3G/6L/2P ; Poisson sec 100g : 0G/5L/60P ; Mangue 100g : 14G/0L/0,5P ; Jus bissap 100ml : 10G/0L/0P.
\textbf{Questions :}
\begin{enumerate}
    \item Calculer les apports quotidiens totaux $A$ (kJ) et la répartition \% G / \% L / \% P.
    \item Estimer son MB (kJ/j) par Mifflin-St Jeor.
    \item Choisir un NAP justifié et calculer $D_{\text{total}}$ (inclure TPP).
    \item Calculer $\Delta E$ et conclure sur l'évolution pondérale prévisible.
    \item Proposer deux modifications alimentaires concrètes pour corriger le déséquilibre si nécessaire.
\end{enumerate}
\end{exercice}

\begin{corrige}[Exercice : Bilan énergétique d'une vendeuse de marché]
\textbf{1. Apports $A$ (calcul ligne par ligne) :}
\begin{align*}
\text{Bouillie (300g)} &: 60G, 6L, 6P \rightarrow 60\times17 + 6\times37 + 6\times17 = 1344 \text{ kJ} \\
\text{Lait conc. (30g)} &: 16,5G, 2,4L, 2,4P \rightarrow 280,5 + 88,8 + 40,8 = 410 \text{ kJ} \\
\text{Sucre (20g)} &: 20G \rightarrow 340 \text{ kJ} \\
\text{Riz (250g)} &: 70G, 0,75L, 6,25P \rightarrow 1190 + 27,75 + 106,25 = 1324 \text{ kJ} \\
\text{Sauce gombo (200g)} &: 8G, 14L, 12P \rightarrow 136 + 518 + 204 = 858 \text{ kJ} \\
\text{Huile midi (15g)} &: 15L \rightarrow 555 \text{ kJ} \\
\text{Plantain frit (100g)} &: 30G, 10L, 1,5P \rightarrow 510 + 370 + 25,5 = 905,5 \text{ kJ} \\
\text{Pâte manioc (200g)} &: 50G, 1L, 3P \rightarrow 850 + 37 + 51 = 938 \text{ kJ} \\
\text{Sauce feuilles (150g)} &: 4,5G, 9L, 3P \rightarrow 76,5 + 333 + 51 = 460,5 \text{ kJ} \\
\text{Huile soir (10g)} &: 10L \rightarrow 370 \text{ kJ} \\
\text{Poisson sec (50g)} &: 0G, 2,5L, 30P \rightarrow 92,5 + 510 = 602,5 \text{ kJ} \\
\text{Mangues (300g)} &: 42G, 0L, 1,5P \rightarrow 714 + 25,5 = 739,5 \text{ kJ} \\
\text{Jus bissap (250ml)} &: 25G \rightarrow 425 \text{ kJ} \\
\textbf{Total } A &= \textbf{9 272 kJ} \\
\textbf{Total G} &= 326 \text{ g } (5542 \text{ kJ}, 60\%) \\
\textbf{Total L} &= 60,65 \text{ g } (2244 \text{ kJ}, 24\%) \\
\textbf{Total P} &= 65,65 \text{ g } (1116 \text{ kJ}, 12\%) \quad (\text{le reste } \approx 4\% \text{ fibres/autres})
\end{align*}
\textit{Note : La ligne « Huile plantain (10g) » du journal initial n'est pas comptabilisée séparément car l'huile absorbée est déjà incluse dans la composition du plantain frit (10L/100g).}

\textbf{2. MB (Femme, 42 ans, 68 kg, 158 cm) :}
$MB_{\text{kcal}} = 10\times68 + 6,25\times158 - 5\times42 - 161 = 680 + 987,5 - 210 - 161 = 1296,5$ kcal.
$MB_{\text{kJ}} = 1296,5 \times 4,184 \approx \textbf{5 425 kJ/j}$.

\textbf{3. NAP et Dépenses $D$ :}
Activité : debout 10h, marche, port charges $\rightarrow$ \textbf{Activité modérée}. NAP = 1,75.
$D_{\text{activite}} = 1,75 \times 5\,425 = 9\,494$ kJ.
$TPP = 0,10 \times 9\,272 = 927$ kJ.
$D_{\text{total}} = 9\,494 + 927 = \textbf{10 421 kJ/j}$.

\textbf{4. Bilan $\Delta E$ :}
$\Delta E = 9\,272 - 10\,421 = \textbf{-1 149 kJ/j}$.
\textbf{Conclusion :} Déficit modéré mais significatif. Perte de poids prévisible $\sim$ 1 kg/mois en graisse (plus au début avec l'eau). Risque de fatigue en fin de journée, fonte musculaire à long terme si apports protéiques insuffisants (65,6 g = 0,96 g/kg, limite basse pour une activité modérée).

\textbf{5. Propositions correctives (pour tendre vers $\Delta E \approx 0$) :}
\begin{itemize}
    \item \textbf{Augmenter l'apport protéique/énergétique du matin :} Remplacer le sucre (20 g, 340 kJ) par 30 g d'arachides grillées (+ 10G, +14L, +7P $\approx$ 850 kJ) ou ajouter un œuf dur (+ 0G, +5L, +6P $\approx$ 485 kJ). Gain net $\approx$ +500 à +850 kJ.
    \item \textbf{Enrichir la collation de l'après-midi :} Ajouter 100 g yaourt nature + 20 g arachides ($\approx$ 600 kJ, 10 g P, calcium).
    \item \textbf{Alternative simple :} Augmenter la portion de riz du midi de 50 g (+850 kJ) et ajouter 10 g d'huile dans la sauce du soir (+370 kJ).
\end{itemize}
Objectif : atteindre $\approx$ 10 500 kJ/j avec $\ge$ 1,2 g P/kg/j (82 g P).
\end{corrige}

\begin{attention}[Erreurs fréquentes à l'examen (Probatoire)]
\begin{itemize}
    \item \textbf{Oublier la TPP} (10\% des apports) : sous-estime $D$ de $\sim$ 800--1 200 kJ $\rightarrow$ faux équilibre ou faux surplus.
    \item \textbf{Mauvaise conversion kcal $\leftrightarrow$ kJ} : 1 kcal = 4,184 kJ (souvent arrondi 4,18). Vérifier l'unité demandée dans l'énoncé.
    \item \textbf{Utiliser les valeurs de combustion} (39,5 ; 23,5) au lieu des valeurs d'Atwater (37 ; 17).
    \item \textbf{Choisir un NAP incohérent} avec le profil décrit (ex. NAP 1,4 pour un paysan, ou 2,0 pour un élève sédentaire). Justifier toujours le choix par une phrase.
    \item \textbf{Confondre masse perdue et graisse perdue} : 1 kg sur la balance $\neq$ 1 kg de graisse (37 000 kJ). Préciser « tissu adipeux » ou « masse corporelle mixte ».
    \item \textbf{Ne pas adapter le MB au nouveau poids} dans un suivi sur plusieurs semaines : le MB baisse quand on maigrit, le déficit effectif diminue $\rightarrow$ courbe asymptotique, pas linéaire indéfiniment.
    \item \textbf{Double-compter les huiles de cuisson} : si la composition de l'aliment cuit (ex. plantain frit, sauce) inclut déjà l'huile, ne pas ajouter une ligne « huile » séparée.
\end{itemize}
\end{attention}

\coursec{Conséquences physiologiques des déséquilibres énergétiques}

Dans la section précédente, nous avons établi le \cle{bilan énergétique} de l'organisme en comparant les apports nutritionnels (glucides, lipides, protides) aux dépenses (métabolisme de base, thermogenèse, activité physique). Lorsque ce bilan n'est pas nul sur la durée, l'organisme doit adapter ses voies métaboliques pour stocker ou mobiliser de l'énergie. Ces adaptations, si elles persistent, engendrent des modifications physiologiques majeures que nous allons analyser.

\courssub{Déséquilibre énergétique positif chronique}

\begin{definition}[Déséquilibre énergétique positif]
Un déséquilibre énergétique positif se définit par des apports énergétiques supérieurs aux dépenses sur une période prolongée (semaines à mois). L'excédent $\Delta E = \text{Apports} - \text{Dépenses} > 0$ est stocké, principalement sous forme de triglycérides dans le tissu adipeux.
\end{definition}

\textbf{Observation.} Au Cameroun, de nombreux jeunes adultes vivant à Douala ou Yaoundé consomment quotidiennement des repas riches en huile de palme raffinée, en riz blanc et en boissons sucrées, tout en ayant une activité physique réduite (transport en moto-taxi, travail de bureau). Sur plusieurs mois, une prise de poids progressive est constatée.

\textbf{Hypothèse.} L'excédent énergétique chronique stimule la \cle{lipogenèse} hépatique et adipocytaire, entraînant une hypertrophie puis une hyperplasie des adipocytes, accompagnée d'une sécrétion altérée d'adipokines.

\textbf{Données expérimentales.} Des études cliniques montrent qu'un surplus constant de \num{2000}~\SI{}{kJ} par jour pendant 30 jours correspond à un excédent total de \num{60000}~\SI{}{kJ}. Sachant que la densité énergétique du tissu adipeux est d'environ \num{37}~\SI{}{kJ/g}, le gain théorique de masse grasse est de $\frac{60000}{37} \approx 1620$~\SI{}{g}, soit \num{1,6}~\SI{}{kg}.

\textbf{Interprétation.} L'excédent favorise la synthèse d'acides gras (lipogenèse de novo) à partir du glucose excédentaire via l'acétyl-CoA carboxylase (ACC) et la synthase des acides gras (FAS). Les adipocytes stockent les triglycérides, augmentant de volume (hypertrophie) et, au-delà d'un seuil critique, de nombre (hyperplasie). Le tissu adipeux devient un organe endocrinien actif : il sécrète de la \cle{leptine} (signal de satiété proportionnel à la masse grasse) et de l'\cle{adiponectine} (sensibilisant à l'insuline). Chez l'obèse, une \cle{résistance à la leptine} se développe : malgré des taux plasmatiques élevés, le signal anorexigène n'atteint plus l'hypothalamus, maintenant l'appétit. Parallèlement, l'adiponectine diminue, contribuant à la \cle{résistance à l'insuline}. Cette résistance favorise l'hyperglycémie, l'hyperinsulinémie compensatrice, la dyslipidémie (augmentation des VLDL, diminution du HDL) et l'hypertension artérielle par activation du système rénine-angiotensine et rétention sodée.

\textbf{Conclusion.} Un déséquilibre positif chronique conduit à l'obésité, au syndrome métabolique et augmente le risque de diabète de type 2, de maladies cardiovasculaires et de certains cancers.

\begin{propriete}[Variation de masse et excédent énergétique (admis)]
La variation de masse corporelle $\Delta m$ (en kg) liée à un déséquilibre énergétique $\Delta E$ (en kJ) sur une période donnée s'approxime par :
\[
\Delta m \approx \frac{\Delta E}{\rho_{\text{tissu}}}
\]
où $\rho_{\text{tissu}}$ est la densité énergétique du tissu stocké ou mobilisé (\num{37}~\SI{}{kJ/g} pour le tissu adipeux, \num{17}~\SI{}{kJ/g} pour le glycogène hydraté, \num{7}~\SI{}{kJ/g} pour les protéines hydratées). Cette relation est admise au Probatoire ; elle ne nécessite pas de démonstration.
\end{propriete}

\begin{attention}[Perte de poids rapide $\neq$ perte de masse grasse]
Une perte de \num{2} à \num{3}~\SI{}{kg} en quelques jours lors d'un régime drastique correspond essentiellement à une perte d'eau (glycogène hydraté : \num{3} à \num{4}~g d'eau par g de glycogène) et de glycogène hépatique et musculaire, non à une mobilisation massive de triglycérides. La vraie perte de masse grasse est lente (\num{0,5} à \num{1}~\SI{}{kg/semaine} en déficit modéré).
\end{attention}

\begin{exoresolu}[Gain de masse grasse théorique]
\textbf{Énoncé.} Un élève de Première D consomme en moyenne \num{12000}~\SI{}{kJ/j} alors que ses dépenses totales sont estimées à \num{10000}~\SI{}{kJ/j}. Calcule le gain de masse grasse théorique après 60 jours si tout l'excédent est stocké sous forme de tissu adipeux.

\textbf{Solution détaillée.}
\begin{align*}
\Delta E_{\text{jour}} &= 12000 - 10000 = 2000~\SI{}{kJ/j} \\
\Delta E_{\text{60 j}} &= 2000 \times 60 = 120000~\SI{}{kJ} \\
\Delta m_{\text{graisse}} &= \frac{120000~\SI{}{kJ}}{37~\SI{}{kJ/g}} \approx 3243~\SI{}{g} \approx 3,2~\SI{}{kg}
\end{align*}
L'élève gagnerait théoriquement \num{3,2}~\SI{}{kg} de tissu adipeux en deux mois. En réalité, une partie de l'excédent augmente la thermogenèse adaptative et le coût énergétique du tissu adipeux lui-même, réduisant légèrement ce gain.
\end{exoresolu}

\begin{savaistu}[La leptine, hormone de la satiété]
La leptine est une protéine de \num{16}~\SI{}{kDa} sécrétée par les adipocytes. Son taux plasmatique est proportionnel à la masse grasse. Elle agit sur les récepteurs Ob-Rb de l'hypothalamus arqué pour inhiber la NPY/AgRP (orexigène) et stimuler la POMC/CART (anorexigène). Chez l'obèse, la barrière hémato-encéphalique limite le transport de la leptine et les voies de signalisation JAK2/STAT3 sont désensibilisées : c'est la \emph{résistance à la leptine}. Cela explique pourquoi l'obèse mange malgré des réserves énergétiques énormes.
\end{savaistu}

\begin{popfigure}
\begin{tikzpicture}[
  node distance=1.5cm and 2cm,
  box/.style={rectangle, draw=popdark, thick, fill=popcream, rounded corners, minimum width=2.5cm, minimum height=1.2cm, align=center},
  arrow/.style={-Stealth, thick, popblue},
  label/.style={font=\small, text=popdark}
]
% Nœuds colonne surplus
\node[box] (gluc) {Glucose excédentaire};
\node[box, below=of gluc] (acCoA) {Acétyl-CoA};
\node[box, below=of acCoA] (lipogen) {Lipogenèse (FAS, ACC)};
\node[box, below=of lipogen] (tg) {Triglycérides};
\node[box, below=of tg, align=center] (adipo){Adipocyte : stockage \\ hypertrophie / hyperplasie};
\node[box, right=of adipo, xshift=1cm, align=center] (adipoSec){Sécrétion adipokines : \\ Leptine $\uparrow$, Adiponectine $\downarrow$};
\node[box, below=of adipoSec, align=center] (resist){Résistance à l'insuline \\ \& à la leptine};
\node[box, below=of resist, align=center] (patho){Syndrome métabolique \\ Diabète type 2, HTA, Dyslipidémie};

% Nœuds colonne déficit
\node[box, right=of gluc, xshift=3cm] (jeune) {Jeûne / Effort prolongé};
\node[box, below=of jeune, align=center] (lipolyse){Lipolyse (ATGL, HSL) \\ $\rightarrow$ Acides gras libres + Glycérol};
\node[box, below=of lipolyse, align=center] (betaox){$\beta$-oxydation mitochondriale \\ $\rightarrow$ Acétyl-CoA + NADH + FADH2};
\node[box, below=of betaox, align=center] (ketogen){Cétogenèse hépatique \\ (corps cétoniques)};
\node[box, below=of ketogen, align=center] (neogluc){Néoglucogenèse (glycérol, AA) \\ $\rightarrow$ Glucose};
\node[box, below=of neogluc, align=center] (mobil){Mobilisation protéique \\ (perte musculaire, immunité)};

% Flèches surplus
\draw[arrow] (gluc) -- (acCoA);
\draw[arrow] (acCoA) -- (lipogen);
\draw[arrow] (lipogen) -- (tg);
\draw[arrow] (tg) -- (adipo);
\draw[arrow] (adipo) -- (adipoSec);
\draw[arrow] (adipoSec) -- (resist);
\draw[arrow] (resist) -- (patho);

% Flèches déficit
\draw[arrow] (jeune) -- (lipolyse);
\draw[arrow] (lipolyse) -- (betaox);
\draw[arrow] (betaox) -- (ketogen);
\draw[arrow] (betaox) -- (neogluc);
\draw[arrow] (neogluc) -- (mobil);

% Étiquettes de colonnes
\node[label, above=0.2cm of gluc] {\textbf{Surplus énergétique (lipogenèse)}};
\node[label, above=0.2cm of jeune] {\textbf{Déficit énergétique (lipolyse / néoglucogenèse)}};
\end{tikzpicture}
\legende{Schéma comparatif des flux métaboliques en situation de surplus énergétique chronique (à gauche) et de déficit énergétique prolongé (à droite). Les flèches indiquent le sens des transformations biochimiques principales.}
\end{popfigure}

\courssub{Déséquilibre énergétique négatif prolongé}

\begin{definition}[Déséquilibre énergétique négatif]
Un déséquilibre énergétique négatif se caractérise par des apports inférieurs aux dépenses sur une période prolongée ($\Delta E < 0$). L'organisme mobilise ses réserves : d'abord le glycogène hépatique, puis les triglycérides adipeux (lipolyse), et enfin les protéines structurelles (catabolisme protéique).
\end{definition}

\textbf{Observation.} Lors d'une restriction calorique sévère (ex. régime \num{5000}~\SI{}{kJ/j} pour un besoin de \num{10000}~\SI{}{kJ/j}) ou d'une pathologie augmentant les dépenses (fièvre prolongée, hyperthyroïdie, cancer), le poids chute rapidement les premiers jours puis plus lentement.

\textbf{Hypothèse.} La mobilisation successive des réserves suit une hiérarchie énergétique : glycogène (court terme), lipides (moyen terme), protéines (long terme). Le catabolisme protéique entraîne une perte de masse musculaire et une altération des fonctions immunitaires.

\textbf{Données physiologiques.}
\begin{itemize}
  \item \textbf{Jours 1-2 :} Glycogène hépatique (\num{100} à \num{120}~\SI{}{g}) et musculaire (\num{300} à \num{400}~\SI{}{g}) épuisés $\rightarrow$ perte d'eau associée (\num{3} à \num{4}~g eau/g glycogène) $\approx$ \num{1,5} à \num{2}~\SI{}{kg}.
  \item \textbf{Jours 3-14 :} Lipolyse prédominante. Les acides gras libres (AGL) alimentent la $\beta$-oxydation hépatique et musculaire ; le foie produit des corps cétoniques (acétoacétate, $\beta$-hydroxybutyrate) qui deviennent le carburant principal du cerveau (jusqu'à \num{60}~\% de ses besoins). La protéolyse est limitée par l'effet épargnant des corps cétoniques.
  \item \textbf{Au-delà de 2-3 semaines :} Diminution de la leptine $\rightarrow$ baisse du métabolisme de base (thermogénèse adaptative). Si le déficit persiste, la néoglucogenèse hépatique utilise massivement les acides aminés (alanine, glutamine) provenant du muscle $\rightarrow$ perte de masse maigre, affaiblissement immunitaire (lymphopénie), aménorrhée chez la femme (arrêt de l'axe gonadotrope par manque de leptine et d'insuline), ostéoporose (baisse de l'ostéocalcine, augmentation du cortisol).
\end{itemize}

\textbf{Interprétation.} L'organisme privilégie la survie du cerveau (glucose puis corps cétoniques) au détriment des tissus périphériques. La perte de protéines structurelles (muscle squelettique, muscle cardiaque, protéines de la matrice osseuse) devient critique si le déficit n'est pas corrigé.

\textbf{Conclusion.} Un déséquilibre négatif prolongé conduit à la dénutrition protéino-énergétique, avec des complications infectieuses, cardiaques, osseuses et reproductives.

\begin{attention}[Effet yoyo et rebond pondéral]
Après une restriction sévère, la reprise alimentaire s'accompagne souvent d'une reprise de poids supérieure à la perte initiale (effet yoyo). La thermogenèse adaptative persiste plusieurs mois, la leptine reste basse, l'appétit est accru (hyperphagie de rebond) et la lipogenèse est favorisée par une sensibilité à l'insuline transitoirement augmentée. Ce phénomène explique l'échec fréquent des régimes hypocaloriques drastiques.
\end{attention}

\begin{popfigure}
\begin{tikzpicture}
\begin{axis}[
  name=axmasse,
  width=\linewidth,
  height=8cm,
  xlabel={Temps (semaines)},
  ylabel={Masse corporelle (kg)},
  ymin=55, ymax=85,
  xmin=0, xmax=20,
  xtick={0,4,8,12,16,20},
  ytick={55,60,65,70,75,80,85},
  grid=major,
  grid style={popline},
  legend pos=north west,
  legend style={fill=popcream, draw=popdark},
  axis line style={popdark},
  tick style={popdark},
  label style={popdark},
  tick label style={popdark},
  scale only axis
]
% Courbe masse
\addplot[popcurve, thick, popblue, mark=*] coordinates {
  (0,78) (2,76.5) (4,74) (6,72) (8,70.5) (10,69.5) (12,69) (14,69.2) (16,70) (18,71.5) (20,73)
};
\addlegendentry{Masse (kg)}
\end{axis}
% Deuxième axe (non imbriqué, superposé) pour tour de taille
\begin{axis}[
  at={(axmasse.south west)}, anchor=south west,
  width=\linewidth,
  height=8cm,
  xlabel={Temps (semaines)},
  ylabel={Tour de taille (cm)},
  ymin=70, ymax=100,
  xmin=0, xmax=20,
  xtick={0,4,8,12,16,20},
  ytick={70,75,80,85,90,95,100},
  axis y line*=right,
  axis x line=none,
  legend pos=north east,
  legend style={fill=popcream, draw=popdark},
  axis line style={poporange},
  tick style={poporange},
  label style={poporange},
  tick label style={poporange},
  scale only axis
]
\addplot[popcurve, thick, poporange, mark=square*] coordinates {
  (0,92) (2,90) (4,87) (6,84) (8,82) (10,81) (12,81) (14,81.5) (16,83) (18,85) (20,87)
};
\addlegendentry{Tour de taille (cm)}
\end{axis}
\end{tikzpicture}
\legende{Évolution superposée de la masse corporelle (kg, courbe bleue) et du tour de taille (cm, courbe orange) chez un sujet suivant un régime hypocalorique de 8 semaines suivi d'une reprise alimentaire. On observe une pente négative initiale, un plateau (semaines 10-14) puis un rebond (effet yoyo) avec reprise du tour de taille plus rapide que la masse.}
\end{popfigure}

\courssub{Signes cliniques et biologiques associés}

\begin{center}
\begin{tabularx}{\linewidth}{|l|X|X|}\hline
\rowcolor{popblueL}
\textbf{Paramètre} & \textbf{Déséquilibre positif (obésité)} & \textbf{Déséquilibre négatif (dénutrition)} \\ \hline
Glycémie à jeun & $\ge 1,10$~\SI{}{g/L} (prédiabète) ou $\ge 1,26$~\SI{}{g/L} (diabète) & $\le 0,70$~\SI{}{g/L} (hypoglycémie de jeûne) \\ \hline
Profil lipidique & Triglycérides $\uparrow$, LDL-c $\uparrow$, HDL-c $\downarrow$ & Triglycérides $\downarrow$, cholestérol total $\downarrow$ (sauf stéatose hépatique de dénutrition) \\ \hline
CRP ultrasensible & $\ge 3$~\SI{}{mg/L} (inflammation bas bruit) & Normal ou $\uparrow$ si infection associée \\ \hline
Hémoglobine & Normale ou $\uparrow$ (hémoconcentration) & $\downarrow$ (anémie ferriprive, anémie inflammatoire) \\ \hline
Leptine & $\uparrow\uparrow$ (résistance) & $\downarrow\downarrow$ (signal de famine) \\ \hline
Adiponectine & $\downarrow$ & $\uparrow$ (mais insuffisante) \\ \hline
Cortisol & Normal ou $\uparrow$ (syndrome de Cushing pseudo) & $\uparrow$ (stress catabolique) \\ \hline
\end{tabularx}
\end{center}

Ces marqueurs permettent au clinicien de confirmer le type de déséquilibre, d'en évaluer la sévérité et de surveiller l'efficacité des mesures correctives.

\courssub{Interprétation de courbes de variation pondérale}

Une courbe de variation de la masse corporelle dans le temps (abscisse : semaines ou mois ; ordonnée : kg) révèle la dynamique du bilan énergétique :

\begin{itemize}
  \item \textbf{Pente positive constante} : excédent énergétique stable $\rightarrow$ prise de poids linéaire (phase d'installation de l'obésité).
  \item \textbf{Pente négative abrupte puis plateau} : perte rapide d'eau/glycogène (premiers jours) puis perte plus lente de masse grasse (lipolyse) ; le plateau correspond à l'équilibre nouveau (thermogénèse adaptative, réduction du métabolisme de base).
  \item \textbf{Rebond (effet yoyo)} : après la fin de la restriction, la courbe remonte au-dessus du point de départ. La reprise du tour de taille précède souvent la reprise de la masse, traduisant une redistribution préférentielle vers le tissu adipeux viscéral.
  \item \textbf{Oscillations répétées} : cycles de restriction/reprise (régimes yo-yo) augmentent le risque métabolique (résistance à l'insuline, inflammation) même si la masse moyenne reste stable.
\end{itemize}

L'analyse conjointe de la masse et du tour de taille (ou de l'impédancemétrie) permet de distinguer perte de masse grasse vs perte de masse maigre.

\courssub{Contexte camerounais : transition nutritionnelle et prévention}

Le Cameroun traverse une \cle{transition nutritionnelle} rapide. L'urbanisation (Douala, Yaoundé, Bafoussam) s'accompagne de :
\begin{itemize}
  \item Remplacement des aliments traditionnels (igname, plantain, légumes-feuilles, poisson fumé) par des produits transformés : huiles de palme raffinées (riches en acides gras saturés), boissons sucrées (jus industriels, sodas), pains blancs, plats rapides frits.
  \item Sédentarité accrue : transport motorisé, emplois tertiaires, loisirs écrans.
  \item Publicité ciblant les jeunes pour les produits ultra-transformés.
\end{itemize}
Conséquence : augmentation de la prévalence du surpoids/obésité chez l'adulte (\num{25}~\% en zone urbaine selon l'enquête STEPS 2017) et apparition du diabète de type 2 chez l'adolescent. Parallèlement, la sous-nutrition persiste en zone rurale et chez les populations déplacées (insécurité alimentaire, maladies infectieuses).

\begin{methode}[Relier un déséquilibre observé à ses conséquences et proposer des mesures correctives]
\begin{enumerate}
  \item \textbf{Quantifier le bilan} : calculer apports (enquête alimentaire 3 jours) et dépenses (formule de Black et al. + coefficient d'activité). Déterminer $\Delta E$ journalier.
  \item \textbf{Analyser la courbe pondérale} : identifier pente, plateaux, rebonds ; corréler avec tour de taille et impédancemétrie si possible.
  \item \textbf{Rechercher les marqueurs biologiques} : glycémie à jeun, profil lipidique, CRPus, hémoglobine, leptine/adiponectine si disponible.
  \item \textbf{Établir le diagnostic} : surpoids/obésité (IMC $\ge 25$), syndrome métabolique (critères IDF), dénutrition (IMC $< 18,5$ ou perte $> 10~\%$ en 6 mois).
  \item \textbf{Proposer un plan individualisé} :
        \begin{itemize}
          \item Si $\Delta E > 0$ : réduire apports de \num{500} à \num{1000}~\SI{}{kJ/j} (limiter huiles, sucres, portions), augmenter activité (\num{150} min/semaine d'endurance modérée), viser perte \num{0,5} à \num{1}~\SI{}{kg/semaine}.
          \item Si $\Delta E < 0$ : augmenter apports protéiques (\num{1,2} à \num{1,5}~g/kg/j) et énergétiques (\num{+500} à \num{+1000}~\SI{}{kJ/j}), supplémentation micronutriments, réhabilitation progressive de l'activité physique.
        \end{itemize}
  \item \textbf{Suivi} : pesée hebdomadaire, tour de taille mensuel, biologie à 3 mois, éducation nutritionnelle continue (cuisson vapeur, huiles non raffinées, fruits locaux, légumineuses).
\end{enumerate}
\end{methode}

\begin{exemplebox}[Cas pratique : jeune conducteur de moto-taxi à Douala]
\textbf{Profil.} Homme, 24 ans, 1,72 m, 92 kg (IMC 31,1), tour de taille 102 cm. Apports estimés \num{13500}~\SI{}{kJ/j} (riz, huile de palme, boissons sucrées, beignets). Dépenses \num{10500}~\SI{}{kJ/j} (métabolisme de base \num{7500} + activité \num{3000}). $\Delta E = +3000$~\SI{}{kJ/j}.
\textbf{Biologie.} Glycémie 1,18 g/L, TG 2,1 g/L, HDL 0,35 g/L, CRPus 4,2 mg/L, leptine 45 ng/mL.
\textbf{Interprétation.} Syndrome métabolique (obésité androïde, dyslipidémie, hyperglycémie, inflammation). Résistance à l'insuline et à la leptine.
\textbf{Plan correctif.} Réduire huile de palme (remplacer par huile d'arachide non raffinée), supprimer boissons sucrées, introduire légumes-feuilles (feuilles d'okra, amarante) et légumineuses (haricots, niébé). Marche rapide 30 min 5 fois/semaine. Objectif : $\Delta E \approx -500$~\SI{}{kJ/j} $\rightarrow$ perte \num{0,5} kg/semaine. Réévaluation à 3 mois.
\end{exemplebox}

\begin{savaistu}[Le paradoxe de l'obésité au Cameroun]
Dans certaines régions, on observe simultanément dans un même ménage une mère en surpoids (apports excessifs en lipides raffinés, sédentarité) et un enfant en retard de croissance (apports insuffisants en protéines et micronutriments, infections à répétition). Ce \emph{double fardeau nutritionnel} illustre la complexité de la transition nutritionnelle et la nécessité d'interventions multisectorielles (agriculture, éducation, santé, urbanisme).
\end{savaistu}

\coursec{Adaptations métaboliques à l'effort et à la privation alimentaire}

Dans la section précédente, nous avons vu qu'un déséquilibre durable entre apports et dépenses énergétiques conduit à des conséquences physiologiques graves : obésité et ses complications en cas de bilan positif, maigreur, marasme ou kwashiorkor en cas de bilan négatif. Mais l'organisme n'est pas passif face aux variations de la dépense ou de l'apport énergétique. Observe un coureur de fond sur la route Yaoundé--Nsimalen un samedi matin, ou un cultivateur de la région de l'Ouest qui travaille sa plantation toute la journée sans repas de midi : leur organisme \emph{adapte} en permanence les voies métaboliques utilisées pour produire l'ATP nécessaire. De même, lors d'une privation alimentaire (jeûne rituel, disette, grève de la faim), le corps met en place des stratégies de survie remarquables. C'est l'objet de cette section : comprendre \emph{comment} l'organisme choisit et ajuste ses carburants énergétiques selon l'intensité de l'effort et la durée du jeûne.

\courssub{Le choix des carburants énergétiques selon l'intensité de l'effort}

\textbf{Situation-problème.} Deux élèves de Première D font un test d'effort progressif sur cycloergomètre au laboratoire de physiologie. On mesure en continu leur consommation de dioxygène ($V_{\ce{O2}}$) et leur production de dioxyde de carbone ($V_{\ce{CO2}}$). Au repos, le rapport $V_{\ce{CO2}}/V_{\ce{O2}}$ vaut environ 0,78 ; à effort maximal, il dépasse 1,1. Que signifie ce rapport, et pourquoi varie-t-il ?

\begin{propriete}[Le quotient respiratoire (QR), signature du carburant utilisé]
Le \cle{quotient respiratoire} est le rapport entre le volume de \ce{CO2} produit et le volume de \ce{O2} consommé par l'organisme :
\[ QR = \frac{V_{\ce{CO2}}}{V_{\ce{O2}}} \]
Sa valeur dépend du substrat oxydé :
\begin{itemize}
\item oxydation des \textbf{glucides purs} : \ce{C6H12O6 + 6O2 -> 6CO2 + 6H2O}, donc $QR = 6/6 = 1{,}0$ ;
\item oxydation des \textbf{lipides} (ex. palmitate) : \ce{C16H32O2 + 23O2 -> 16CO2 + 16H2O}, donc $QR = 16/23 \approx 0{,}7$ ;
\item alimentation \textbf{mixte} au repos : $QR \approx 0{,}85$.
\end{itemize}
Un $QR > 1$ à l'effort signale une production supplémentaire de \ce{CO2} par tamponnage de l'acide lactique : c'est le signe d'un métabolisme anaérobie important.
\end{propriete}

\begin{methode}[Interpréter des courbes de $V_{\ce{O2}}$ et $V_{\ce{CO2}}$ pour identifier le substrat principal]
\begin{enumerate}
\item Relever sur les courbes, à l'instant considéré, les valeurs de $V_{\ce{CO2}}$ et $V_{\ce{O2}}$ (mêmes unités, même instant).
\item Calculer le quotient $QR = V_{\ce{CO2}}/V_{\ce{O2}}$.
\item Comparer aux valeurs de référence : proche de 0,7 $\rightarrow$ lipides dominants ; proche de 1,0 $\rightarrow$ glucides dominants ; intermédiaire $\rightarrow$ mélange des deux.
\item Conclure en reliant le substrat identifié au contexte (repos, intensité d'effort, jeûne) et vérifier la cohérence avec les réserves de l'organisme.
\end{enumerate}
\end{methode}

\courssub{Au repos et à effort faible : le règne des acides gras}

Au repos, la dépense énergétique est faible et le temps ne manque pas : l'organisme privilégie alors l'oxydation des \cle{acides gras}, carburant très riche en énergie (environ \SI{38}{kJ/g} contre \SI{17}{kJ/g} pour les glucides) et stocké en grande quantité dans le tissu adipeux. Les acides gras, libérés des triglycérides par la \textbf{lipolyse}, pénètrent dans les mitochondries des cellules musculaires et y subissent la \cle{$\beta$-oxydation} : ils sont coupés en fragments à deux atomes de carbone (acétyl-CoA) qui alimentent le cycle de Krebs, avec production abondante d'ATP en présence d'\ce{O2}.

Cette stratégie a une fonction vitale : \textbf{épargner le glucose}. En effet, le cerveau (système nerveux central) et les globules rouges dépendent quasi exclusivement du glucose en conditions normales ; les réserves de glycogène étant limitées (environ 400 à \SI{500}{g} chez l'adulte), il faut les préserver. Au repos, le QR d'environ 0,78--0,85 traduit ce mélange où les lipides dominent.

\begin{attention}[Piège classique]
Ne crois pas que les graisses sont le \emph{seul} carburant au repos ! Même au repos, le cerveau consomme environ \SI{120}{g} de glucose par jour. Les lipides dominent quantitativement, mais le glucose reste \emph{indispensable} pour le système nerveux central et les hématies, qui ne peuvent pas utiliser les acides gras (ceux-ci ne traversent pas la barrière hémato-encéphalique, et les globules rouges n'ont pas de mitochondries).
\end{attention}

\courssub{Effort modéré (50--65 \% de la VO\textsubscript{2}max) : montée en puissance des glucides}

Lorsque l'intensité de l'effort augmente (footing soutenu, match de football, travail agricole soutenu), la demande en ATP s'accélère. Or l'oxydation des lipides est \emph{lente} : elle exige de nombreuses étapes et beaucoup d'\ce{O2} par unité d'ATP. Le muscle recrute alors de plus en plus le \textbf{glucose sanguin} et le \textbf{glycogène musculaire} stocké localement, dont la dégradation glycolytique fournit l'ATP plus rapidement. C'est dans cette zone d'intensité que l'oxydation des glucides atteint son \textbf{pic} en valeur absolue. Le QR remonte vers 0,85--0,95.

\courssub{Effort intense (au-delà de 80 \% de la VO\textsubscript{2}max) : l'anaérobie prend le relais}

Au-delà d'environ 80 \% de la $\dot{V}\ce{O2}$max, l'apport d'oxygène aux muscles ne suffit plus à couvrir les besoins. Le muscle recrute massivement les \cle{fibres rapides} (fibres de type II, glycolytiques), pauvres en mitochondries mais capables de contractions puissantes. La \textbf{glycolyse anaérobie} devient la voie majeure : le glucose est dégradé jusqu'au pyruvate, qui, faute d'\ce{O2} suffisant pour entrer dans le cycle de Krebs, est réduit en \cle{lactate} :
\[ \ce{C6H12O6 -> 2 CH3-CHOH-COOH} + 2\,\text{ATP} \]
Le rendement est faible (2 ATP par glucose contre environ 30 à 32 en aérobie), mais la vitesse de production est élevée. L'accumulation de lactate et d'ions \ce{H+} provoque la fatigue musculaire (sensation de brûlure). Après l'effort, la consommation d'\ce{O2} reste élevée pendant un certain temps : c'est la \cle{dette d'oxygène}, qui sert à reconvertir le lactate en glucose (foie, cycle de Cori), à reconstituer les réserves d'ATP et de phosphocréatine et à réoxygéner la myoglobine.

\begin{definition}[Seuil aérobie (seuil lactique)]
Le \cle{seuil aérobie} (ou seuil lactique) est l'intensité d'effort au-delà de laquelle la \textbf{production de lactate par les muscles dépasse sa capacité d'élimination} (réutilisation, oxydation, captation hépatique). Le lactate sanguin s'élève alors brutalement, l'acide lactique est tamponné par les bicarbonates avec libération de \ce{CO2} supplémentaire, et le $QR$ devient supérieur à 1. Au-delà de ce seuil, l'effort ne peut être maintenu que peu de temps.
\end{definition}

\begin{popfigure}
\begin{center}
\begin{tikzpicture}
\begin{axis}[
  width=0.92\linewidth, height=7.2cm,
  xlabel={Intensité de l'effort (\% $\dot{V}\ce{O2}$max)},
  ylabel={$\dot{V}\ce{O2}$ (L/min)},
  ylabel style={color=popblue},
  axis y line*=left, axis x line*=bottom,
  xmin=0, xmax=100, ymin=0, ymax=4.2,
  xtick={0,20,40,60,80,100},
  legend style={at={(0.03,0.97)}, anchor=north west, font=\small},
  grid=both, grid style={popline!40},
]
\addplot[popcurve, color=popblue, very thick, domain=0:100, samples=100] {0.3 + 0.035*x};
\addlegendentry{$\dot{V}\ce{O2}$ (L/min)}
\addplot[popcurve, color=poporange, very thick, domain=0:100, samples=100] {0.24 + 0.026*x + 0.00016*x^2};
\addlegendentry{$\dot{V}\ce{CO2}$ (L/min)}
\draw[dashed, popdark] (axis cs:80,0) -- (axis cs:80,4.2) node[pos=0.95, anchor=south east, font=\small\itshape, text=popdark]{seuil aérobie};
\end{axis}
\begin{axis}[
  width=0.92\linewidth, height=7.2cm,
  axis y line*=right, axis x line=none,
  ylabel={QR}, ylabel style={color=popgreen},
  xmin=0, xmax=100, ymin=0.6, ymax=1.2,
  ytick={0.7,0.85,1.0},
  legend style={at={(0.03,0.80)}, anchor=north west, font=\small},
]
\addplot[popcurve, color=popgreen, very thick, dashed, domain=0:100, samples=100] {0.78 + 0.0012*x + 0.000014*x^2};
\addlegendentry{QR = $V_{\ce{CO2}}/V_{\ce{O2}}$}
\end{axis}
\end{tikzpicture}
\end{center}
\legende{Évolution de la consommation d'\ce{O2}, de la production de \ce{CO2} (axe de gauche) et du quotient respiratoire (axe de droite) en fonction de l'intensité de l'effort. Au repos, $QR \approx 0{,}78$ (lipides dominants) ; le QR remonte vers 1 (glucides dominants) puis dépasse 1 au-delà du seuil aérobie (tamponnage de l'acide lactique).}
\end{popfigure}

\begin{exemplebox}[Exemple chiffré : le coureur de 70 kg]
Un coureur de \SI{70}{kg} a une $\dot{V}\ce{O2}$max de \SI{55}{mL/kg/min}, soit $\dot{V}\ce{O2}\text{max} = 55 \times 70 = \SI{3850}{mL/min} \approx \SI{3,85}{L/min}$. À 60 \% de sa $\dot{V}\ce{O2}$max :
\[ \dot{V}\ce{O2} = 0{,}60 \times 3{,}85 \approx \SI{2,3}{L/min} \]
Le QR mesuré est d'environ 0,88. Les tables de correspondance QR/substrats indiquent alors une oxydation d'environ \textbf{55 \% de glucides et 45 \% de lipides}. Interprétation : à cette intensité modérée, les deux carburants contribuent presque également ; si le coureur accélère vers 90 \% de sa $\dot{V}\ce{O2}$max, le QR approchera 1 et les glucides deviendront quasi exclusifs, épuisant rapidement le glycogène.
\end{exemplebox}

\courssub{La régulation hormonale du métabolisme énergétique}

Le choix des carburants n'est pas laissé au hasard : il est piloté par un système hormonal précis, dont le chef d'orchestre est le pancréas endocrine (îlots de Langerhans).

\begin{center}
\begin{tabularx}{\linewidth}{|l|X|X|}
\hline
\rowcolor{popblueL} \textbf{Hormone} & \textbf{Glande / situation de sécrétion} & \textbf{Effet métabolique principal} \\
\hline
\cle{Insuline} & Cellules $\beta$ du pancréas ; sécrétée après un repas (hyperglycémie) & Hormone de \textbf{stockage} : entrée du glucose dans les cellules, glycogénèse, lipogenèse ; inhibe la lipolyse \\
\hline
\cle{Glucagon} & Cellules $\alpha$ du pancréas ; sécrété à jeun (hypoglycémie) & Hormone de \textbf{mobilisation} : glycogénolyse hépatique, néoglucogenèse, lipolyse \\
\hline
\cle{Adrénaline} & Médullo-surrénales ; effort, stress, hypoglycémie & Mobilisation rapide : glycogénolyse musculaire et hépatique, lipolyse ; prépare à l'action \\
\hline
\cle{Cortisol} & Cortico-surrénales ; stress prolongé, jeûne prolongé & \textbf{Néoglucogenèse} à partir des acides aminés (protéolyse musculaire), lipolyse ; adaptation aux situations durables \\
\hline
\end{tabularx}
\end{center}

Retiens la logique d'ensemble : l'\textbf{insuline} gère l'abondance (elle range), le \textbf{glucagon} et l'\textbf{adrénaline} gèrent le besoin immédiat (ils sortent les réserves), le \textbf{cortisol} gère la pénurie durable (il fabrique du glucose neuf, même au prix des protéines).

\courssub{Les adaptations métaboliques à la privation alimentaire}

\textbf{Situation-problème.} Lors de certaines périodes de jeûne religieux, ou dans des situations de disette alimentaire que connaissent malheureusement certaines zones du pays, une personne peut rester des jours sans manger tout en continuant à marcher, penser, travailler. D'où vient l'énergie ? Et surtout, comment le cerveau, grand consommateur de glucose, survit-il alors que les réserves de glycogène ne couvrent qu'environ une journée de besoins ?

\courssub{Le jeûne court (moins de 24 h) : la glycogénolyse hépatique}

Dès les premières heures qui suivent le dernier repas, la glycémie tend à baisser. Sous l'action du \textbf{glucagon} (et de l'adrénaline), le foie dégrade son glycogène : c'est la \cle{glycogénolyse hépatique}, qui libère du glucose dans le sang pour maintenir la glycémie autour de \SI{1}{g/L}. Le muscle, lui, garde jalousement son glycogène (il ne possède pas l'enzyme glucose-6-phosphatase permettant d'exporter le glucose). En parallèle, la lipolyse s'intensifie : les muscles et la plupart des organes passent progressivement aux acides gras, épargnant le glucose pour le cerveau.

\courssub{Le jeûne prolongé (au-delà de 24 h) : néoglucogenèse et cétogenèse}

Les réserves hépatiques de glycogène s'épuisent en 18 à 24 heures. Deux voies de secours prennent alors le relais :

\begin{itemize}
\item La \cle{néoglucogenèse} : le foie (et, à terme, le rein) \emph{fabrique du glucose neuf} à partir de précurseurs non glucidiques : le \textbf{glycérol} issu de la lipolyse, le \textbf{lactate} (cycle de Cori) et surtout les \textbf{acides aminés} provenant de la dégradation des protéines musculaires, sous l'action du cortisol. C'est ce qui explique la fonte musculaire observée lors des dénutritions prolongées.
\item La \cle{cétogenèse} : l'oxydation massive des acides gras dans le foie produit un excès d'acétyl-CoA, transformé en \textbf{corps cétoniques} (acétoacétate, $\beta$-hydroxybutyrate, acétone). Après quelques jours de jeûne, le cerveau \emph{adapte son métabolisme} et utilise les corps cétoniques pour couvrir jusqu'aux deux tiers de ses besoins : adaptation capitale, car elle réduit la néoglucogenèse et donc \textbf{préserve les protéines musculaires}. L'haleine odorante (odeur d'acétone) des personnes en jeûne prolongé en est un signe clinique.
\end{itemize}

\begin{popfigure}
\begin{center}
\begin{tikzpicture}[
  node distance=0.55cm and 1.1cm,
  etape/.style={rectangle, rounded corners=4pt, draw=popblue, fill=popblueL, thick, align=center, font=\small, minimum width=3.1cm, minimum height=1.15cm},
  voie/.style={-{Stealth[length=3mm]}, very thick, popdark},
  temps/.style={font=\small\itshape, text=popmuted}
]
\node[etape, align=center] (repas){\textbf{Dernier repas}\\ absorption intestinale\\ insuline $\uparrow$ (stockage)};
\node[etape, below=of repas, fill=popgreenL, draw=popgreen, align=center] (glyco){\textbf{Glycogénolyse hépatique}\\ glycogène du foie $\rightarrow$ glucose\\ glucagon $\uparrow$};
\node[etape, below=of glyco, fill=poporangeL, draw=poporange, align=center] (neo){\textbf{Néoglucogenèse}\\ acides aminés, glycérol, lactate\\ $\rightarrow$ glucose (foie, rein) ; cortisol $\uparrow$};
\node[etape, below=of neo, fill=poppurpleL, draw=poppurple, align=center] (ceto){\textbf{Cétogenèse}\\ acides gras $\rightarrow$ corps cétoniques\\ carburant de secours du cerveau};
\draw[voie] (repas) -- (glyco) node[midway, right, temps]{0 -- 6 h : glycémie maintenue};
\draw[voie] (glyco) -- (neo) node[midway, right, temps]{6 -- 24 h : épuisement du glycogène};
\draw[voie] (neo) -- (ceto) node[midway, right, temps]{$> 24$ h : jeûne prolongé};
\node[etape, right=1.6cm of glyco, fill=popgoldL, draw=popgold, align=center] (lipo){\textbf{Lipolyse} (continue)\\ triglycérides $\rightarrow$ acides gras\\ + glycérol : carburant des muscles};
\draw[voie, popgold] (lipo.south) |- (neo.east);
\draw[voie, popgold] (lipo.south) |- (ceto.east);
\end{tikzpicture}
\end{center}
\legende{Succession des voies métaboliques activées au cours du jeûne. La glycogénolyse hépatique assure les premières heures ; la néoglucogenèse prend le relais à partir de précurseurs non glucidiques ; la cétogenèse fournit des corps cétoniques qui alimentent le cerveau lors du jeûne prolongé, épargnant ainsi les protéines. La lipolyse fonctionne en toile de fond dès les premières heures.}
\end{popfigure}

\begin{experience}[Démarche d'investigation : comment sait-on que le cerveau jeune utilise les corps cétoniques ?]
\textbf{Observation} : des patients en jeûne prolongé restent conscients et lucides alors que leur glycémie est basse et leur glycogène épuisé. \textbf{Hypothèse} : le cerveau utilise un carburant de substitution. \textbf{Données expérimentales} (études classiques de type Cahill) : le dosage des corps cétoniques sanguins montre leur multiplication par 20 à 50 après une semaine de jeûne ; la mesure des prélèvements artério-veineux cérébraux montre que le cerveau \emph{consomme} ces corps cétoniques, et sa consommation de glucose chute d'environ deux tiers. \textbf{Interprétation} : les corps cétoniques traversent la barrière hémato-encéphalique et sont oxydés par les neurones. \textbf{Conclusion} : la cétogenèse hépatique est une adaptation de survie qui préserve à la fois le cerveau et les protéines de l'organisme.
\end{experience}

\begin{savaistu}[L'entraînement repousse les limites]
L'entraînement d'endurance régulier (course, football, cyclisme — pensons aux coureurs camerounais qui s'entraînent sur les pentes du Mont Cameroun lors de la célèbre course de l'Espoir) provoque des adaptations remarquables : augmentation de la \textbf{densité mitochondriale} dans les fibres musculaires, des stocks de glycogène et de la \textbf{capacité d'oxydation des lipides}. Résultat : le sportif entraîné utilise davantage les graisses à intensité égale, épargne son glycogène et \textbf{retarde son seuil aérobie} — il peut courir plus vite, plus longtemps, avant d'accumuler du lactate. C'est la preuve vivante que le métabolisme énergétique est \emph{plastique} : il s'adapte à ce qu'on lui demande.
\end{savaistu}

\begin{exoresolu}[Identifier le carburant d'un sportif à partir de données respiratoires]
\textbf{Énoncé.} Lors d'un test d'effort, un élève de \SI{65}{kg} présente à un palier donné : $V_{\ce{CO2}} = \SI{2,04}{L/min}$ et $V_{\ce{O2}} = \SI{2,40}{L/min}$. 1) Calculer le QR à ce palier. 2) En déduire le ou les substrats majoritairement oxydés. 3) Que prévoit-on pour le QR si l'intensité augmente encore au-delà du seuil aérobie ? Justifier.
\tcblower
\textbf{Solution détaillée.}
1) $QR = \dfrac{V_{\ce{CO2}}}{V_{\ce{O2}}} = \dfrac{2{,}04}{2{,}40} = 0{,}85$.

2) Un QR de 0,85 correspond à une oxydation \textbf{mixte} : glucides et lipides contribuent de façon comparable (environ 50 \% chacun). C'est typique d'un effort d'intensité modérée (50--65 \% de la $\dot{V}\ce{O2}$max).

3) Au-delà du seuil aérobie, le QR va \textbf{dépasser 1}. En effet, la glycolyse anaérobie produit du lactate ; l'acide lactique est tamponné par les bicarbonates sanguins selon $\ce{H+} + \ce{HCO3-} \rightarrow \ce{CO2} + \ce{H2O}$, ce qui libère du \ce{CO2} \emph{supplémentaire}, non issu de l'oxydation. Le numérateur $V_{\ce{CO2}}$ augmente donc plus vite que $V_{\ce{O2}}$, et le QR dépasse 1 : c'est le signal que le métabolisme anaérobie domine et que la fatigue musculaire est proche.
\end{exoresolu}

\begin{exercice}[Jeûne et adaptation métabolique]
Un adulte de \SI{70}{kg} jeûne complètement pendant 3 jours (apport hydrique maintenu). 1) Quelle voie métabolique maintient la glycémie pendant les 20 premières heures ? Quelle hormone la commande ? 2) Pourquoi la néoglucogenèse devient-elle ensuite indispensable, et quels en sont les précurseurs ? 3) Expliquer pourquoi l'apparition des corps cétoniques est bénéfique pour la survie à long terme. 4) Quel signe clinique de l'haleine peut trahir la cétogenèse ?
\end{exercice}

\begin{corrige}[Exercice : Jeûne et adaptation métabolique]
1) Pendant les 20 premières heures, la \textbf{glycogénolyse hépatique} (dégradation du glycogène du foie) maintient la glycémie ; elle est commandée par le \textbf{glucagon} (sécrété en réponse à l'hypoglycémie), relayé par l'adrénaline.

2) Les réserves de glycogène hépatique sont épuisées en 18--24 h, or le cerveau et les globules rouges exigent du glucose en permanence. La \textbf{néoglucogenèse} devient donc indispensable : le foie synthétise du glucose à partir du \textbf{glycérol} (lipolyse), du \textbf{lactate} et des \textbf{acides aminés} issus de la protéolyse musculaire, sous l'action du cortisol.

3) Les corps cétoniques, produits par le foie à partir des acides gras, deviennent utilisables par le \textbf{cerveau} après quelques jours. Ils couvrent jusqu'aux deux tiers de ses besoins, ce qui réduit la néoglucogenèse glucidique et donc \textbf{limite la fonte des protéines musculaires} : l'organisme survit plus longtemps en puisant dans ses graisses, quasi inépuisables, plutôt que dans ses muscles.

4) L'\textbf{acétone}, corps cétonique volatil éliminé par les poumons, donne à l'haleine une odeur caractéristique (fruitée, de « pomme reinette ») : c'est l'haleine acétonémique du jeûne prolongé (retrouvée aussi dans le diabète déséquilibré).
\end{corrige}

\begin{aretenir}[L'essentiel de la section]
\begin{itemize}
\item Le carburant utilisé dépend de l'intensité de l'effort : \textbf{lipides} au repos (β-oxydation, épargne du glucose pour le cerveau), \textbf{glucides croissants} à effort modéré, \textbf{glycolyse anaérobie et lactate} à effort intense, avec apparition d'une \textbf{dette d'oxygène}.
\item Le \textbf{quotient respiratoire} $QR = V_{\ce{CO2}}/V_{\ce{O2}}$ identifie le substrat : 0,7 (lipides), 0,85 (mixte), 1,0 (glucides), $> 1$ au-delà du \textbf{seuil aérobie}.
\item La régulation hormonale : \textbf{insuline} = stockage ; \textbf{glucagon} et \textbf{adrénaline} = mobilisation ; \textbf{cortisol} = néoglucogenèse.
\item Au jeûne : \textbf{glycogénolyse hépatique} ($< 24$ h), puis \textbf{néoglucogenèse} (glycérol, lactate, acides aminés) et \textbf{cétogenèse} (corps cétoniques pour le cerveau), qui épargne les protéines.
\end{itemize}
\end{aretenir}

\coursec{Mesures correctives, prévention et éducation nutritionnelle}

Dans la section précédente, nous avons vu que l'organisme humain possède de remarquables \cle{adaptations métaboliques} : il mobilise ses réserves glycogéniques puis lipidiques à l'effort, et ralentit son \cle{métabolisme de base} en cas de privation alimentaire prolongée. Ces adaptations, héritées de millions d'années d'évolution, expliquent un paradoxe que tu as peut-être observé autour de toi : certaines personnes suivent un « régime » très strict, perdent quelques kilogrammes... puis en reprennent davantage quelques mois plus tard. À l'inverse, d'autres modifient durablement leurs habitudes et stabilisent leur poids. \textbf{Quelles mesures correctives sont réellement efficaces, et comment les mettre en œuvre à l'échelle individuelle et communautaire ?} C'est l'objet de cette dernière section, qui fait de toi un acteur de santé publique capable de proposer des solutions concrètes, adaptées au contexte camerounais.

\courssub{Principes du rééquilibrage alimentaire}

\textbf{Situation concrète.} Mme Etoundi, 42 ans, commerçante à Yaoundé, consulte au Centre Hospitalier d'Essos : son médecin diagnostique un surpoids (\cle{IMC} = 29 kg/m²) et une glycémie à jeun limite. Son alimentation quotidienne : beignets-haricots le matin, un grand plat de ndolé très huilé avec du bâton de manioc le midi, du riz sauté à l'huile de palme le soir, et deux bouteilles de boisson sucrée dans la journée. Faut-il lui interdire de manger ? Non : il faut \emph{rééquilibrer}, non supprimer.

\begin{definition}[Éducation nutritionnelle]
L'\cle{éducation nutritionnelle} est un processus d'acquisition de connaissances, de compétences et d'attitudes permettant à un individu ou à une communauté de choisir une alimentation adaptée à ses besoins physiologiques, dans le respect de sa culture et de ses ressources. Elle vise des changements de comportement durables, et non des restrictions temporaires.
\end{definition}

Le rééquilibrage repose sur trois principes complémentaires :

\begin{enumerate}
\item \textbf{Réduire la densité énergétique} des repas : diminuer les huiles (une cuillère à soupe d'huile de palme $\approx$ 450 kJ), supprimer ou limiter les boissons sucrées (une bouteille de 50 cL $\approx$ 900 kJ, sans aucune satiété), réduire les fritures au profit des cuissons à l'eau, à la vapeur ou braisées.
\item \textbf{Augmenter la part des fibres} : légumes-feuilles (ndolé, eru, folong, kelen-kelen), fruits locaux (mangue, goyave, papaye), céréales complètes ou peu raffinées, légumineuses (niébé, haricots, arachides). Les fibres augmentent le volume du bol alimentaire sans apporter beaucoup d'énergie, accélèrent la satiété et ralentissent l'absorption du glucose.
\item \textbf{Répartir les apports sur 3 à 4 repas} par jour, sans sauter le petit-déjeuner : une alimentation concentrée en un seul gros repas favorise le stockage lipidique et les fringales.
\end{enumerate}

\begin{popfigure}
\begin{center}
\begin{tikzpicture}[scale=1]
% Pyramide alimentaire en couches (trapèzes)
\fill[popgreen!70] (0,0) -- (6,0) -- (5.1,1.2) -- (0.9,1.2) -- cycle;
\fill[popgold!75] (0.9,1.2) -- (5.1,1.2) -- (4.2,2.4) -- (1.8,2.4) -- cycle;
\fill[poporange!75] (1.8,2.4) -- (4.2,2.4) -- (3.6,3.4) -- (2.4,3.4) -- cycle;
\fill[poppink!70] (2.4,3.4) -- (3.6,3.4) -- (3.3,4.2) -- (2.7,4.2) -- cycle;
\fill[popmuted!60] (2.7,4.2) -- (3.3,4.2) -- (3,5) -- cycle;
% Étiquettes à droite
\draw[thin,popdark] (6,0.6) -- (7.1,0.6) node[right,align=left,font=\small] {\textbf{Base :} céréales et tubercules (maïs, mil,\\ riz complet, igname, plantain) : 4 à 6 portions/j};
\draw[thin,popdark] (5.1,1.8) -- (7.1,1.8) node[right,align=left,font=\small] {\textbf{Légumes et fruits} (ndolé, folong,\\ goyave, papaye) : 5 portions/j minimum};
\draw[thin,popdark] (4.2,2.9) -- (7.1,2.9) node[right,align=left,font=\small] {\textbf{Protéines} (niébé, haricots, poisson,\\ viande maigre, œufs) : 2 à 3 portions/j};
\draw[thin,popdark] (3.6,3.8) -- (7.1,3.8) node[right,align=left,font=\small] {\textbf{Laitages et corps gras} : avec modération\\ (1 à 2 cuillères à soupe d'huile/j)};
\draw[thin,popdark] (3.15,4.6) -- (7.1,4.6) node[right,align=left,font=\small] {\textbf{Sucres, boissons sucrées, fritures :}\\ occasionnels, en très petites quantités};
% Titre
\node[font=\small\bfseries,popdark] at (3,-0.5) {Pyramide alimentaire adaptée au contexte camerounais};
\end{tikzpicture}
\end{center}
\legende{Pyramide alimentaire adaptée aux aliments locaux du Cameroun. La base (céréales et tubercules peu raffinés) fournit l'essentiel de l'énergie ; le sommet (sucres, huiles, fritures) doit rester exceptionnel. Une portion $\approx$ 80 à 100 g.}
\end{popfigure}

\begin{attention}[Les régimes drastiques sont contre-productifs]
Un régime très hypocalorique ($< \num{1200}$ kcal/j, soit environ \num{5000} kJ/j) déclenche les adaptations à la privation étudiées en section 4 : fonte musculaire (les protéines sont dégradées pour la \cle{néoglucogenèse}), \textbf{ralentissement du métabolisme de base} (diminution de la \cle{thermogénèse} et de l'activité thyroïdienne), fatigue et irritabilité. À l'arrêt du régime, l'organisme « économisé » stocke d'autant plus facilement : c'est l'\cle{effet yo-yo}, avec reprise pondérale souvent supérieure au poids initial. Un déficit modéré et durable est toujours préférable.
\end{attention}

\begin{attention}[Méfie-toi des « brûle-graisse »]
Les compléments alimentaires présentés comme « brûle-graisse », vendus sans contrôle sur les marchés ou sur les réseaux sociaux, ne font l'objet d'aucune évaluation rigoureuse. Certains contiennent des stimulants (amphétamines-like, extraits thyroïdiens, diurétiques) provoquant tachycardie, hypertension, troubles du rythme cardiaque, voire hépatites toxiques. Aucun produit ne « brûle » les graisses : seul un déficit énergétique obtenu par l'alimentation réelle et l'activité physique mobilise durablement les réserves du \cle{tissu adipeux}.
\end{attention}

\courssub{L'activité physique : l'autre volet du bilan}

Corriger un déséquilibre énergétique ne consiste pas seulement à manger différemment : il faut aussi \textbf{augmenter les dépenses}. L'Organisation mondiale de la santé (OMS) recommande aux enfants et adolescents \textbf{au moins 60 minutes par jour d'activité physique d'intensité modérée à vigoureuse}, et aux adultes au moins 150 minutes par semaine.

L'activité physique présente un double avantage métabolique :
\begin{itemize}
\item elle augmente directement la \cle{dépense énergétique} (30 min de marche rapide $\approx$ 600 kJ ; 1 h de football $\approx$ \num{1800} kJ) ;
\item elle préserve et développe la \textbf{masse musculaire}, tissu très consommateur d'énergie, ce qui maintient le métabolisme de base et améliore la \cle{sensibilité à l'insuline}.
\end{itemize}

Les activités recommandées sont accessibles à tous, sans matériel coûteux : marche rapide (aller à l'école ou au marché à pied d'un pas soutenu), football, danse traditionnelle (bikutsi, makossa, assiko — une heure de danse équivaut à un footing léger !), vélo, corvées actives, jardinage.

\begin{exemplebox}[Exemple chiffré : un petit changement, un grand effet]
Mme Etoundi applique deux mesures simples :
\begin{itemize}
\item elle réduit de 2 cuillères à soupe l'huile de palme de ses plats : $-900$ kJ/j ;
\item elle marche rapidement 30 min chaque matin : $-600$ kJ/j de dépense supplémentaire.
\end{itemize}
Déficit quotidien : $\Delta E \approx -\num{1500}$ kJ/j. Or 1 kg de tissu adipeux $\approx$ \num{30000} kJ. La perte hebdomadaire attendue est donc :
\[ \frac{1500 \times 7}{30000} \approx 0,35 \ \text{kg/semaine} \]
Soit environ 1,4 kg par mois, \textbf{sans faim, sans régime drastique}, et de façon durable.
\end{exemplebox}

\courssub{Suivi pondéral et comportemental}

Un rééquilibrage ne se pilote pas « au feeling » : il s'évalue avec des outils simples.

\begin{itemize}
\item \textbf{La pesée hebdomadaire} : une fois par semaine, le matin à jeun, dans les mêmes conditions. Peser chaque jour est inutile et anxiogène : les fluctuations hydriques masquent l'évolution réelle des réserves.
\item \textbf{Le journal alimentaire} : noter chaque jour les repas, les grignotages et les boissons. Cet outil révèle les apports « invisibles » (boissons sucrées, beignets entre les repas) qui expliquent souvent à eux seuls le déséquilibre.
\item \textbf{L'auto-évaluation de la satiété} : apprendre à distinguer la faim physiologique (ventre qui gargouille, creux gastrique) de l'envie de manger (stress, ennui, vue d'un aliment). S'arrêter de manger quand la satiété apparaît, environ 15 à 20 minutes après le début du repas — le temps que les signaux hormonaux (\cle{leptine}, \cle{insuline}, \cle{cholécystokinine}) atteignent l'\cle{hypothalamus}.
\end{itemize}

\begin{propriete}[Bénéfice d'une perte de poids modérée]
Une perte de \textbf{5 à 10 \% du poids initial} — par exemple 4 à 8 kg pour une personne de 80 kg — suffit à améliorer significativement la \cle{sensibilité à l'insuline} (diminution de la glycémie à jeun et du risque de diabète de type 2) et le \cle{profil lipidique} (baisse des triglycérides et du LDL-cholestérol, hausse du HDL-cholestérol). Il n'est donc pas nécessaire d'atteindre un « poids idéal » pour obtenir des bénéfices majeurs de santé.
\end{propriete}

\begin{popfigure}
\begin{center}
\begin{tikzpicture}
\begin{axis}[
  width=0.95\linewidth, height=6.5cm,
  xlabel={Durée de suivi (mois)}, ylabel={Masse corporelle (kg)},
  xmin=0, xmax=6, ymin=80, ymax=94,
  xtick={0,1,2,3,4,5,6}, ytick={80,82,84,86,88,90,92,94},
  grid=both, grid style={popline!40},
  legend style={at={(0.98,0.98)},anchor=north east,font=\small,draw=popline},
  axis line style={popdark}, tick label style={font=\small}, label style={font=\small}
]
% Sans intervention : légère hausse
\addplot[color=popink, very thick, mark=*, mark size=1.5pt] coordinates {(0,90) (1,90.3) (2,90.6) (3,91) (4,91.3) (5,91.6) (6,92)};
% Avec intervention : perte progressive puis stabilisation
\addplot[color=popgreen, very thick, mark=square*, mark size=1.5pt] coordinates {(0,90) (1,88.8) (2,87.7) (3,86.7) (4,85.9) (5,85.3) (6,85)};
\legend{Sans intervention, Avec conseil diététique + 3 séances de sport/semaine}
\end{axis}
\end{tikzpicture}
\end{center}
\legende{Évolution de la masse corporelle sur 6 mois chez un adulte en surpoids (90 kg au départ). Sans intervention, la masse augmente lentement (+2 kg). Avec un rééquilibrage alimentaire et 3 séances hebdomadaires d'activité physique, la perte est régulière ($\approx$ 1 kg/mois) puis se stabilise : perte totale de 5 kg, soit 5,6 \% du poids initial — suffisante pour améliorer la sensibilité à l'insuline.}
\end{popfigure}

\begin{methode}[Proposer des mesures correctives individualisées]
Face à un cas de déséquilibre énergétique (exercice type Probatoire ou situation réelle), procède en 5 étapes :
\begin{enumerate}
\item \textbf{Établis le bilan} : calcule les apports (journal alimentaire) et les dépenses (métabolisme de base + activité). Détermine le signe et l'ampleur du déséquilibre $\Delta E$.
\item \textbf{Fixe un objectif raisonnable} : pour une perte de 0,5 kg/semaine, vise $\Delta E \approx -500$ kJ/j par l'alimentation, complété par l'activité physique (au total $-1000$ à $-1500$ kJ/j). Ne descends jamais sous \num{1200} kcal/j d'apports.
\item \textbf{Propose des modifications concrètes et culturellement acceptables} : remplacer plutôt qu'interdire (eau à la place des boissons sucrées, cuisson à l'eau plutôt que friture, ajout de légumes-feuilles).
\item \textbf{Ajoute l'activité physique} : au moins 60 min/j d'activité modérée à vigoureuse, en choisissant des activités accessibles et plaisantes.
\item \textbf{Planifie le suivi} : pesée hebdomadaire, journal alimentaire, réévaluation après 1 mois et ajustement des mesures.
\end{enumerate}
\end{methode}

\courssub{L'approche communautaire : agir au-delà de l'individu}

Les choix alimentaires ne dépendent pas seulement de la volonté individuelle : ils sont façonnés par l'environnement (offre alimentaire, prix, publicité, habitudes familiales). La prévention doit donc être aussi \textbf{collective} :

\begin{itemize}
\item \textbf{Cantines scolaires équilibrées} : proposer des repas combinant une céréale ou un tubercule, une légumineuse (niébé, haricots) et un légume-feuille, plutôt que des beignets et boissons sucrées. Plusieurs programmes de cantines scolaires au Cameroun (régions du Nord et de l'Adamaoua) montrent une amélioration conjointe de l'état nutritionnel et de la concentration en classe.
\item \textbf{Sensibilisation aux étiquettes nutritionnelles} : apprendre à lire la composition des produits emballés (teneur en sucres, en graisses, valeur énergétique pour 100 g) permet de comparer et de choisir en connaissance de cause.
\item \textbf{Valorisation des aliments locaux} : le niébé, le maïs, le mil, les légumes-feuilles (ndolé, eru, folong), l'igname et le plantain sont des aliments de grande qualité nutritionnelle, riches en fibres, protéines végétales et micronutriments, et économiquement accessibles. L'éducation nutritionnelle combat l'idée fausse selon laquelle « manger moderne » signifie consommer des produits industriels sucrés et gras.
\item \textbf{Relais communautaires} : chefs traditionnels, associations de femmes, églises et mosquées, clubs sportifs de quartier peuvent diffuser les messages de prévention et organiser des activités physiques collectives.
\end{itemize}

\begin{savaistu}[Le microbiote intestinal, allié méconnu de ton bilan énergétique]
Ton intestin héberge environ $10^{14}$ bactéries : le \cle{microbiote intestinal}. Or toutes les flores ne se valent pas : des études ont montré que, à alimentation identique, deux personnes peuvent extraire des quantités d'énergie différentes de leurs repas selon la composition de leur microbiote. Certaines bactéries fermentent les fibres en \textbf{acides gras à chaîne courte} (butyrate, propionate, acétate) qui protègent la muqueuse intestinale et améliorent la sensibilité à l'insuline. Une alimentation riche en \cle{prébiotiques} — fibres nourrissant les « bonnes » bactéries — comme la banane plantain, l'igname, le niébé ou l'oignon, favorise un profil métabolique sain. Manger local et fibreux, c'est aussi nourrir son microbiote !
\end{savaistu}

\begin{exoresolu}[Corriger le déséquilibre d'un lycéen]
\textbf{Énoncé.} Kevin, 16 ans, en Première D, pèse 78 kg pour 1,70 m. Son journal alimentaire sur 3 jours donne un apport moyen de \num{12500} kJ/j, dont \num{1800} kJ/j de boissons sucrées. Ses dépenses sont estimées à \num{10800} kJ/j (métabolisme de base \num{7200} kJ/j + activité \num{3600} kJ/j). (1) Calcule son déséquilibre énergétique quotidien et prédis son évolution pondérale en 6 mois (1 kg de tissu adipeux $\approx$ \num{30000} kJ). (2) Propose deux mesures correctives chiffrées visant une perte de 0,5 kg/semaine.

\tcblower
\textbf{Solution détaillée.}

\textbf{(1) Bilan énergétique.}
\[ \Delta E = \text{apports} - \text{dépenses} = 12500 - 10800 = +1700 \ \text{kJ/j} \]
Le bilan est \textbf{positif} : Kevin stocke l'excédent sous forme de triglycérides dans le tissu adipeux. Sur 6 mois ($\approx$ 180 jours) :
\[ \text{Excédent stocké} = 1700 \times 180 = \num{306000} \ \text{kJ} \quad\Rightarrow\quad \text{gain} \approx \frac{306000}{30000} \approx 10 \ \text{kg} \]
Sans intervention, Kevin atteindrait environ 88 kg, avec un risque accru de diabète de type 2 et de stéatose hépatique.

\textbf{(2) Mesures correctives.} Objectif : perte de 0,5 kg/semaine, soit un déficit de $0,5 \times 30000 = \num{15000}$ kJ/semaine $\approx$ \num{2150} kJ/j. Proposition :
\begin{itemize}
\item remplacer les boissons sucrées par de l'eau : $-1800$ kJ/j ;
\item ajouter 45 min de football ou de marche rapide par jour : $\approx -900$ kJ/j de dépense supplémentaire.
\end{itemize}
Nouveau bilan : $\Delta E = (12500 - 1800) - (10800 + 900) = -1000$ kJ/j... insuffisant. On ajoute alors une réduction des fritures et de l'huile ($-1200$ kJ/j) : $\Delta E = (12500 - 1800 - 1200) - 11700 = -2200$ kJ/j, ce qui correspond exactement à l'objectif ($\approx$ 0,5 kg/semaine). Les apports restent à \num{9500} kJ/j ($\approx$ \num{2270} kcal/j), très au-dessus du seuil dangereux de \num{1200} kcal/j : la démarche est sûre et durable.
\end{exoresolu}

\begin{aretenir}[L'essentiel de la section]
\begin{itemize}
\item Le rééquilibrage alimentaire repose sur trois principes : réduire la densité énergétique (huiles, boissons sucrées, fritures), augmenter les fibres (légumes, fruits, céréales complètes), répartir les apports sur 3 à 4 repas.
\item L'OMS recommande $\geq$ 60 min/j d'activité physique modérée à vigoureuse : marche rapide, football, danse traditionnelle sont accessibles à tous.
\item Le suivi repose sur la pesée hebdomadaire, le journal alimentaire et l'écoute de la satiété ; une perte de 5 à 10 \% du poids initial améliore déjà nettement la santé métabolique.
\item Les régimes drastiques ($< \num{1200}$ kcal/j) et les « brûle-graisse » sont dangereux et inefficaces à long terme (fonte musculaire, effet yo-yo).
\item La prévention est aussi collective : cantines équilibrées, lecture des étiquettes, valorisation des aliments locaux (niébé, maïs, légumes-feuilles, plantain).
\end{itemize}
\end{aretenir}

\begin{exercice}[À toi de jouer]
Une ménagère de Bafoussam consomme chaque jour 3 cuillères à soupe d'huile rouge en excès ($\approx$ \num{1350} kJ) et ne pratique aucune activité physique. (1) Calcule le gain de masse attendu en un an si rien ne change. (2) Élabore un programme correctif complet (alimentation + activité + suivi) visant une perte de 0,5 kg/semaine, en utilisant uniquement des aliments et activités disponibles localement. (3) Explique pourquoi on lui déconseille un régime à \num{1000} kcal/j.
\end{exercice}

\begin{corrige}[Exercice : À toi de jouer]
\textbf{(1)} Excédent annuel : $1350 \times 365 \approx \num{492750}$ kJ, soit un gain de $\frac{492750}{30000} \approx 16$ kg en un an.

\textbf{(2)} Programme : supprimer 2 cuillères à soupe d'huile ($-900$ kJ/j) et cuire les légumes à l'eau ; remplacer les beignets du goûter par une goyave ou une banane ($-600$ kJ/j) ; marcher rapidement 40 min/j (marché à pied, $-800$ kJ/j). Bilan : $\Delta E \approx -2300$ kJ/j, soit $\approx$ 0,5 kg/semaine. Suivi : pesée chaque lundi matin, journal alimentaire, réévaluation mensuelle.

\textbf{(3)} Un régime à \num{1000} kcal/j ($< \num{1200}$ kcal/j) provoquerait une fonte musculaire (protéines dégradées en néoglucogenèse), un ralentissement du métabolisme de base (adaptation à la privation) et une reprise pondérale quasi certaine à l'arrêt (effet yo-yo), avec des carences en fer, calcium et vitamines.
\end{corrige}

\newpage

\coursec{Activité d'intégration}

\courssub{Objectifs de l'activité}
\par
Cette activité d'intégration vise à mobiliser l'ensemble des savoirs et savoir-faire de la leçon ND05 dans une tâche complexe, ancrée dans la réalité d'un élève de Première D au Cameroun. À l'issue de cette activité, tu seras capable de :
\begin{itemize}
    \item \textbf{Calculer} les apports énergétiques quotidiens moyens (en kJ) à partir d'un journal alimentaire réel (menus camerounais) en utilisant des tables de composition nutritionnelle.
    \item \textbf{Estimer} les dépenses énergétiques quotidiennes (Métabolisme de Base + Activités Physiques) par la méthode des coefficients d'activité (MET) ou du \cle{NAP} (Niveau d'Activité Physique).
    \item \textbf{Établir} le bilan énergétique journalier ($\Delta E = \text{Apports} - \text{Dépenses}$) et en déduire la variation de masse corporelle prévisionnelle sur 4 semaines.
    \item \textbf{Interpréter} une courbe de variation pondérale simulée et relier le signe du bilan (positif, négatif, nul) aux conséquences physiologiques (stockage/déstockage, adaptation métabolique, risques pathologiques).
    \item \textbf{Proposer} un plan d'action correctif concret, chiffré et réaliste (modifications alimentaires, activité physique adaptée) pour rétablir l'homéostasie énergétique.
    \item \textbf{Communiquer} oralement les résultats et le plan d'action lors d'une restitution structurée de 5 minutes.
\end{itemize}

\courssub{Situation : Le cas d'Armand, élève de 1ère D au Lycée Bilingue de Nkolbisson}
\par
\textbf{Contexte.} Armand, 17 ans, est élève en Première D au Lycée Bilingue de Nkolbisson (Yaoundé). Il mesure 1,72 m pour une masse de 65 kg (IMC $= 65 / 1{,}72^2 \approx 22,0$ kg/m², situation normale). Ses parents, inquiets d'une prise de poids récente (+3 kg en 3 mois) et d'une fatigue vespérale, lui demandent de faire le point sur son hygiène de vie. Armand accepte de tenir un \textbf{journal alimentaire et d'activité sur 3 jours représentatifs} (un jour de sport, un jour de cours classique, un week-end) pour analyser son bilan énergétique.

\par
\textbf{Données anthropométriques et physiologiques :}
\begin{itemize}
    \item Sexe : Masculin ; Âge : 17 ans ; Taille : 1,72 m ; Masse : 65 kg.
    \item Métabolisme de Base (MB) estimé par la formule de \textbf{Harris-Benedict} (homme) :
    $MB (\text{kcal/j}) = 66,5 + 13,75 \times P(\text{kg}) + 5,003 \times T(\text{cm}) - 6,755 \times A(\text{ans})$.
    \item Conversion : 1 kcal = 4,184 kJ.
\end{itemize}

\par
\textbf{Journal des 3 jours (menus camerounais typiques et emploi du temps) :}

\begin{center}
\begin{tabularx}{\linewidth}{|l|X|X|X|}\hline
\rowcolor{popblueL}
\textbf{Moment} & \textbf{Jour 1 : Lundi (jour de sport -- football 1 h 30)} & \textbf{Jour 2 : Mercredi (jour de cours classique)} & \textbf{Jour 3 : Samedi (week-end -- repos/écrans)} \\ \hline
\textbf{6 h 30 -- Réveil / toilette} & Marche 15 min (aller-retour point taxi) & Marche 15 min & Lever tardif, toilette (15 min debout) \\ \hline
\textbf{7 h 00 -- Petit-déjeuner} & \textbf{Bouillie de maïs} (300 g) + \textbf{lait entier} (50 mL) + \textbf{sucre} (20 g) + \textbf{beignets haricots} (2 unités, 100 g) + \textbf{café} (sans sucre) & \textbf{Pain blanc} (150 g) + \textbf{beurre} (20 g) + \textbf{chocolat en poudre} (15 g dans lait) + \textbf{lait entier} (200 mL) + \textbf{banane} (100 g) & \textbf{Bouillie de soja} (350 g) + \textbf{sucre} (25 g) + \textbf{beignets haricots} (3 unités, 150 g) + \textbf{jus d'ananas industriel} (250 mL) \\ \hline
\textbf{7 h 30--12 h 30 -- Cours} & Assis (cours, TP) -- pause 10 h : \textbf{arachides grillées} (50 g) & Assis (cours) -- pause 10 h : \textbf{bonbons} (30 g) + \textbf{soda} (300 mL) & \textbf{Ménage léger} (balayage, vaisselle) 45 min + \textbf{écrans} (téléphone/TV) assis 3 h \\ \hline
\textbf{12 h 30 -- Déjeuner} & \textbf{Ndolé} (feuilles 100 g, viande de bœuf 80 g, crevettes séchées 20 g, huile rouge 15 mL, arachide 10 g) + \textbf{riz blanc cuit} (300 g) + \textbf{mangue} (150 g) + \textbf{eau} (500 mL) & \textbf{Poulet DG} : riz (250 g), plantain frit (150 g), poulet (100 g), légumes (carottes/haricots verts 100 g), huile (10 mL) + \textbf{papaye} (150 g) + \textbf{jus de bissap} (250 mL, sucré) & \textbf{Koki} (haricots niébé 150 g cuits, huile rouge 20 mL, épices) + \textbf{banane plantain cuite} (200 g) + \textbf{avocat} (50 g) + \textbf{eau} \\ \hline
\textbf{13 h 30--16 h 30 -- Cours} & Assis (cours) & Assis (cours) & \textbf{Sieste} (1 h) + \textbf{écrans} (jeux vidéo/TV) 2 h assis \\ \hline
\textbf{16 h 30--18 h 00 -- Activité physique} & \textbf{Football} (entraînement intense : course, dribbles, matchs) \textbf{1 h 30} & \textbf{Marche retour domicile} (30 min, allure modérée) + \textbf{devoirs} assis 1 h & \textbf{Football loisir} entre amis (intensité modérée) \textbf{1 h} \\ \hline
\textbf{18 h 00--20 h 00 -- Soirée} & Douche + \textbf{devoirs/révisions} assis 1 h 30 + \textbf{téléphone} 30 min & Douche + \textbf{devoirs} assis 2 h & Douche + \textbf{sortie amicale} (marche 30 min, discussion debout 1 h) + \textbf{écrans} 30 min \\ \hline
\textbf{20 h 00 -- Dîner} & \textbf{Pâtes} (200 g cuites) + \textbf{tomates/oignons} (100 g) + \textbf{œuf dur} (1, 50 g) + \textbf{huile végétale} (10 mL) + \textbf{jus d'orange pressé} (200 mL) & \textbf{Riz blanc} (250 g) + \textbf{haricots rouges} (150 g cuits) + \textbf{queue de bœuf} (80 g, sauce) + \textbf{huile} (10 mL) + \textbf{ananas} (100 g) & \textbf{Eru} (feuilles 100 g, viande fumée 50 g, peau de vache/kpem 30 g, huile rouge 15 mL, crevettes 10 g) + \textbf{miondo/water fufu} (300 g) + \textbf{boisson maltée} (330 mL) \\ \hline
\textbf{21 h 00--22 h 30 -- Coucher} & Lecture/détente assis 1 h + sommeil & Lecture/détente assis 1 h 30 + sommeil & Discussion/écrans 1 h 30 + sommeil \\ \hline
\end{tabularx}
\end{center}

\par
\textbf{Table de composition nutritionnelle simplifiée (valeurs moyennes pour 100 g ou 100 mL de partie comestible) :}

\begin{center}
\begin{tabular}{|l|c|c|c|c|c|}\hline
\rowcolor{popgreenL}
\textbf{Aliment} & \textbf{Eau (g)} & \textbf{Protéines (g)} & \textbf{Lipides (g)} & \textbf{Glucides (g)} & \textbf{Énergie (kJ)} \\ \hline
Riz blanc cuit & 65 & 2,5 & 0,3 & 28 & 540 \\ \hline
Pâtes cuites & 62 & 5 & 1 & 26 & 580 \\ \hline
Banane plantain cuite & 65 & 1,2 & 0,3 & 30 & 540 \\ \hline
Miondo/water fufu (cuit) & 60 & 1 & 0,2 & 35 & 600 \\ \hline
Bouillie maïs/soja (préparée) & 80 & 2 & 1 & 15 & 350 \\ \hline
Pain blanc & 35 & 8 & 1 & 50 & 1050 \\ \hline
Beurre & 16 & 0,8 & 82 & 0,6 & 3050 \\ \hline
Huile rouge / végétale & 0 & 0 & 100 & 0 & 3700 \\ \hline
Arachide grillée & 3 & 25 & 45 & 15 & 2300 \\ \hline
Viande bœuf / poulet (cuit) & 60 & 25 & 10 & 0 & 850 \\ \hline
Poisson fumé / viande fumée & 40 & 30 & 10 & 0 & 950 \\ \hline
Queue de bœuf / peau (kpem) & 50 & 20 & 25 & 0 & 1200 \\ \hline
Crevettes séchées & 10 & 60 & 5 & 5 & 1350 \\ \hline
Haricots rouges/niébé (cuits) & 65 & 9 & 0,5 & 20 & 520 \\ \hline
Œuf entier (cuit) & 75 & 12,5 & 10 & 0,5 & 650 \\ \hline
Lait entier & 87 & 3,3 & 3,6 & 4,8 & 270 \\ \hline
Sucre / chocolat poudre / confiture & 0 & 0 & 0 & 100 & 1700 \\ \hline
Fruits (mangue, papaye, banane, ananas) & 80--85 & 0,5--1,5 & 0,2 & 10--20 & 180--350 \\ \hline
Avocat & 75 & 2 & 15 & 5 & 700 \\ \hline
Légumes verts (ndolé, eru, carottes) & 85--90 & 2--4 & 0,3 & 3--5 & 100--150 \\ \hline
Beignets haricots (farine + huile) & 30 & 8 & 20 & 40 & 1800 \\ \hline
Soda / jus industriel / bissap sucré / boisson maltée & 90 & 0 & 0 & 10--12 & 170--200 \\ \hline
Bonbons & 5 & 0 & 0 & 90 & 1500 \\ \hline
Café / thé / eau & 100 & 0 & 0 & 0 & 0 \\ \hline
\end{tabular}
\end{center}

\par
\textbf{Coefficients d'activité (MET -- \emph{Metabolic Equivalent of Task}) :}
\par
Le MET exprime l'intensité d'une activité par rapport au repos : 1 MET correspond au métabolisme de base horaire. Pour Armand, on montrera (étape 1 du corrigé) que $MB_{\text{horaire}} \approx 297$ kJ/h. La dépense horaire d'une activité est donc : $D_h = \text{MET} \times 297$ kJ/h. Le tableau ci-dessous donne les valeurs arrondies à utiliser.

\begin{center}
\begin{tabular}{|l|c|c|}\hline
\rowcolor{poporangeL}
\textbf{Activité} & \textbf{MET} & \textbf{Dépense horaire (kJ/h pour 65 kg)} \\ \hline
Sommeil & 0,95 & 280 \\ \hline
Assis (cours, écrans, repas, devoirs) & 1,3 & 385 \\ \hline
Debout léger (toilette, cuisine, discussion) & 1,8 & 535 \\ \hline
Marche lente / ménage léger & 2,5 & 745 \\ \hline
Marche modérée (4--5 km/h) & 3,5 & 1040 \\ \hline
Football loisir / course lente & 7,0 & 2080 \\ \hline
Football entraînement intense & 9,0 & 2675 \\ \hline
\end{tabular}
\end{center}

\courssub{Consignes}
\par
Travailler en groupe de 4 élèves. Répondre aux questions ci-dessous en montrant toutes les étapes de calcul. Présenter les résultats dans un rapport structuré (tableaux, graphiques, texte argumenté). Préparer une restitution orale de 5 minutes.

\begin{enumerate}
    \item \textbf{Apports énergétiques :} calculer l'apport énergétique total (en kJ) pour chacun des 3 jours à l'aide de la table de composition. En déduire l'apport quotidien moyen $\bar{A}$ (en kJ/j). Comparer $\bar{A}$ aux Apports Nutritionnels Conseillés (ANC) pour un garçon de 17 ans, activité modérée (environ 12 500 kJ/j).
    \item \textbf{Dépenses énergétiques :}
    \begin{enumerate}
        \item Calculer le Métabolisme de Base (MB) d'Armand en kJ/j (formule de Harris-Benedict).
        \item Pour chaque jour, comptabiliser les 24 heures et calculer la Dépense Énergétique Totale : $\text{DET} = \sum (\text{Dépense horaire}_i \times \text{Durée}_i)$ en utilisant le tableau MET. Vérifier que la somme des durées vaut 24 h.
        \item Calculer la moyenne $\overline{\text{DET}}$.
    \end{enumerate}
    \item \textbf{Bilan énergétique et variation pondérale :}
    \begin{enumerate}
        \item Calculer le bilan journalier $\Delta E_j = A_j - \text{DET}_j$ pour chaque jour et le bilan moyen $\overline{\Delta E}$.
        \item En admettant que 1 kg de masse corporelle stockée/déstockée correspond à environ 32 000 kJ (équivalent énergétique mixte graisse/eau/protéines), estimer la variation de masse théorique sur 28 jours : $\Delta m_{28} = 28 \times \overline{\Delta E} / 32\,000$ (en kg).
        \item Tracer la courbe prévisionnelle de la masse $m(t)$ sur 4 semaines (abscisse : temps en jours ; ordonnée : masse en kg), en supposant le bilan constant et l'hydratation stable.
    \end{enumerate}
    \item \textbf{Analyse physiologique :}
    \begin{enumerate}
        \item Qualifier le déséquilibre (positif, négatif, nul).
        \item Expliquer les mécanismes physiologiques mis en jeu (lipogenèse vs lipolyse et gluconéogenèse, rôle de l'insuline et du glucagon, adaptation du MB).
        \item Citer 3 risques pathologiques à moyen/long terme si ce déséquilibre persiste.
    \end{enumerate}
    \item \textbf{Plan d'action correctif :} proposer des modifications \textbf{chiffrées} (objectifs en kJ ou en grammes/portions) portant sur :
    \begin{enumerate}
        \item l'alimentation (réduction des lipides ajoutés et des sucres ajoutés, augmentation des fibres, baisse de la densité énergétique) ;
        \item l'activité physique (fréquence, durée, intensité -- viser un \cle{PAL}, \emph{Physical Activity Level}, cible) ;
        \item l'hygiène de vie (sommeil, écrans, hydratation).
    \end{enumerate}
    L'objectif est de ramener $\overline{\Delta E}$ proche de 0 (stabilisation) ou légèrement négatif (environ $-500$ à $-1000$ kJ/j) pour perdre progressivement les 3 kg en excès en 3 mois environ.
\end{enumerate}

\courssub{Ressources à mobiliser}
\begin{propriete}[Formules et références utiles]
\begin{itemize}
    \item \textbf{Métabolisme de Base (Harris-Benedict, homme) :} $MB (\text{kcal/j}) = 66,5 + 13,75\,P + 5,003\,T - 6,755\,A$. Conversion : $MB (\text{kJ/j}) = MB (\text{kcal/j}) \times 4,184$.
    \item \textbf{Dépense horaire d'une activité de MET $M$ :} $D_h = M \times MB_{\text{horaire}}$, avec $MB_{\text{horaire}} = MB(\text{kJ/j}) / 24$. (Le tableau MET fourni donne directement les valeurs arrondies pour 65 kg.)
    \item \textbf{Bilan énergétique :} $\Delta E = \text{Apports totaux} - \text{Dépense Énergétique Totale (DET)}$.
    \item \textbf{Équivalent énergétique de la masse corporelle :} $\approx 32\,000$ kJ/kg (mixte graisse/eau/protéines). Pour le tissu adipeux pur : environ 37 000 kJ/kg.
    \item \textbf{Voies énergétiques cellulaires :}
    \begin{itemize}
        \item Respiration cellulaire aérobie (mitochondries) : $\ce{C6H12O6 + 6 O2 -> 6 CO2 + 6 H2O} + \text{énergie}$ (environ 30 à 38 ATP par glucose selon les estimations ; la valeur scolaire classique est $\sim$38 ATP).
        \item Fermentation lactique (cytoplasme, effort intense et bref, en anaérobiose) : $\ce{C6H12O6 -> 2 C3H6O3} + 2$ ATP.
        \item $\beta$-oxydation des acides gras (effort prolongé, jeûne) : acyl-CoA $\rightarrow$ acétyl-CoA + NADH,H$^+$ + FADH$_2$, puis cycle de Krebs et chaîne respiratoire mitochondriale.
    \end{itemize}
    \item \textbf{Régulation hormonale :} insuline (hormone de stockage, anabolisme) sécrétée en période post-prandiale ; glucagon, adrénaline et cortisol (hormones de mobilisation, catabolisme) sécrétés à jeun et à l'effort.
\end{itemize}
\end{propriete}

\courssub{Démarche et corrigé commenté}
\begin{exoresolu}[Corrigé détaillé de l'activité intégrée -- cas d'Armand]
\textbf{Étape 1 : calcul du Métabolisme de Base (MB)}
\par
$P = 65$ kg, $T = 172$ cm, $A = 17$ ans.
\begin{align*}
MB (\text{kcal/j}) &= 66,5 + 13,75 \times 65 + 5,003 \times 172 - 6,755 \times 17 \\
&= 66,5 + 893,75 + 860,52 - 114,84 \\
&\approx 1705,9 \text{ kcal/j}
\end{align*}
\[
MB (\text{kJ/j}) = 1705,9 \times 4,184 \approx \textbf{7137 kJ/j} \quad (\text{arrondi à } 7140 \text{ kJ/j})
\]
\[
MB_{\text{horaire}} = 7140 / 24 \approx \textbf{297 kJ/h}
\]
C'est cette valeur qui définit 1 MET pour Armand et qui a servi à construire le tableau des dépenses horaires (par exemple : assis, 1,3 MET $\rightarrow 1,3 \times 297 \approx 385$ kJ/h ; football intense, 9 MET $\rightarrow 9 \times 297 \approx 2675$ kJ/h). Le tableau MET est donc \textbf{cohérent avec le MB calculé} : on l'utilisera pour toute la suite.

\textbf{Étape 2 : apports énergétiques quotidiens}
\par
Méthode : $\text{Énergie} = \dfrac{\text{masse (g ou mL)}}{100} \times \text{valeur de la table (kJ/100 g)}$. Pour les huiles, on retient 10 mL $\approx$ 9 g $\approx$ 333 kJ (soit 15 mL $\approx$ 500 kJ, 20 mL $\approx$ 666 kJ).
\par
\textbf{Jour 1 (lundi) :}
\begin{itemize}
    \item Petit-déjeuner : bouillie 300 g ($3 \times 350 = 1050$) + lait 50 mL ($0,5 \times 270 = 135$) + sucre 20 g ($0,2 \times 1700 = 340$) + beignets 100 g ($1 \times 1800 = 1800$) = \textbf{3325 kJ}.
    \item Collation 10 h : arachides 50 g ($0,5 \times 2300 = 1150$ kJ).
    \item Déjeuner : ndolé (feuilles 100 g $\approx 120$ ; viande 80 g $\approx 0,8 \times 850 = 680$ ; crevettes 20 g $\approx 0,2 \times 1350 = 270$ ; huile 15 mL $\approx 500$ ; arachide 10 g $\approx 230$) $\approx 1800$ kJ + riz 300 g ($3 \times 540 = 1620$) + mangue 150 g ($\approx 375$) = \textbf{3795 kJ}.
    \item Dîner : pâtes 200 g ($2 \times 580 = 1160$) + légumes 100 g ($\approx 100$) + œuf 50 g ($0,5 \times 650 = 325$) + huile 10 mL ($\approx 333$) + jus d'orange 200 mL ($\approx 360$) = \textbf{2278 kJ}.
    \item \textbf{Total Jour 1 :} $3325 + 1150 + 3795 + 2278 = \textbf{10 548 kJ} \approx \textbf{10 550 kJ}$.
\end{itemize}
\par
\textbf{Jour 2 (mercredi) :}
\begin{itemize}
    \item Petit-déjeuner : pain 150 g ($1,5 \times 1050 = 1575$) + beurre 20 g ($0,2 \times 3050 = 610$) + chocolat 15 g ($0,15 \times 1700 = 255$) + lait 200 mL ($2 \times 270 = 540$) + banane 100 g ($\approx 250$) = \textbf{3230 kJ}.
    \item Collation : bonbons 30 g ($0,3 \times 1500 = 450$) + soda 300 mL ($3 \times 180 = 540$) = \textbf{990 kJ}.
    \item Déjeuner : riz 250 g ($1350$) + plantain frit 150 g ($\approx 810$) + poulet 100 g ($850$) + légumes 100 g ($120$) + huile 10 mL ($333$) + papaye 150 g ($\approx 270$) + bissap sucré 250 mL ($\approx 450$) = \textbf{4183 kJ}.
    \item Dîner : riz 250 g ($1350$) + haricots 150 g ($1,5 \times 520 = 780$) + queue de bœuf 80 g ($0,8 \times 1200 = 960$) + huile 10 mL ($333$) + ananas 100 g ($\approx 180$) = \textbf{3603 kJ}.
    \item \textbf{Total Jour 2 :} $3230 + 990 + 4183 + 3603 = \textbf{12 006 kJ} \approx \textbf{12 000 kJ}$.
\end{itemize}
\par
\textbf{Jour 3 (samedi) :}
\begin{itemize}
    \item Petit-déjeuner : bouillie 350 g ($3,5 \times 350 = 1225$) + sucre 25 g ($0,25 \times 1700 = 425$) + beignets 150 g ($1,5 \times 1800 = 2700$) + jus industriel 250 mL ($\approx 450$) = \textbf{4800 kJ}.
    \item Déjeuner : koki (haricots 150 g $= 780$ ; huile 20 mL $\approx 666$) + plantain cuit 200 g ($2 \times 540 = 1080$) + avocat 50 g ($0,5 \times 700 = 350$) = \textbf{2876 kJ}.
    \item Dîner : eru (feuilles 100 g $\approx 120$ ; viande fumée 50 g $= 475$ ; kpem 30 g $\approx 0,3 \times 1200 = 360$ ; huile 15 mL $\approx 500$ ; crevettes 10 g $= 135$) + miondo 300 g ($3 \times 600 = 1800$) + boisson maltée 330 mL ($\approx 600$) = \textbf{3990 kJ}.
    \item \textbf{Total Jour 3 :} $4800 + 2876 + 3990 = \textbf{11 666 kJ} \approx \textbf{11 670 kJ}$.
\end{itemize}
\par
\textbf{Moyenne des apports :}
\[
\bar{A} = \frac{10\,550 + 12\,000 + 11\,670}{3} = \frac{34\,220}{3} \approx \textbf{11 400 kJ/j}
\]
\textit{Comparaison aux ANC (garçon de 17 ans, activité modérée : $\approx 12\,500$ kJ/j) :} les apports moyens sont \textbf{légèrement inférieurs aux ANC} (environ $-9\,\%$). Mais attention : les ANC supposent une activité physique \emph{modérée et régulière}, ce qui n'est pas le cas d'Armand (voir étape 3). De plus, la \textbf{qualité} des apports est déséquilibrée : part élevée de lipides ajoutés (huiles, beignets, fritures) et de sucres rapides (sodas, jus industriels, bonbons), apports en fibres limités aux légumes cuits.

\textbf{Étape 3 : Dépenses Énergétiques Totales (DET) par la méthode MET}
\par
On comptabilise les 24 h de chaque journée avec les dépenses horaires du tableau MET.
\par
\textbf{Jour 1 (lundi, sport) :}
\begin{center}
\begin{tabular}{|l|c|c|c|}\hline
\rowcolor{poppinkL}
\textbf{Activité} & \textbf{Durée (h)} & \textbf{Dépense horaire (kJ/h)} & \textbf{Dépense (kJ)} \\ \hline
Sommeil (22 h 30--6 h 30) & 8,0 & 280 & 2240 \\ \hline
Toilette (debout) + marche matin & 0,5 & 535 / 745 & 320 \\ \hline
Cours assis (7 h 30--12 h 30 ; 13 h 30--16 h 30) & 8,0 & 385 & 3080 \\ \hline
Repas et collation (assis) & 1,5 & 385 & 578 \\ \hline
Football intense (16 h 30--18 h) & 1,5 & 2675 & 4013 \\ \hline
Devoirs, téléphone, lecture (assis) & 3,5 & 385 & 1348 \\ \hline
Douche, détente (debout léger) & 1,0 & 535 & 535 \\ \hline
\textbf{Total J1} & \textbf{24,0} & & $\approx$ \textbf{12 114 kJ} \\ \hline
\end{tabular}
\end{center}
On retiendra $\text{DET}_1 \approx \textbf{12 100 kJ}$.
\par
\textbf{Jour 2 (mercredi, sans sport) :} sommeil 8 h ($2240$) ; toilette debout 0,5 h ($268$) ; cours assis 8 h ($3080$) ; repas assis 1,5 h ($578$) ; marche modérée 0,5 h ($520$) ; devoirs assis 3 h ($1155$) ; écrans assis 2 h ($770$) ; douche/détente debout 0,5 h ($268$) ; lecture assise 1 h ($385$). Total : 24 h.
\[
\text{DET}_2 = 2240 + 268 + 3080 + 578 + 520 + 1155 + 770 + 268 + 385 \approx \textbf{9264 kJ} \approx \textbf{9260 kJ}
\]
\par
\textbf{Jour 3 (samedi) :} sommeil 9 h ($2520$) ; toilette debout 0,25 h ($134$) ; ménage léger 0,75 h ($559$) ; écrans assis 5 h ($1925$) ; sieste 1 h ($280$) ; football loisir 1 h ($2080$) ; marche lente 0,5 h ($373$) ; discussion debout 1 h ($535$) ; repas assis 1,5 h ($578$) ; douche debout 0,5 h ($268$) ; discussion/écrans assis 1,5 h ($578$) ; détente assise 2 h ($770$). Total : 24 h.
\[
\text{DET}_3 = 2520 + 134 + 559 + 1925 + 280 + 2080 + 373 + 535 + 578 + 268 + 578 + 770 \approx \textbf{10 600 kJ}
\]
\par
\textbf{Moyenne des dépenses :}
\[
\overline{\text{DET}} = \frac{12\,100 + 9\,260 + 10\,600}{3} = \frac{31\,960}{3} \approx \textbf{10 650 kJ/j}
\]
\textit{Remarque :} le rapport $\overline{\text{DET}} / MB \approx 10\,650 / 7\,140 \approx 1,5$ correspond à un \cle{PAL} (NAP) de 1,5, c'est-à-dire une activité « légère à modérée » -- et non « modérée » comme le supposent les ANC.

\textbf{Étape 4 : bilan énergétique et variation pondérale}
\begin{center}
\begin{tabular}{|c|c|c|c|}\hline
\rowcolor{popgoldL}
\textbf{Jour} & \textbf{Apports A (kJ)} & \textbf{Dépenses DET (kJ)} & \textbf{Bilan $\Delta E = A - \text{DET}$ (kJ)} \\ \hline
Lundi (sport) & 10 550 & 12 100 & \textbf{$-1\,550$} \\ \hline
Mercredi (classique) & 12 000 & 9 260 & \textbf{$+2\,740$} \\ \hline
Samedi (week-end) & 11 670 & 10 600 & \textbf{$+1\,070$} \\ \hline
\textbf{Moyenne} & \textbf{11 400} & \textbf{10 650} & \textbf{$+750$} \\ \hline
\end{tabular}
\end{center}
\par
\textbf{Analyse :} le bilan moyen est \textbf{positif} ($\overline{\Delta E} \approx +750$ kJ/j). Le jour de sport, Armand est même en déficit ($-1550$ kJ) : ce n'est pas le sport qui pose problème, mais les \textbf{jours sans sport}, où les apports restent élevés alors que la dépense chute fortement (sédentarité : écrans, sieste, cours assis).
\par
\textbf{Variation de masse sur 28 jours :}
\[
\Delta m_{28} = \frac{28 \times 750}{32\,000} = \frac{21\,000}{32\,000} \approx \textbf{+0,66 kg}
\]
Sur 3 mois ($\approx 13$ semaines), on obtient $\Delta m \approx 750 \times 91 / 32\,000 \approx +2,1$ kg, ordre de grandeur \textbf{compatible avec les +3 kg en 3 mois} constatés par les parents (l'écart s'explique par les approximations des tables et par les jours de fête non comptés).
\par
\textbf{Courbe prévisionnelle $m(t)$ :} droite affine croissante $m(t) = 65 + \dfrac{0,66}{28}\,t \approx 65 + 0,024\,t$ ($t$ en jours).
\begin{popfigure}
\begin{tikzpicture}
\begin{axis}[
    width=\linewidth, height=8cm,
    axis lines=middle,
    xlabel={Temps (jours)}, ylabel={Masse (kg)},
    xmin=0, xmax=28, ymin=64.5, ymax=66.5,
    xtick={0,7,14,21,28}, ytick={65,65.5,66},
    grid=major, grid style={popline}
]
\addplot[popcurve, thick, domain=0:28, samples=2] {65 + 0.0236*x};
\node[above right, popdark] at (axis cs:10,65.45) {$\Delta m \approx +0,66$ kg en 28 j};
\end{axis}
\end{tikzpicture}
\legende{Évolution prévisionnelle de la masse d'Armand sur 4 semaines si le bilan énergétique moyen ($+750$ kJ/j) se maintient. La pente positive traduit un stockage net (lipogenèse hépatique et adipocytaire). Prolongée sur 3 mois, cette tendance redonne l'ordre de grandeur de la prise de poids observée (+2 à +3 kg).}
\end{popfigure}

\textbf{Étape 5 : analyse physiologique du déséquilibre positif}
\begin{itemize}
    \item \textbf{Type :} déséquilibre \textbf{positif chronique modéré} (apports $>$ dépenses en moyenne), masqué par le déficit du jour de sport.
    \item \textbf{Mécanismes cellulaires du stockage :} en période d'excédent, le glucose excédentaire est dégradé par la glycolyse en pyruvate puis en acétyl-CoA ; lorsque les besoins énergétiques sont couverts, le citrate s'accumule dans la mitochondrie et est exporté vers le cytosol, où il régénère l'acétyl-CoA, point de départ de la \textbf{lipogenèse de novo} : acétyl-CoA $\xrightarrow{\text{acétyl-CoA carboxylase (ACC)}}$ malonyl-CoA $\xrightarrow{\text{acide gras synthase (FAS)}}$ acides gras $\rightarrow$ triglycérides. Le malonyl-CoA inhibe en outre la carnitine palmitoyltransférase I (CPT1), ce qui bloque l'entrée des acides gras dans la mitochondrie et donc leur $\beta$-oxydation : la cellule est en mode « stockage ». Les triglycérides sont stockés dans les adipocytes (hypertrophie, puis hyperplasie si l'excès persiste).
    \item \textbf{Régulation hormonale :} les repas riches en glucides rapides et en lipides provoquent des pics d'\textbf{insuline} répétés ; l'insuline stimule la lipoprotéine lipase (LPL) des adipocytes (captation des acides gras circulants) et inhibe la lipase hormono-sensible (HSL), qui catalyse la lipolyse. À long terme, la \textbf{leptine} (hormone de satiété sécrétée par les adipocytes) augmente, mais une résistance à la leptine peut s'installer.
    \item \textbf{Risques pathologiques à moyen/long terme (3 exemples) :}
    \begin{enumerate}
        \item \textbf{Surcharge pondérale puis obésité} (en particulier abdominale), favorisant le syndrome métabolique (hypertension artérielle, dyslipidémie, hyperglycémie).
        \item \textbf{Stéatose hépatique non alcoolique (NAFLD)} : accumulation de triglycérides dans les hépatocytes, favorisée par la lipogenèse de novo et les boissons sucrées.
        \item \textbf{Insulinorésistance puis diabète de type 2} : l'accumulation de lipides dans le muscle et le foie perturbe la voie de signalisation de l'insuline ; le pancréas finit par ne plus compenser, d'où une hyperglycémie chronique.
    \end{enumerate}
\end{itemize}

\textbf{Étape 6 : plan d'action correctif}
\par
\textit{Objectif : passer de $\overline{\Delta E} \approx +750$ kJ/j à environ $-500$ à $-1000$ kJ/j, soit une perte de 0,5 à 1 kg par mois, donc environ 3 kg en 3 mois. Stratégie combinée : réduire les apports d'environ 1000 kJ/j (qualité) et augmenter les dépenses d'environ 500 kJ/j (quantité).}

\begin{center}
\begin{tabularx}{\linewidth}{|l|X|X|}\hline
\rowcolor{poppurpleL}
\textbf{Levier} & \textbf{Action concrète (chiffrée)} & \textbf{Impact estimé (kJ/j)} \\ \hline
\textbf{Alimentation} & \textbf{Supprimer les boissons sucrées} (soda, bissap sucré, jus industriels, boisson maltée) : environ 2 à 3 verres/j $\approx$ 800 à 1000 kJ. Remplacer par de l'eau ou des infusions non sucrées. & $-900$ \\ \hline
& \textbf{Réduire les huiles ajoutées} (cuisson, sauces) : passer d'environ 45 mL/j à 20 mL/j (cuisson vapeur, braisage sans bain d'huile, mesure de l'huile à la cuillère). Gain $\approx 25$ mL $\times$ 33 kJ/mL. & $-800$ \\ \hline
& \textbf{Modifier le petit-déjeuner} : remplacer les beignets (1800 kJ/100 g) par des fruits frais + pain complet ou bouillie peu sucrée + une source de protéines (œuf, lait). & $-800$ \\ \hline
& \textbf{Augmenter légumes et fibres} (crudités en entrée, légumes feuilles à chaque repas) : satiété accrue, densité énergétique du repas diminuée. & effet indirect \\ \hline
\textbf{Activité physique} & \textbf{Ajouter 2 séances/semaine} de course ou renforcement musculaire de 45 min (MET 7--9, $\approx$ 1500 à 2000 kJ/séance) : moyenne $\approx +500$ kJ/j. & $+500$ \\ \hline
& \textbf{Marche active quotidienne} de 30 min (trajet école ou marche du soir, MET 3,5) : $0,5 \times 1040$. & $+520$ \\ \hline
& \textbf{Réduire la sédentarité} : se lever 5 min toutes les 45 min pendant devoirs et écrans (thermogenèse des activités non sportives, NEAT). & $+200$ \\ \hline
\textbf{Hygiène de vie} & \textbf{Sommeil} : coucher à 22 h pour 8 h 30--9 h de sommeil ; le manque de sommeil augmente la ghréline (faim) et diminue la leptine (satiété). & prévention de la prise de poids \\ \hline
& \textbf{Écrans} : arrêt 30 min avant le coucher (la lumière bleue retarde la sécrétion de mélatonine). & qualité du sommeil \\ \hline
& \textbf{Hydratation} : 1,5 à 2 L d'eau par jour ; la soif est parfois confondue avec la faim. & contrôle de l'appétit \\ \hline
\end{tabularx}
\end{center}
\par
\textbf{Bilan prévisionnel après intervention :} $\bar{A}_{\text{new}} \approx 11\,400 - 1\,500 = 9\,900$ kJ/j (en ne cumulant pas toutes les mesures alimentaires pour rester réaliste) ; $\overline{\text{DET}}_{\text{new}} \approx 10\,650 + 700 = 11\,350$ kJ/j. D'où $\Delta E_{\text{new}} \approx -1\,450$ kJ/j au départ, à ajuster vers $-1000$ kJ/j : perte théorique $\approx 1000 \times 90 / 32\,000 \approx 2,8$ kg en 3 mois. \textbf{Objectif atteint.} Une fois le poids cible atteint, on stabilise en réaugmentant légèrement les apports (féculents complets, fruits) pour retrouver $\Delta E \approx 0$.
\par
\textbf{Restitution orale (5 min) :}
1. Présentation du profil et de la méthode (30 s).
2. Résultats clés : apports, dépenses, bilan (tableau + courbe) (1 min 30 s).
3. Interprétation physiologique (lipogenèse, insuline, risques) (1 min).
4. Plan d'action SMART (Spécifique, Mesurable, Atteignable, Réaliste, Temporellement défini) (1 min 30 s).
5. Conclusion motivante (30 s).
\end{exoresolu}

\courssub{Bilan de l'activité}
\par
Cette activité d'intégration a permis de mettre en évidence plusieurs points cruciaux du programme de bioénergétique en Première D :
\begin{enumerate}
    \item \textbf{La quantification rigoureuse :} passer du qualitatif (« manger trop », « bouger peu ») au quantitatif (kJ, MET, $\Delta E$) est indispensable pour objectiver un déséquilibre. L'utilisation de tables de composition adaptées au contexte alimentaire camerounais (huile rouge, beignets, plats en sauce, féculents locaux) ancre la biologie dans le réel de l'élève.
    \item \textbf{L'interdépendance apports/dépenses :} le cas d'Armand illustre qu'un apport moyen apparemment « normal » (11 400 kJ, proche des ANC) peut masquer un déséquilibre si la dépense est faible les jours sans sport. Le \cle{NAP} (ou PAL) est le principal levier de variabilité de la DET : il fait passer celle-ci d'environ $1,3 \times MB$ (sédentaire) à $2 \times MB$ ou plus (très actif).
    \item \textbf{La dynamique pondérale :} la modélisation linéaire ($\Delta m = \Delta E \times t / 32\,000$) est une approximation pédagogique puissante : un excédent « invisible » de 750 kJ/j (l'équivalent d'un soda et d'une poignée d'arachides) suffit à expliquer une prise de 2 à 3 kg en 3 mois.
    \item \textbf{Les mécanismes moléculaires du stockage :} l'excédent glucidique et lipidique converge vers l'acétyl-CoA cytosolique pour la lipogenèse (enzymes ACC et FAS), sous contrôle de l'insuline. Ce ne sont pas seulement « les calories » qui font grossir, mais le \textbf{signal hormonal anabolique} (insuline) maintenu élevé de façon chronique par des apports fréquents à index glycémique élevé et à forte densité énergétique.
    \item \textbf{L'approche systémique de la correction :} le plan d'action ne se résume pas à « manger moins ». Il combine :
    \begin{itemize}
        \item la \textbf{densité nutritionnelle} (réduire lipides cachés et sucres ajoutés, augmenter fibres et protéines, donc la satiété) ;
        \item la \textbf{dépense active} (sport programmé + thermogenèse des activités non sportives, NEAT : marche, pauses actives) ;
        \item la \textbf{chronobiologie} (sommeil suffisant, limitation des écrans le soir) pour préserver la sensibilité à l'insuline et l'équilibre leptine/ghréline.
    \end{itemize}
    \item \textbf{Compétences transversales :} calculs avec unités (kJ, kcal, MET), lecture de tableaux, tracé et interprétation de courbes, argumentation scientifique écrite et orale, travail collaboratif -- autant de compétences au cœur de l'approche par les compétences du MINESEC et attendues au Probatoire.
\end{enumerate}
\par
\emph{En tant que futur scientifique ou citoyen éclairé, tu comprends maintenant que l'homéostasie énergétique n'est pas un compte bancaire statique, mais un système dynamique finement régulé. Agir sur l'un des paramètres (apport, dépense, sommeil, stress) fait basculer l'équilibre. Ta mission, si tu l'acceptes, est d'être l'ingénieur de ta propre machine métabolique !}

\newpage

\coursec{Méthodes et exercices résolus}

\coursec{Sensibilisation sur les déséquilibres énergétiques au sein des organismes et leurs conséquences}

\courssub{Sources et voies de production de l'énergie cellulaire}

\begin{definition}[Bioénergétique cellulaire et ATP]
La \textbf{bioénergétique} étudie les transformations d'énergie au sein de la cellule. La molécule universelle de transfert d'énergie utilisable est l'\textbf{ATP} (adénosine triphosphate). Son hydrolyse en ADP + P$_i$ libère de l'énergie ($\approx \SI{-30,5}{kJ/mol}$ dans les conditions cellulaires) qui alimente les réactions endothermiques (biosynthèse, transport actif, contraction musculaire, division cellulaire). La régénération de l'ATP à partir d'ADP + P$_i$ nécessite un apport d'énergie fourni par l'oxydation des nutriments (glucides, lipides, protéines).
\end{definition}

\begin{propriete}[Les deux grandes voies de production d'ATP]
\begin{itemize}
\item \textbf{Respiration cellulaire (aérobie)} : oxydation complète du glucose (ou acides gras, acides aminés) en \ce{CO2} et \ce{H2O} en présence d'\ce{O2}. Elle se déroule dans le cytoplasme (glycolyse) et les mitochondries (cycle de Krebs, chaîne respiratoire). Rendement élevé : \textbf{36 à 38 ATP par glucose}.
\item \textbf{Fermentation (anaérobie)} : dégradation incomplète du glucose sans \ce{O2}. Seule la glycolyse a lieu (cytoplasme), le pyruvate est réduit pour régénérer le \ce{NAD+}. Rendement faible : \textbf{2 ATP par glucose}.
\begin{itemize}
\item Fermentation lactique (muscle, globules rouges, certaines bactéries) : \ce{C6H12O6 -> 2 CH3-CHOH-COOH + 2 ATP} (acide lactique).
\item Fermentation alcoolique (levures, plantes) : \ce{C6H12O6 -> 2 C2H5OH + 2 CO2 + 2 ATP}.
\end{itemize}
\end{itemize}
\end{propriete}

\begin{popfigure}
\begin{tikzpicture}[
    node distance=1.2cm and 1.5cm,
    box/.style={draw=popdark, thick, rounded corners, fill=popcream, minimum width=2.5cm, minimum height=1cm, align=center},
    arrow/.style={-Stealth, thick, popdark},
    smallbox/.style={draw=popblue, fill=popblueL, rounded corners, minimum width=2cm, minimum height=0.7cm, align=center, font=\small}
]
% Nodes
\node[box] (glucose) {Glucose \ce{C6H12O6}};
\node[box, below left=of glucose, align=center] (glycolyse){Glycolyse (Cytoplasme)\\\ce{2 Pyruvate + 2 ATP + 2 NADH}};
\node[smallbox, below right=of glycolyse, align=center] (ferment){Fermentation\\Anaérobie};
\node[smallbox, below left=of glycolyse, align=center] (resp){Respiration\\Aérobie};
\node[box, below=of ferment, align=center] (lact){Fermentation lactique\\\ce{2 Lactate + 2 ATP}};
\node[box, below=of resp, align=center] (mito){Mitochondrie\\\ce{2 Acetyl-CoA -> Krebs -> Chaine respiratoire}};
\node[box, below=of mito] (atp) {\textbf{~36-38 ATP} + \ce{6 CO2 + 6 H2O}};

% Arrows
\draw[arrow] (glucose) -- (glycolyse);
\draw[arrow] (glycolyse) -- node[above, sloped, font=\small] {Si \ce{O2} présent} (resp);
\draw[arrow] (glycolyse) -- node[above, sloped, font=\small] {Si \ce{O2} absent} (ferment);
\draw[arrow] (ferment) -- (lact);
\draw[arrow] (resp) -- (mito);
\draw[arrow] (mito) -- (atp);
\end{tikzpicture}
\legende{Schémas des voies métaboliques du glucose : bifurcation vers la fermentation (anaérobie, 2 ATP) ou la respiration (aérobie, ~36-38 ATP) selon la disponibilité en oxygène.}
\end{popfigure}

\begin{methode}[Exploiter les bilans moléculaires des voies énergétiques]
\begin{enumerate}
\item \textbf{Écrire l'équation bilan correcte} selon la voie et le substrat (glucide, lipide) :
\begin{itemize}
\item Respiration aérobie (glucide) : \ce{C6H12O6 + 6 O2 -> 6 CO2 + 6 H2O} \quad ($\Delta G^\circ \approx \SI{-2870}{kJ/mol}$, $\approx$ 36-38 ATP).
\item Fermentation lactique : \ce{C6H12O6 -> 2 C3H6O3 + 2 ATP} (acide lactique).
\item Fermentation alcoolique : \ce{C6H12O6 -> 2 C2H5OH + 2 CO2 + 2 ATP}.
\item $\beta$-oxydation d'un acide gras (ex. palmitate C16) : \ce{C16H32O2 + 23 O2 -> 16 CO2 + 16 H2O} ($\approx$ 106 ATP).
\end{itemize}
\item \textbf{Calculer le Quotient Respiratoire (QR)} : $QR = \dfrac{V(\ce{CO2})}{V(\ce{O2})}$.
\begin{itemize}
\item QR $= 1$ : oxydation quasi exclusive de glucides.
\item QR $\approx 0,7$ : oxydation quasi exclusive de lipides.
\item QR $\approx 0,8-0,9$ : mélange ou oxydation de protéines.
\end{itemize}
\item \textbf{Relier au contexte physiologique} : effort intense $\rightarrow$ apport \ce{O2} insuffisant $\rightarrow$ fermentation lactique $\rightarrow$ accumulation lactate + \ce{H+} $\rightarrow$ baisse pH $\rightarrow$ fatigue.
\end{enumerate}
\textbf{Réflexe Probatoire} : Ne jamais écrire « la fermentation produit de l'énergie sans glucose » (elle consomme du glucose). Ne pas confondre \textbf{pyruvate} (produit glycolyse) et \textbf{lactate} (produit fermentation lactique). Préciser le rendement ATP à chaque fois.
\end{methode}

\begin{exoresolu}[Respiration ou fermentation lors d'un sprint]
Lors d'un match de football, un joueur effectue un sprint de 80 m. La concentration en lactate sanguin passe de 1 mmol/L (repos) à 12 mmol/L (post-effort).
\begin{enumerate}
\item Écris l'équation de la voie métabolique responsable.
\item Explique le recours à cette voie pendant le sprint.
\item Compare le rendement ATP de cette voie à la respiration.
\end{enumerate}
\tcblower
\textbf{1. Équation de la fermentation lactique musculaire :}
\[ \ce{C6H12O6 -> 2 CH3-CHOH-COOH + 2 ATP} \quad (\text{Glucose} \rightarrow \text{2 Lactate} + \text{2 ATP}) \]

\textbf{2. Pourquoi cette voie ?}
Le sprint est un effort \textbf{bref et maximal}. La demande en ATP explose. L'apport d'\ce{O2} (ventilation, débit cardiaque, diffusion) ne peut suivre instantanément. La respiration aérobie (nécessitant \ce{O2} comme accepteur final d'électrons) est saturée. Le muscle poursuit la glycolyse : le pyruvate accepte les électrons du NADH (via lactate déshydrogénase) pour régénérer le \ce{NAD+}, indispensable à la glycolyse. Le lactate (acide lactique) s'accumule dans la fibre puis diffuse dans le sang (1 $\rightarrow$ 12 mmol/L).

\textbf{3. Comparaison des rendements :}
\begin{itemize}
\item Fermentation lactique : \textbf{2 ATP / glucose}.
\item Respiration cellulaire : \textbf{36 à 38 ATP / glucose}.
\end{itemize}
La fermentation est \textbf{18 à 19 fois moins rentable}. Elle permet une puissance ATPique maximale (vitesse) mais épuise vite le glycogène. L'acidose (lactate + \ce{H+}) inhibe les enzymes glycolytiques et la contraction $\rightarrow$ fatigue. La récupération nécessite une \textbf{dette d'oxygène} (EPOC) pour oxyder le lactate.

\textbf{Conclusion} : La fermentation lactique est une adaptation d'urgence à l'hypoxie tissulaire relative, efficace en puissance mais coûteuse en substrat.
\end{exoresolu}

\courssub{Bilan énergétique de l'organisme : apports et dépenses}

\begin{definition}[Composantes du bilan énergétique]
Le \textbf{bilan énergétique} (BE) d'un organisme sur une période donnée s'exprime en kJ (ou kcal) :
\[ \text{BE} = \text{Apports énergétiques (AE)} - \text{Dépenses énergétiques (DE)} \]
\begin{itemize}
\item \textbf{Apports (AE)} : Énergie fournie par l'oxydation des macronutriments absorbés.
\begin{center}
\begin{tabular}{|c|c|}\hline \rowcolor{popblueL} Nutriment & Valeur énergétique de référence \\ \hline Glucides & \SI{17}{kJ/g} (\approx \SI{4}{kcal/g}) \\ \hline Lipides & \SI{38}{kJ/g} (\approx \SI{9}{kcal/g}) \\ \hline Protéines & \SI{17}{kJ/g} (\approx \SI{4}{kcal/g}) \\ \hline Alcool (éthanol) & \SI{29}{kJ/g} (\approx \SI{7}{kcal/g}) \\ \hline
\end{tabular}
\end{center}
\item \textbf{Dépenses (DE)} = Métabolisme de Base (MB) + Thermogenèse post-prandiale (10\% AE) + Dépense d'Activité Physique (DAP) + Thermorégulation (souvent incluse dans MB/DAP).
\begin{itemize}
\item \textbf{Métabolisme de Base (MB)} : Énergie minimale pour les fonctions vitales au repos complet, à jeun, thermoneutralité. Représente 60-70\% DE. Dépend de la masse maigre, âge, sexe, thyroïde.
\item \textbf{DAP} : Variable (15-30\% DE), seule composante volontairement modulable.
\end{itemize}
\end{itemize}
\end{definition}

\begin{methode}[Établir et interpréter un bilan énergétique quantitatif]
\begin{enumerate}
\item \textbf{Lister les apports} : Convertir les masses (g) de chaque nutriment en kJ à l'aide des valeurs données dans l'énoncé. $E = m \times V_e$.
\item \textbf{Sommer les dépenses} : MB + DAP (+ thermorégulation si distincte). Vérifier que toutes les composantes sont incluses une seule fois.
\item \textbf{Calculer le bilan} : $\text{BE} = \sum \text{AE} - \sum \text{DE}$.
\item \textbf{Interpréter le signe} :
\begin{itemize}
\item BE $> 0$ : \textbf{Déséquilibre positif} $\rightarrow$ Stockage (triglycérides dans adipocytes, glycogène hépatique/musculaire) $\rightarrow$ Prise de masse.
\item BE $< 0$ : \textbf{Déséquilibre négatif} $\rightarrow$ Déstockage (glycogène, puis lipides, puis protéines) $\rightarrow$ Perte de masse.
\item BE $\approx 0$ : \textbf{Équilibre} $\rightarrow$ Masse stable (homéostasie énergétique).
\end{itemize}
\item \textbf{Proposer une correction chiffrée} : Quantifier l'ajustement (réduire apports de X kJ ou augmenter DE de Y kJ) pour ramener BE $\rightarrow$ 0.
\end{enumerate}
\textbf{Réflexes} : Unités homogènes (tout en kJ ou tout en kcal ; $1\ \text{kcal} = \num{4,18}\ \text{kJ}$). Présenter les calculs en tableau. Conclure par une phrase complète reliant le chiffre à la physiologie.
\end{methode}

\begin{exoresolu}[Bilan énergétique d'un élève de Première D]
Un élève de 17 ans consomme : 400 g glucides, 90 g lipides, 80 g protéines. DE = \SI{10500}{kJ/jour}. $V_e$ : Glucides \SI{17}{kJ/g}, Lipides \SI{38}{kJ/g}, Protéines \SI{17}{kJ/g}.
\begin{enumerate}
\item Calcule l'apport énergétique journalier.
\item Établis le bilan et interprète-le.
\item Propose une mesure corrective chiffrée.
\end{enumerate}
\tcblower
\textbf{1. Apports :}
\begin{align*}
E_G &= 400 \times 17 = \SI{6800}{kJ} \\
E_L &= 90 \times 38 = \SI{3420}{kJ} \\
E_P &= 80 \times 17 = \SI{1360}{kJ} \\
\text{AE}_{\text{total}} &= 6800 + 3420 + 1360 = \SI{11580}{kJ/jour}
\end{align*}

\textbf{2. Bilan :}
\[ \text{BE} = 11580 - 10500 = +\SI{1080}{kJ/jour} \]
Bilan \textbf{positif}. Excédent stocké (principalement triglycérides). Risque de surpoids/obésité à terme (diabète type 2, HTA, CVD).

\textbf{3. Correction :}
Annuler $+\SI{1080}{kJ/jour}$.
Option A : Réduire lipides de $\frac{1080}{38} \approx \SI{28}{g/jour}$.
Option B : Augmenter DAP (ex. 35 min course $\approx \SI{1100}{kJ}$).
Option C (recommandée) : Combinaison (ex. -15 g lipides + 20 min marche rapide).
\textbf{Conclusion} : Bilan $+\SI{1080}{kJ/j}$ $\rightarrow$ stockage. Rééquilibrage nécessaire par hygiène de vie.
\end{exoresolu}

\begin{methode}[Interpréter une courbe de variation pondérale]
\begin{enumerate}
\item \textbf{Décrire} : Tendance (croissante, décroissante, palier), valeurs initiales/finales, durée.
\item \textbf{Relier au bilan} : Masse stable $\leftrightarrow$ BE $\approx$ 0 ; Hausse $\leftrightarrow$ BE $>$ 0 ; Baisse $\leftrightarrow$ BE $<$ 0.
\item \textbf{Quantifier} : 1 kg tissu adipeux $\approx \SI{30000}{kJ}$ (adipocytes $\approx$ 87\% lipides). Déficit/Excédent journalier moyen $= \frac{\Delta m \times 30000}{\text{nb jours}}$.
\item \textbf{Diagnostiquer} : Perte $> 10\%$ masse initiale = \textbf{dénutrition} (kwashiorkor/marasme chez l'enfant). IMC $\geq 30$ = \textbf{obésité} ($IMC = \frac{m}{t^2}$).
\item \textbf{Attention} : Variation rapide (jours) = eau/glycogène. Variation lente (semaines) = graisses/protéines.
\end{enumerate}
\end{methode}

\begin{exoresolu}[Dénutrition infectieuse à l'hôpital de Yaoundé]
Patient infection digestive chronique. Masse : S0=65 kg, S2=62, S4=59, S6=56, S8=54.
\begin{enumerate}
\item Perte totale et \%.
\item Interprétation bilan.
\item Déficit journalier moyen (1 kg $\approx \SI{30000}{kJ}$).
\end{enumerate}
\tcblower
\textbf{1.} $\Delta m = 11\ \text{kg}$. $\% = \frac{11}{65}\times 100 \approx 16,9\%$.
\textbf{2.} Baisse continue $\rightarrow$ \textbf{BE négatif chronique}. Malabsorption (apports utilisables $\downarrow$) + besoins $\uparrow$ (fièvre, réponse immune). Mobilisation réserves : glycogène $\rightarrow$ lipides $\rightarrow$ protéines (fonte musculaire, $\uparrow$ urémie). Perte $> 10\%$ = \textbf{dénutrition sévère} $\rightarrow$ immunodépression, retard cicatrisation, cercle vicieux.
\textbf{3.} Durée = 56 jours. Déficit total $= 11 \times 30000 = \SI{330000}{kJ}$. Déficit moyen $= \frac{330000}{56} \approx \SI{5893}{kJ/jour} \approx \SI{5,9}{MJ/jour}$.
\textbf{Conclusion} : Dénutrition (17\% perte) par déficit $\approx \SI{5,9}{MJ/j}$. Prise en charge : nutrition entérale/parentérale enrichie + traitement étiologique.
\end{exoresolu}

\courssub{Conséquences physiologiques des déséquilibres énergétiques}

\begin{propriete}[Conséquences pathologiques des déséquilibres chroniques]
\begin{center}
\begin{tabularx}{\linewidth}{|>{\bfseries}l|X|X|}\hline
\rowcolor{popblueL} & \textbf{Déséquilibre POSITIF chronique (BE $> 0$)} & \textbf{Déséquilibre NÉGATIF chronique (BE $< 0$)} \\ \hline
\textbf{Mécanisme central} & Stockage excédent $\rightarrow$ Hypertrophie/Hyperplasie adipocytes (tissu adipeux blanc) & Mobilisation réserves : Glycogène $\rightarrow$ Lipides (lipolyse) $\rightarrow$ Protéines (protéolyse musculaire) \\ \hline
\textbf{Conséquences métaboliques} & Insulinorésistance, dyslipidémie (HTG, LDL$\uparrow$, HDL$\downarrow$), hyperuricémie & Hypoglycémie, cétonémie (corps cétoniques), hyperurémie (catabolisme protéique), perte électrolytes \\ \hline
\textbf{Pathologies associées} & \textbf{Obésité} (IMC $\geq 30$), Diabète type 2, HTA, Athérosclérose/MCV, Stéatose hépatique (NASH), Apnée du sommeil, Cancers (sein, côlon, endomètre) & \textbf{Dénutrition} (Marasme/Kwashiorkor), Anémie, Immunodépression (infections), Ostéoporose, Aménorrhée, Retard croissance (enfant), Insuffisance cardiaque/réspiratoire \\ \hline
\end{tabularx}
\end{center}
\end{propriete}

\begin{savaistu}[Transition nutritionnelle au Cameroun]
Le Cameroun fait face à un \textbf{double fardeau nutritionnel} : persistance de la dénutrition (régions Extrême-Nord, Nord, Adamaoua : retard de croissance $\approx 30\%$ enfants $<5$ ans) et explosion du surpoids/obésité en milieu urbain (Yaoundé, Douala : $\approx 30\%$ adultes en surpoids, $10-15\%$ obèses). Causes : sédentarisation, alimentation ultra-transformée (huiles de palme raffinées, boissons sucrées, pains blancs), réduction activité physique. Le Programme National de Lutte contre la Malnutrition (PNLM) et les campagnes « Manger Bouger » visent à promouvoir l'alimentation locale diversifiée (feuilles vertes, légumineuses, poisson, fruits) et l'activité physique.
\end{savaistu}

\begin{methode}[Construire l'argumentation : Déséquilibre $\rightarrow$ Conséquences $\rightarrow$ Correction]
\begin{enumerate}
\item \textbf{Qualifier} le déséquilibre (positif/négatif, aigu/chronique) à partir des données (BE calculé, courbe masse, IMC, biologie).
\item \textbf{Enchaîner logiquement} : Excédent $\rightarrow$ Stockage TG $\rightarrow$ Adiposité $\rightarrow$ Inflammation bas bruit/Insulinorésistance $\rightarrow$ Pathologies. Déficit $\rightarrow$ Catabolisme $\rightarrow$ Amaigrissement $\rightarrow$ Dénutrition $\rightarrow$ Complications.
\item \textbf{Proposer des mesures correctives \textbf{chiffrées, durables, globales}} :
\begin{itemize}
\item \textit{Apports} : Réduire densité énergétique (limiter lipides ajoutés, sucres rapides), augmenter fibres/légumes, fractionner repas. \textbf{Jamais} exclure une catégorie.
\item \textit{Dépenses} : Activité physique régulière (OMS : 150-300 min/semaine intensité modérée). Protège masse musculaire et sensibilité à l'insuline.
\item \textit{Éducation} : Repères alimentaires (guide alimentaire camerounais), lecture étiquettes, gestion stress/sommeil.
\end{itemize}
\item \textbf{Conclure} : Lier cause (quantifiée), mécanisme (biologique) et solution (réaliste).
\end{enumerate}
\end{methode}

\begin{exoresolu}[Obésité d'un commerçant sédentaire]
Homme 45 ans, 1,70 m, 98 kg. AE = \SI{14000}{kJ/j}, DE = \SI{9800}{kJ/j}.
\begin{enumerate}
\item IMC et qualification.
\item Bilan énergétique.
\item Mécanisme et 2 conséquences.
\item 2 mesures correctives cohérentes.
\end{enumerate}
\tcblower
\textbf{1.} $IMC = \frac{98}{1,70^2} = \frac{98}{2,89} \approx 33,9\ \text{kg/m}^2$. \textbf{Obésité} (classe I).
\textbf{2.} BE $= 14000 - 9800 = +\SI{4200}{kJ/j}$. Excédent majeur.
\textbf{3.} Excédent glucides/lipides $\rightarrow$ Lipogenèse hépatique $\rightarrow$ VLDL $\rightarrow$ Stockage TG adipocytes (hypertrophie). Adipocytes hypertrophiés secrètent adipokines pro-inflammatoires (TNF-$\alpha$, IL-6, leptine$\uparrow$, adiponectine$\downarrow$) $\rightarrow$ \textbf{Insulinorésistance} $\rightarrow$ Diabète type 2. Excès lipides sanguins $\rightarrow$ \textbf{Athérosclérose} $\rightarrow$ Risque IDM/AVC.
\textbf{4.} Objectif : Annuler $+\SI{4200}{kJ/j}$.
\begin{itemize}
\item \textbf{Alimentation} : Réduire \SI{2280}{kJ} lipides (-60 g, ex. huile, beurre, fritures) + \SI{1920}{kJ} glucides (-110 g, ex. sodas, pain blanc, sucreries) = \SI{4200}{kJ}. Privilégier légumes, fruits, protéines maigres (poisson, volaille, légumineuses).
\item \textbf{Activité} : Marche rapide 45 min/j ($\approx \SI{1000}{kJ}$) + renforcement musculaire 2$\times$/sem. Permet de moins restreindre l'alimentation et améliore profil lipidique/glycémique.
\end{itemize}
\textbf{Conclusion} : Obésité (IMC 33,9) due à BE $+\SI{4200}{kJ/j}$. Correction durable = restriction modérée ciblée + activité physique régulière.
\end{exoresolu}

\courssub{Adaptations métaboliques à l'effort et à la privation alimentaire}

\begin{propriete}[Adaptations métaboliques à l'effort musculaire]
L'effort impose une demande en ATP brutale. Trois systèmes s'enchaînent selon l'intensité/durée :
\begin{center}
\begin{tabularx}{\linewidth}{|>{\bfseries}l|X|X|X|}\hline
\rowcolor{popblueL} \textbf{Système} & \textbf{Alactique anaérobie (Phosphagène)} & \textbf{Lactique anaérobie (Glycolyse)} & \textbf{Aérobie (Respiration)} \\ \hline
\textbf{Durée / Intensité} & $< 10$ s / Maximale & 10 s - 2 min / Sous-maximale & $> 2-3$ min / Modérée à forte \\ \hline
\textbf{Substrat} & Phosphocréatine (CrP) + ATP stocké & Glycogène musculaire / Glucose sanguin & Glycogène, Glucose, Acides gras, Protéines \\ \hline
\textbf{O$_2$ requis} & Non & Non & \textbf{Oui} \\ \hline
\textbf{Rendement ATP} & Très faible (stocks limités) & Faible (2 ATP/glucose) & Élevé (36-38 ATP/glucose) \\ \hline
\textbf{Produits} & Cr $\rightarrow$ CrP (récup) & \textbf{Lactate + H$^+$} (acidose) & \ce{CO2 + H2O} \\ \hline
\textbf{Fatigue} & Épuisement CrP & Acidose (inhibition PFK, Ca$^{2+}$) & Épuisement glycogène, déshydratation, hyperthermie \\ \hline
\end{tabularx}
\end{center}
\textbf{Quotient Respiratoire (QR) à l'effort} : Au repos QR $\approx 0,8$. Effort modéré QR $\rightarrow 1,0$ (glucides). Effort très intense QR $> 1,0$ (hyperventilation + tamponnage lactate $\rightarrow$ \ce{CO2} excédentaire).
\end{propriete}

\begin{methode}[Lire et interpréter un graphique d'échanges gazeux ($\dot{V}\ce{O2}$, $\dot{V}\ce{CO2}$)]
\begin{enumerate}
\item \textbf{Lire les axes} : Temps (min), $\dot{V}\ce{O2}$ et $\dot{V}\ce{CO2}$ (L/min ou mL/min/kg). Repérer phases : Repos, Effort, Récupération.
\item \textbf{Décrire chiffré} : Valeurs repos, cinétique montée (déficit O$_2$), plateau (régime permanent / $\dot{V}\ce{O2}$ max), cinétique récupération (dette O$_2$). Citer au moins 3 valeurs numériques.
\item \textbf{Comparer courbes} : $\dot{V}\ce{CO2} > \dot{V}\ce{O2}$ au début effort (tamponnage lactate) ou fin effort intense. QR $= \dot{V}\ce{CO2}/\dot{V}\ce{O2}$.
\item \textbf{Interpréter} :
\begin{itemize}
\item Montée $\dot{V}\ce{O2}$ : Augmentation ventilation/débit cardiaque $\rightarrow$ apport O$_2$ muscles.
\item Plateau $\dot{V}\ce{O2}$ : $\dot{V}\ce{O2}$ max (puissance aérobie max) ou régime permanent (seuil lactique).
\item $\dot{V}\ce{O2}$ élevée récupération : \textbf{Paiement dette O$_2$} (réoxydation lactate $\rightarrow$ pyruvate/glucose, resynthèse CrP/ATP, réoxygénation Mb/Hb, thermorégulation).
\end{itemize}
\end{enumerate}
\textbf{Réflexe} : Décrire \textbf{avant} d'interpréter. « La courbe montre... » $\rightarrow$ « Cela traduit... ».
\end{methode}

\begin{exoresolu}[Analyse cinétique $\dot{V}\ce{O2}$]
Données : Repos 0,25 ; 2 min 1,8 ; 5 min 2,8 ; 10 min (fin) 3,0 ; 12 min 1,5 ; 20 min 0,30 (L/min).
\begin{enumerate}
\item Décris l'évolution.
\item Rapport max/repos + interprétation.
\item Pourquoi $\dot{V}\ce{O2}$ élevée à 12 min ?
\end{enumerate}
\tcblower
\textbf{1.} Repos : \SI{0,25}{L/min}. Montée rapide : \SI{1,8}{L/min} (2 min), \SI{2,8}{L/min} (5 min). Plateau effort : \SI{3,0}{L/min} (10 min). Récupération : \SI{1,5}{L/min} (12 min), retour quasi basal \SI{0,30}{L/min} (20 min).
\textbf{2.} Rapport $= 3,0 / 0,25 = 12$. Multiplié par 12. Besoins ATP musculaires $\uparrow \uparrow$ $\rightarrow$ respiration cellulaire $\uparrow \uparrow$ $\rightarrow$ ventilation/débit cardiaque $\uparrow \uparrow$.
\textbf{3.} \textbf{Dette d'oxygène (EPOC)}. Pendant effort, apport O$_2$ < demande $\rightarrow$ fermentation lactique (lactate $\uparrow$). Post-effort, $\dot{V}\ce{O2}$ supra-basale sert à : (1) Oxydation lactate (foie/cœur/muscle), (2) Resynthèse CrP/ATP, (3) Réoxygénation Myoglobine/Hémoglobine, (4) Thermorégulation.
\textbf{Conclusion} : Courbe type adaptation aérobie + dette O$_2$ significative (effort intense).
\end{exoresolu}

\begin{methode}[Analyser l'adaptation métabolique au jeûne (privation alimentaire)]
\begin{enumerate}
\item \textbf{Ordonner les phases chronologiques} (hiérarchisation des réserves pour sauver la glycémie/cerveau) :
\begin{enumerate}
\item \textbf{Phase 1 (0-24 h) : Post-absorptif / Glycogénolyse.} Baisse glycémie $\rightarrow$ $\downarrow$ Insuline, $\uparrow$ \textbf{Glucagon}. Glycogénolyse hépatique $\rightarrow$ Glucose sanguin. Glycogène musculaire $\rightarrow$ usage local. Perte masse rapide (glycogène + 3g eau/g).
\item \textbf{Phase 2 (1-3 à 10-14 j) : Lipolyse / Cétonémie.} Glycogène hépatique épuisé. $\uparrow\uparrow$ Glucagon, $\downarrow\downarrow$ Insuline $\rightarrow$ Lipolyse adipose (TG $\rightarrow$ AG + Glycérol). Foie : $\beta$-oxydation AG $\rightarrow$ Acétyl-CoA $\rightarrow$ \textbf{Corps cétoniques} (acétoacétate, $\beta$-hydroxybutyrate, acétone). Cerveau utilise corps cétoniques ($\uparrow$ à 50-70\% besoins) $\rightarrow$ épargne glucose/protéines. Haleine acétone.
\item \textbf{Phase 3 (Prolongé $> 2$ sem) : Protéolyse / Détresse.} Réserves lipidiques basses. Protéolyse musculaire massive $\rightarrow$ Acides aminés $\rightarrow$ Néoglucogenèse (foie/rein) + Énergie. $\uparrow\uparrow$ Urémie. Fonte musculaire (cœur, diaphragme) $\rightarrow$ décès.
\end{enumerate}
\item \textbf{Hormone clé} : \textbf{Glucagon} (cellules $\alpha$ pancréas). Stimulus = \textbf{Hypoglycémie}. Antagoniste : Insuline.
\item \textbf{Marqueurs biologiques} : Glycémie maintenue ($\approx 0,7-0,8$ g/L), Cétonémie $\uparrow\uparrow$, Urémie $\uparrow$ (protéolyse), $\downarrow$ Insulinémie.
\end{enumerate}
\textbf{Réflexe} : L'organisme ne « s'arrête pas ». MB continue. Le cerveau \textbf{doit} avoir du glucose (ou corps cétoniques). La néoglucogenèse est vitale.
\end{methode}

\begin{popfigure}
\begin{tikzpicture}[
    node distance=1.5cm,
    phase/.style={draw=popdark, thick, rounded corners, fill=popcream, minimum width=3.5cm, minimum height=2.5cm, align=center},
    arrow/.style={-Stealth, thick, popdark},
    label/.style={font=\small, text width=3cm, align=center}
]
\node[phase, align=center] (p1){\textbf{Phase 1 : 0-24h}\\\textbf{Glycogénolyse hépatique}\\Glycogène $\to$ Glucose\\\textit{Hormone : Glucagon $\uparrow$}\\\textit{Perte masse : Eau + Glycogène}};
\node[phase, right=of p1, align=center] (p2){\textbf{Phase 2 : Jours 1-14}\\\textbf{Lipolyse + Cétonémie}\\\ce{TG -> AG + Glyc}\\AG $\to$ Corps cétoniques\\\textit{Cerveau utilise cétones}\\\textit{Épargne protéines}};
\node[phase, right=of p2, align=center] (p3){\textbf{Phase 3 : $>$ 2 sem}\\\textbf{Protéolyse sévère}\\Protéines $\to$ AA $\to$ Néoglucogenèse\\\textit{Urémie $\uparrow\uparrow$}\\\textit{Fonte musculaire (cœur), risque vital}};

\draw[arrow] (p1) -- node[above, label] {Épuisement glycogène} (p2);
\draw[arrow] (p2) -- node[above, label] {Réserves lipides basses} (p3);
\end{tikzpicture}
\legende{Phases successives de l'adaptation au jeûne prolongé : priorité au maintien de la glycémie pour le cerveau via le glucagon.}
\end{popfigure}

\begin{exoresolu}[Jeûne expérimental 5 jours]
Volontaire sain, jeûne 5 jours. Paramètres sanguins :
\begin{center}
\begin{tabular}{|c|c|c|c|}\hline
\rowcolor{popblueL} Paramètre & Jour 0 & Jour 2 & Jour 5 \\ \hline
Glycémie (g/L) & 0,95 & 0,80 & 0,75 \\ \hline
Corps cétoniques (mmol/L) & 0,1 & 1,5 & 4,0 \\ \hline
Urée sanguine (g/L) & 0,25 & 0,30 & 0,45 \\ \hline
\end{tabular}
\end{center}
\begin{enumerate}
\item Décris l'évolution de chaque paramètre.
\item Explique l'origine de l'augmentation des corps cétoniques et de l'urée.
\item Hormone responsable et stimulus.
\end{enumerate}
\tcblower
\textbf{1. Évolutions :}
\begin{itemize}
\item Glycémie : Légère baisse (0,95 $\rightarrow$ 0,75 g/L), maintenue dans la norme (homéostasie).
\item Corps cétoniques : Hausse massive (0,1 $\rightarrow$ 4,0 mmol/L, $\times 40$). Signature de la phase 2.
\item Urémie : Hausse modérée (0,25 $\rightarrow$ 0,45 g/L). Début phase 3 (protéolyse mesurable).
\end{itemize}
\textbf{2. Origines :}
\begin{itemize}
\item \textbf{Corps cétoniques} : Épuisement glycogène hépatique (J1) $\rightarrow$ Lipolyse adipose massive $\rightarrow$ Flux AG vers foie $\rightarrow$ $\beta$-oxydation $\rightarrow$ Excès Acétyl-CoA $\rightarrow$ Cétonégénèse (HMG-CoA synthase). Carburant substitut cerveau/cœur.
\item \textbf{Urée} : Protéolyse musculaire $\rightarrow$ Acides aminés $\rightarrow$ Désamination hépatique (transamination/désamination oxydative) $\rightarrow$ NH$_3$ $\rightarrow$ Cycle de l'urée $\rightarrow$ Urémie $\uparrow$. Squelettes carbonés $\rightarrow$ Néoglucogenèse (maintien glycémie).
\end{itemize}
\textbf{3. Hormone :} \textbf{Glucagon} (cellules $\alpha$ îlots Langerhans). Stimulus : \textbf{Baisse glycémie} (hypoglycémie relative) détectée par cellules $\alpha$. Effets : Glycogénolyse, Néoglucogenèse, Lipolyse, Cétonégénèse.
\textbf{Conclusion} : Adaptation hiérarchisée (Glycogène $\rightarrow$ Lipides/Cétones $\rightarrow$ Protéines) sous contrôle glucagonique. Jour 5 = transition Phase 2/3.
\end{exoresolu}

\courssub{Mesures correctives, prévention et éducation nutritionnelle}

\begin{aretenir}[Principes de la correction des déséquilibres énergétiques]
\begin{itemize}
\item \textbf{Objectif unique} : Ramener le Bilan Énergétique vers 0 (équilibre) de façon \textbf{durable}.
\item \textbf{Déséquilibre Positif (Surpoids/Obésité)} :
\begin{itemize}
\item \textit{Apports} : Déficit modéré (\SI{-500}{kJ} à \SI{-2000}{kJ/jour}). Réduire densité énergétique (lipides cachés, sucres ajoutés, alcool). Augmenter volume (légumes, fibres, eau) pour satiété. \textbf{Ne pas descendre sous MB}.
\item \textit{Dépenses} : Activité physique \textbf{régulière} (endurance + renforcement). Cible : \SI{300}{kJ} à \SI{1000}{kJ}/séance. Augmenter NEAT (thermogenèse non-exercice : marche, escaliers).
\item \textit{Comportement} : Fractionnement (3 repas + 1 collation), manger lentement, écouter satiété, sommeil suffisant (leptine/ghréline), gestion stress (cortisol).
\end{itemize}
\item \textbf{Déséquilibre Négatif (Dénutrition/Maigreur)} :
\begin{itemize}
\item \textit{Apports} : Excédent progressif (+\SI{500}{kJ} à \SI{1500}{kJ/jour}). Enrichir ration (huiles, pâtes d'arachide, lait en poudre, œufs) sans augmenter volume excessif. Fractionner (5-6 prises). Supplémentation micronutriments (vit A, Zn, Fe, iodure).
\item \textit{Pathologie} : Traiter cause (infection, malabsorption, anorexie, dépression).
\item \textit{Réhabilitation} : Réintroduction progressive (risque syndrome de réalimentation : hypophosphatémie, œdème, insuffisance cardiaque).
\end{itemize}
\item \textbf{Prévention (Niveau population - Cameroun)} : Promotion alimentation locale diversifiée (feuilles, légumineuses, céréales complètes, poisson, fruits), lutte contre publicité aliments ultra-transformés, urbanisme actif (pistes cyclables, espaces sport), éducation scolaire (cantine équilibrée, EPS), surveillance croissance (carnets de santé, IMC école).
\end{itemize}
\end{aretenir}

\begin{experience}[TP : Détermination de la valeur énergétique d'un aliment par calorimétrie (bombe calorimétrique simplifiée)]
\textbf{Objectif} : Mesurer l'énergie libérée par la combustion d'un aliment (ex. cacahuète) et comparer à la valeur théorique.
\textbf{Matériel} : Boîte métallique (calorimètre), thermomètre (0,1°C), support, aiguille/liège, cacahuète pesée ($m_a$), eau (masse $m_e$ connue, ex. \SI{50}{g}), allumette.
\textbf{Protocole} :
\begin{enumerate}
\item Mesurer $T_i$ de l'eau.
\item Enflammer la cacahuète, la placer sous la boîte immédiatement.
\item Remuer l'eau, mesurer $T_f$ max après combustion complète.
\item Peser résidu ($m_r$). Masse brûlée $m_b = m_a - m_r$.
\end{enumerate}
\textbf{Calculs} :
Énergie absorbée par l'eau : $Q = m_e \times c_e \times (T_f - T_i)$ ($c_e = \SI{4,18}{J/g/^\circ C}$).
Valeur énergétique mesurée : $V_{e,mes} = \frac{Q}{m_b}$ (en J/g $\rightarrow$ kJ/g).
\textbf{Interprétation} : $V_{e,mes} < V_{e,théo}$ (pertes chaleur air/boîte, combustion incomplète, eau évaporée). Discuter sources d'erreur.
\textbf{Lien cours} : Illustre la conversion énergie chimique $\rightarrow$ chaleur. Base des valeurs de référence (Atwater) utilisées dans les bilans.
\end{experience}

\begin{exercice}[Entraînement Probatoire - Bilan \& Adaptations]
\textbf{Exercice 1.} Une jeune fille de 16 ans, 55 kg, 1,60 m, pratique la danse 1h/jour (DEA $\approx \SI{1500}{kJ}$). Son MB est estimé à \SI{5500}{kJ/j}. Ses apports : 250 g G, 60 g L, 70 g P ($V_e$ : G=17, L=38, P=17 kJ/g).
\begin{enumerate}
\item Calcule son IMC et qualifie son état.
\item Calcule ses Apports, Dépenses totales et Bilan.
\item Prédis l'évolution de sa masse à 3 mois si ce rythme persiste.
\end{enumerate}

\textbf{Exercice 2.} Lors d'un test d'effort, un cycliste présente un QR = 0,95 à 150 W et QR = 1,05 à 300 W.
\begin{enumerate}
\item Interprète le substrat principal à 150 W.
\item Explique le QR $> 1$ à 300 W (mécanisme biochimique précis).
\end{enumerate}

\textbf{Exercice 3.} Un enfant de 3 ans en zone rurale présente : masse 9 kg (référence 14 kg), œdèmes membres inférieurs, dermatose, cheveux décolorés, apathie.
\begin{enumerate}
\item Identifie le type de dénutrition (nom clinique).
\item Explique le mécanisme physiopathologique de l'œdème (protéines plasmatiques).
\item Propose une prise en charge nutritionnelle d'urgence (type aliment, précautions).
\end{enumerate}
\end{exercice}

\begin{corrige}[Exercice 1]
\textbf{1.} $IMC = \frac{55}{1,60^2} = \frac{55}{2,56} \approx 21,5$. \textbf{Normal} (18,5-25).
\textbf{2.} AE $= 250\times17 + 60\times38 + 70\times17 = 4250 + 2280 + 1190 = \SI{7720}{kJ}$.
DE $= MB + DEA = 5500 + 1500 = \SI{7000}{kJ}$.
BE $= 7720 - 7000 = +\SI{720}{kJ/j}$.
\textbf{3.} Bilan positif modéré. Stockage $\approx \SI{720}{kJ/j}$. Sur 90 jours : $90 \times 720 = \SI{64800}{kJ}$. Gain masse adipeuse $\approx \frac{64800}{30000} \approx \SI{2,2}{kg}$. Prise de masse prévisible $\approx$ 2 kg (majoritairement graisse). Risque surpoids à long terme si inactivité augmente.
\end{corrige}

\begin{corrige}[Exercice 2]
\textbf{1.} QR = 0,95 $\approx$ 1,0 $\rightarrow$ Oxydation prédominante des \textbf{glucides} (glycogène musculaire + glucose sanguin). Effort modéré/soutenu.
\textbf{2.} QR = 1,05 $> 1$. À haute intensité (300 W), glycolyse rapide $\rightarrow$ Lactate$^-$ + H$^+$ accumulés. Tamponnage : $\ce{H+ + HCO3- -> H2CO3 -> CO2 + H2O}$. $\ce{CO2}$ produit \textbf{non métabolique} (issu bicarbonate) s'ajoute au $\ce{CO2}$ métabolique $\rightarrow$ $\dot{V}\ce{CO2}$ artificiellement augmentée par rapport à $\dot{V}\ce{O2}$. Hyperventilation compensatoire (excrétion $\ce{CO2}$) amplifie le phénomène. Ne signifie pas oxydation de substrat « sur-oxydé ».
\end{corrige}

\begin{corrige}[Exercice 3]
\textbf{1.} \textbf{Kwashiorkor} (dénutrition protéinée avec apport calorique relatif maintenu par glucides). Signes : œdèmes, dermatose « peau craquelée », cheveux décolorés (dyschromie), hépatomégalie (stéatose), apathie. Marasme = amaigrissement sévère sans œdème (déficit global énergétique+protéique).
\textbf{2.} Hypoprotéinémie sévère (synthèse hépatique albumine $\downarrow$ par manque AA essentiels) $\rightarrow$ Pression oncotique plasmatique $\downarrow$ $\rightarrow$ Fuite liquide intravasculaire $\rightarrow$ interstitiel $\rightarrow$ \textbf{Œdèmes} (démarrant membres inférieurs). Vasodilatation NO-dépendante aggrave.
\textbf{3.} \textbf{Urgence} : Lait thérapeutique F-75 (faible protéines, faible osmolalité, riche K, Mg, Zn, Cu, vitamines) par sonde nasogastrique ou gobelet. \textbf{Phase de stabilisation} (2-7 j) : \SI{100}{kcal/kg/j} (\SI{420}{kJ/kg/j}), volume \SI{130}{mL/kg/j}. Surveillance glycémie, T°, signes réalimentation. \textbf{Puis réhabilitation} : F-100 / RUTF (Plumpy'nut) progressif. Traitement infections associées (antibiotiques systémiques), vermifuge, vitamine A, folates.
\end{corrige}

\begin{attention}[Les 5 erreurs fatales au Probatoire (SVT 1ère D)]
\begin{enumerate}
\item \textbf{Signe du bilan inversé} : Écrire « Bilan positif = bonne santé » ou « Bilan négatif = perte de poids donc c'est bien ». \textbf{Règle} : BE = Apports - Dépenses. Positif = Stockage (Obésité). Négatif = Déstockage (Maigreur/Dénutrition). Toujours expliciter le sens physiologique du signe.
\item \textbf{Unités mélangées / Oubli conversion masse $\rightarrow$ énergie} : Additionner 400 g glucides + 90 g lipides sans multiplier par $V_e$. Ou additionner kJ et kcal. \textbf{Règle} : Tableau de conversion systématique. $1\ \text{kcal} = \num{4,18}\ \text{kJ}$.
\item \textbf{Confusion Fermentation / Respiration} : « Fermentation produit \ce{O2} » ou « Respiration sans \ce{O2} » ou « Fermentation = 38 ATP ». \textbf{Règle} : Fermentation = Glycolyse seule (2 ATP), sans \ce{O2}, produit Lactate OU Éthanol+\ce{CO2}. Respiration = Glycolyse + Krebs + Chaîne respiratoire (36-38 ATP), avec \ce{O2}.
\item \textbf{Description de courbe sans chiffres} : « La courbe monte puis descend ». \textbf{Règle} : Citer valeurs (ex: « $\dot{V}\ce{O2}$ passe de 0,25 à 3,0 L/min entre 0 et 10 min »), pentes, plateaux. \textbf{Puis} interpréter.
\item \textbf{Ordre des réserves au jeûne inversé / Oubli Glucagon} : « Le muscle fond dès le 1er jour » ou « Le foie fait du glucose avec des graisses ». \textbf{Règle} : Ordre strict : Glycogène hépatique (J1) $\rightarrow$ Lipides/Corps cétoniques (J1-J14) $\rightarrow$ Protéines musculaires (fin). Hormone pilote = \textbf{Glucagon} (stimulus = hypoglycémie). Cerveau \textbf{ne peut pas} utiliser AG directement (barrière hémato-encéphalique), utilise corps cétoniques.
\end{enumerate}
\end{attention}

\newpage

\coursec{Exercices}

\courssub{Je vérifie mes connaissances}

\begin{exercice}[QCM : les voies de production d'énergie]
Pour chaque affirmation, choisis la ou les bonne(s) réponse(s) en justifiant brièvement ton choix.
\begin{enumerate}
\item La respiration cellulaire aérobie se déroule principalement :
\begin{enumerate}
\item dans le cytoplasme ; \item dans le noyau ; \item dans les mitochondries ; \item dans les lysosomes.
\end{enumerate}
\item Le rendement énergétique de l'oxydation complète d'une molécule de glucose en présence d'oxygène est d'environ :
\begin{enumerate}
\item 2 ATP ; \item 8 ATP ; \item 36 à 38 ATP ; \item 100 ATP.
\end{enumerate}
\item La fermentation lactique chez le muscle en effort intense produit :
\begin{enumerate}
\item de l'éthanol et du \ce{CO2} ; \item de l'acide lactique ; \item de l'eau et du \ce{CO2} ; \item de l'urée.
\end{enumerate}
\item L'équation globale de la respiration cellulaire est :
\begin{enumerate}
\item \ce{C6H12O6 + 6O2 -> 6CO2 + 6H2O} ; \item \ce{6CO2 + 6H2O -> C6H12O6 + 6O2} ;
\item \ce{C6H12O6 -> 2C3H6O3} ; \item \ce{C6H12O6 -> 2C2H5OH + 2CO2}.
\end{enumerate}
\end{enumerate}
\end{exercice}

\begin{exercice}[Vrai ou faux, à justifier]
Indique si chacune des affirmations suivantes est vraie ou fausse, puis justifie ta réponse en une ou deux phrases.
\begin{enumerate}
\item « Un bilan énergétique positif prolongé entraîne obligatoirement une perte de masse corporelle. »
\item « Au repos, l'organisme d'un adulte dépense de l'énergie uniquement pour ses mouvements volontaires. »
\item « Le quotient respiratoire (QR) égal à 1 indique que le substrat majoritairement oxydé est le glucose. »
\item « Lors d'un jeûne prolongé, l'organisme peut fabriquer du glucose à partir d'acides aminés : c'est la néoglucogenèse. »
\item « La fermentation permet de produire autant d'ATP que la respiration, mais plus lentement. »
\end{enumerate}
\end{exercice}

\begin{exercice}[Définitions éclair]
Définis en deux lignes maximum chacun des termes suivants, en donnant un exemple concret pour chacun :
\begin{enumerate}
\item métabolisme de base ;
\item bilan énergétique ;
\item déséquilibre énergétique négatif ;
\item quotient respiratoire ;
\item néoglucogenèse.
\end{enumerate}
\end{exercice}

\begin{exercice}[Calculs directs]
\begin{enumerate}
\item Un élève ingère en une journée \SI{12000}{kJ} et dépense \SI{10500}{kJ}. Calcule son bilan énergétique journalier en kJ, puis en kcal (rappel : \SI{1}{kcal} = \SI{4,18}{kJ}). Ce bilan est-il positif ou négatif ?
\item Sachant que \SI{1}{g} de lipides stockés correspond à environ \SI{38}{kJ}, estime la masse de réserve lipidique théoriquement accumulée en 30 jours si ce bilan se maintient chaque jour.
\item Un muscle au repos consomme \SI{0,25}{L} de \ce{O2} par minute. Calcule sa consommation en une heure, puis le volume de \ce{CO2} produit si le QR vaut 0,85.
\end{enumerate}
\end{exercice}

\courssub{Je m'entraîne}

\begin{exercice}[Bilan énergétique d'un adolescent de 17 ans]
Marius, élève de Première D à Yaoundé, âgé de 17 ans (masse : \SI{68}{kg}), a noté sa ration alimentaire d'une journée typique et son emploi du temps. Les données sont rassemblées ci-dessous.

\begin{center}
\rowcolors{1}{popblueL}{white}
\begin{tabularx}{\linewidth}{|l|X|c|}
\hline
\rowcolor{popblue!40} \textbf{Repas} & \textbf{Composition} & \textbf{Énergie (kJ)} \\
\hline
Petit-déjeuner & 2 beignets haricots + bouillie de maïs sucrée & 2 400 \\
\hline
Déjeuner & Ndolé (viande, arachides, feuilles) + 300 g de riz + 1 banane mûre & 4 300 \\
\hline
Collation & 1 boisson sucrée (33 cL) + 1 beignet & 1 600 \\
\hline
Dîner & Poisson braisé (150 g) + bâton de manioc + 1 mangue & 3 600 \\
\hline
\end{tabularx}
\end{center}

Dépenses énergétiques estimées : métabolisme de base \SI{7100}{kJ/jour} ; marche à pied (aller-retour école, 50 min) \SI{900}{kJ} ; activités scolaires et domestiques \SI{1800}{kJ} ; match de football (45 min) \SI{1600}{kJ}.

\begin{enumerate}
\item Calcule les apports énergétiques totaux de la journée.
\item Calcule les dépenses énergétiques totales de la journée.
\item Établis le bilan énergétique de Marius et conclus sur le risque de surpoids à long terme si ce mode de vie se maintient.
\item Identifie dans le menu les deux postes d'apport les plus énergétiques et propose une modification réaliste pour chacun.
\end{enumerate}
\end{exercice}

\begin{exercice}[Test d'effort progressif et quotient respiratoire]
Lors d'un test d'effort sur cycloergomètre réalisé dans un laboratoire de physiologie, on mesure chez un sujet les échanges gazeux à trois intensités d'exercice. Les résultats sont les suivants :

\begin{center}
\rowcolors{1}{popgreenL}{white}
\begin{tabularx}{\linewidth}{|l|c|c|X|}
\hline
\rowcolor{popgreen!40} \textbf{Intensité} & \textbf{\ce{O2} consommé (L/min)} & \textbf{\ce{CO2} produit (L/min)} & \textbf{QR = \ce{CO2}/\ce{O2}} \\
\hline
Repos & 0,30 & 0,24 & \\
\hline
Effort modéré (100 W) & 1,50 & 1,35 & \\
\hline
Effort intense (220 W) & 2,80 & 2,94 & \\
\hline
\end{tabularx}
\end{center}

\begin{enumerate}
\item Calcule le quotient respiratoire à chacune des trois intensités (arrondis au centième).
\item Sachant que QR $\approx 0,7$ correspond à l'oxydation des lipides, QR $\approx 1$ à celle des glucides et QR $> 1$ à une production supplémentaire de \ce{CO2} liée au tamponnement de l'acide lactique, déduis les substrats majoritairement oxydés à chaque intensité.
\item Explique pourquoi le QR dépasse 1 à l'effort intense.
\item Trace l'allure de la courbe du QR en fonction de l'intensité de l'effort (axes légendés).
\end{enumerate}
\end{exercice}

\begin{exercice}[Interprétation d'une courbe de variation pondérale]
On suit la masse corporelle d'une jeune femme pendant 12 semaines. Les relevés hebdomadaires sont les suivants :

\begin{center}
\begin{tabular}{|l|c|c|c|c|c|c|c|c|c|c|c|c|}
\hline
Semaine & 0 & 1 & 2 & 3 & 4 & 5 & 6 & 7 & 8 & 9 & 10 & 11 \\
\hline
Masse (kg) & 58 & 58,8 & 59,6 & 60,4 & 61,2 & 62 & 63 & 64 & 62,8 & 61,8 & 60,9 & 61 \\
\hline
\end{tabular}
\end{center}

\begin{enumerate}
\item Trace la courbe de variation de la masse en fonction du temps (échelles au choix, axes légendés).
\item Identifie sur la courbe les deux phases distinctes et précise pour chacune le signe du bilan énergétique.
\item Calcule la prise de masse totale de la phase 1 et la perte de masse de la phase 2.
\item Propose une explication physiologique de la prise de masse (nature du tissu stocké) et de la perte de masse (pourquoi la perte n'est-elle jamais constituée uniquement de lipides ?).
\item La perte de masse de la phase 2 est de \SI{3}{kg} en 4 semaines. En supposant que 70 \% de cette perte provient des lipides (\SI{38}{kJ/g}), estime le déficit énergétique journalier moyen.
\end{enumerate}
\end{exercice}

\begin{exercice}[Programme de rééquilibrage]
Sandrine, 16 ans, mesure \SI{1,60}{m} et pèse \SI{69}{kg}. Son médecin scolaire lui conseille un rééquilibrage sur 8 semaines. Son métabolisme de base est estimé à \SI{5900}{kJ/jour} et ses dépenses d'activité actuelles à \SI{2100}{kJ/jour}. Ses apports actuels sont de \SI{10500}{kJ/jour}.

\begin{enumerate}
\item Calcule l'IMC de Sandrine (rappel : IMC = masse/taille$^2$). À l'aide des classes de référence (maigreur $< 18,5$ ; normal : 18,5 à 25 ; surpoids : 25 à 30 ; obésité $> 30$), situe son état nutritionnel.
\item Calcule son bilan énergétique journalier actuel.
\item Propose un objectif calorique quotidien réaliste (déficit modéré de l'ordre de \SI{1500}{kJ/jour}) et précise la répartition conseillée des macronutriments (glucides 50 à 55 \%, lipides 30 à 35 \%, protéines 12 à 15 \% de l'énergie totale) en kJ pour chaque catégorie.
\item Propose un type et une fréquence d'activité physique adaptés à une élève de Première, en justifiant.
\item Indique trois indicateurs de suivi à relever chaque semaine pour évaluer l'efficacité du programme sans tomber dans les excès.
\end{enumerate}
\end{exercice}

\courssub{Je me prépare au Probatoire}

\begin{exercice}[Sujet type Probatoire : le jeûne de 48 heures]
\textbf{Contexte.} Dans le cadre d'une étude physiologique menée dans un hôpital universitaire, un volontaire sain de 22 ans (masse : \SI{70}{kg}) jeûne pendant 48 heures sous surveillance médicale. On mesure l'évolution de plusieurs paramètres sanguins :

\begin{center}
\rowcolors{1}{poporangeL}{white}
\begin{tabularx}{\linewidth}{|X|c|c|c|}
\hline
\rowcolor{poporange!40} \textbf{Paramètre} & \textbf{Début du jeûne (0 h)} & \textbf{24 h} & \textbf{48 h} \\
\hline
Glycémie (g/L) & 0,95 & 0,82 & 0,78 \\
\hline
Insulinémie (unités arbitraires) & 12 & 6 & 4 \\
\hline
Glucagon (unités arbitraires) & 8 & 15 & 20 \\
\hline
Corps cétoniques (mmol/L) & 0,1 & 0,8 & 2,5 \\
\hline
\end{tabularx}
\end{center}

On rappelle que les réserves de glycogène hépatique d'un adulte (environ \SI{100}{g}, soit \SI{1700}{kJ}) sont largement épuisées après 24 h de jeûne, et que le cerveau consomme normalement environ \SI{120}{g} de glucose par jour.

\begin{enumerate}
\item Définis la glycémie et précise sa valeur normale à jeun.
\item À partir des données du tableau, décris l'évolution de la glycémie au cours du jeûne et montre qu'elle reste dans des limites compatibles avec la survie : quel nom donne-t-on à ce type de régulation ?
\item Explique, en t'appuyant sur les variations hormonales observées, les mécanismes qui maintiennent la glycémie pendant les 24 premières heures (rôle du foie et du glycogène).
\item Au-delà de 24 h, quelles voies métaboliques prennent le relais ? Nomme précisément les substrats mobilisés (tissu adipeux, muscles) et les organes impliqués.
\item L'alanine, acide aminé libéré par les muscles, est un précurseur majeur de la néoglucogenèse hépatique. Complète l'équation globale simplifiée de la synthèse du glucose à partir de l'alanine : \ce{2 C3H7NO2 + 2 H2O -> C6H12O6 + ...} (identifie le produit azoté formé et précise le nom de la voie qui l'élimine dans l'urée).
\item Interprète l'augmentation des corps cétoniques à 48 h : quelle est leur origine et quelle économie réalisent-ils pour l'organisme ?
\end{enumerate}
\end{exercice}

\begin{exercice}[Sujet type Probatoire : effort sportif et métabolisme au Mont Cameroun]
\textbf{Contexte.} Lors de la course d'endurance « Course de l'Espoir » au Mont Cameroun, des chercheurs mesurent chez deux athlètes la consommation d'\ce{O2} et la production de \ce{CO2} avant le départ, pendant l'ascension et à l'arrivée. L'athlète A est bien entraîné ; l'athlète B est peu entraîné.

\begin{center}
\rowcolors{1}{popblueL}{white}
\begin{tabularx}{\linewidth}{|X|c|c|c|}
\hline
\rowcolor{popblue!40} \textbf{Mesure} & \textbf{Repos} & \textbf{Ascension (effort intense)} & \textbf{Arrivée (récupération 10 min)} \\
\hline
\ce{O2} consommé par A (L/min) & 0,28 & 3,60 & 0,90 \\
\hline
\ce{CO2} produit par A (L/min) & 0,22 & 3,42 & 0,85 \\
\hline
\ce{O2} consommé par B (L/min) & 0,30 & 2,90 & 1,40 \\
\hline
\ce{CO2} produit par B (L/min) & 0,24 & 3,10 & 1,35 \\
\hline
\end{tabularx}
\end{center}

\begin{enumerate}
\item Calcule le quotient respiratoire de chaque athlète au repos et pendant l'ascension.
\item Compare les QR pendant l'ascension et indique quel athlète oxyde davantage les lipides. Quel avantage cela représente-t-il pour une course de longue durée (réserves lipidiques $\approx$ \SI{300000}{kJ} contre réserves glucidiques $\approx$ \SI{8000}{kJ}) ?
\item L'athlète B présente un QR supérieur à 1 pendant l'ascension. Interprète ce résultat en faisant intervenir la fermentation lactique et le tamponnement de l'acide lactique par les bicarbonates. Écris l'équation de la fermentation lactique : \ce{C6H12O6 -> 2 C3H6O3}.
\item Après l'arrivée, la consommation d'\ce{O2} reste élevée pendant la récupération, surtout chez B. Nomme ce phénomène et explique son rôle (élimination de l'acide lactique, reconstitution des réserves).
\item Calcule l'énergie dépensée par l'athlète A pendant une ascension de 3 h, sachant que la consommation de \SI{1}{L} d'\ce{O2} libère environ \SI{20}{kJ}.
\item Propose deux adaptations métaboliques et cardiovasculaires acquises par l'entraînement qui expliquent les meilleures performances de A.
\end{enumerate}
\end{exercice}

\begin{exercice}[Sujet type Probatoire : dénutrition et obésité, deux déséquilibres]
\textbf{Contexte.} Une enquête nutritionnelle menée dans deux localités du Cameroun porte sur deux groupes d'adolescents. Le groupe 1 vit dans une zone touchée par l'insécurité alimentaire ; le groupe 2 vit en milieu urbain avec une alimentation riche en produits industriels sucrés et gras et une forte sédentarité.

\begin{center}
\rowcolors{1}{popgreenL}{white}
\begin{tabularx}{\linewidth}{|X|c|c|}
\hline
\rowcolor{popgreen!40} \textbf{Paramètre moyen} & \textbf{Groupe 1} & \textbf{Groupe 2} \\
\hline
Apports énergétiques (kJ/jour) & 6 200 & 12 800 \\
\hline
Dépenses énergétiques (kJ/jour) & 8 400 & 9 600 \\
\hline
IMC moyen (kg/m$^2$) & 16,8 & 27,5 \\
\hline
Masse grasse (\% de la masse corporelle) & 8 & 32 \\
\hline
\end{tabularx}
\end{center}

\begin{enumerate}
\item Calcule le bilan énergétique journalier moyen de chaque groupe et précise son signe.
\item À l'aide des classes d'IMC, caractérise l'état nutritionnel de chaque groupe.
\item Pour le groupe 1 : décris trois conséquences physiologiques attendues d'un bilan énergétique négatif chronique (fonte musculaire, retard de croissance, immunité, cycles hormonaux). Explique pourquoi l'organisme puise d'abord dans le glycogène, puis dans les lipides, puis dans les protéines.
\item Pour le groupe 2 : explique le mécanisme de stockage de l'excédent énergétique et cite trois complications à long terme de l'obésité (diabète de type 2, hypertension artérielle, maladies cardiovasculaires).
\item Le groupe 2 présente un excédent énergétique de \SI{3200}{kJ/jour}. En supposant une perte visée de \SI{0,5}{kg} de masse grasse par semaine (\SI{38}{kJ/g}), calcule le déficit journalier réellement nécessaire et montre qu'une restriction modérée associée à une activité physique suffit.
\item Rédige un court texte (10 à 15 lignes) de sensibilisation destiné aux élèves de ton lycée, expliquant que dénutrition et obésité sont deux faces d'un même problème : le déséquilibre entre apports et dépenses énergétiques, et proposant trois mesures de prévention adaptées au contexte camerounais.
\end{enumerate}
\end{exercice}



\newpage

\coursec{Corrigés des exercices}

\courssub{Je vérifie mes connaissances}

\begin{corrige}[Exercice 1 --- QCM : les voies de production d'énergie]
\begin{enumerate}
\item Réponse \textbf{c} : les mitochondries sont le siège de la respiration cellulaire (cycle de Krebs dans la matrice, chaîne respiratoire dans les crêtes mitochondriales). Seule la glycolyse, première étape commune à toutes les voies, a lieu dans le cytoplasme (hyaloplasme).
\item Réponse \textbf{c} : l'oxydation complète d'une molécule de glucose en présence d'\ce{O2} fournit environ 36 à 38 ATP, contre seulement 2 ATP en fermentation.
\item Réponse \textbf{b} : la fermentation lactique musculaire produit de l'acide lactique (lactate) ; c'est la fermentation alcoolique des levures qui produit éthanol et \ce{CO2}.
\item Réponse \textbf{a} : \ce{C6H12O6 + 6O2 -> 6CO2 + 6H2O}. L'équation b est celle de la photosynthèse ; c est la fermentation lactique ; d est la fermentation alcoolique.
\end{enumerate}
\textbf{Points de vigilance :} ne pas confondre le lieu de la glycolyse (cytoplasme) avec celui de la respiration proprement dite (mitochondries) ; ne pas inverser les deux types de fermentation.
\end{corrige}

\begin{corrige}[Exercice 2 --- Vrai ou faux, à justifier]
\begin{enumerate}
\item \textbf{Faux.} Un bilan énergétique positif (apports $>$ dépenses) entraîne le stockage de l'excédent, principalement sous forme de triglycérides dans le tissu adipeux, donc une \emph{prise} de masse.
\item \textbf{Faux.} Au repos, le métabolisme de base (activité cardiaque et respiratoire, maintien de la température corporelle, renouvellement cellulaire) représente la plus grande part des dépenses, sans aucun mouvement volontaire.
\item \textbf{Vrai.} D'après l'équation \ce{C6H12O6 + 6O2 -> 6CO2 + 6H2O}, l'oxydation du glucose donne QR $= \ce{CO2}/\ce{O2} = 6/6 = 1$.
\item \textbf{Vrai.} La néoglucogenèse hépatique synthétise du glucose à partir de précurseurs non glucidiques (acides aminés comme l'alanine, lactate, glycérol) lorsque le glycogène est épuisé.
\item \textbf{Faux.} La fermentation ne fournit que 2 ATP par glucose contre 36 à 38 en respiration : elle est à la fois moins rentable et limitée dans le temps par l'accumulation d'acide lactique.
\end{enumerate}
\end{corrige}

\begin{corrige}[Exercice 3 --- Définitions éclair]
\begin{enumerate}
\item \textbf{Métabolisme de base :} dépense énergétique minimale de l'organisme au repos, à jeun, nécessaire au maintien des fonctions vitales (cœur, respiration, température). Exemple : environ \SI{7100}{kJ/jour} pour un adolescent de \SI{68}{kg}.
\item \textbf{Bilan énergétique :} différence entre les apports énergétiques alimentaires et les dépenses totales de l'organisme sur une période donnée. Exemple : $12\,000 - 10\,500 = +\SI{1500}{kJ/jour}$.
\item \textbf{Déséquilibre énergétique négatif :} situation où les dépenses dépassent les apports, obligeant l'organisme à puiser dans ses réserves (glycogène, lipides, puis protéines). Exemple : jeûne, dénutrition.
\item \textbf{Quotient respiratoire (QR) :} rapport du volume de \ce{CO2} produit au volume d'\ce{O2} consommé ; il renseigne sur le substrat oxydé (QR $\approx 0,7$ : lipides ; QR $\approx 1$ : glucides). Exemple : QR $= 0,24/0,30 = 0,80$ au repos.
\item \textbf{Néoglucogenèse :} synthèse de glucose par le foie à partir de précurseurs non glucidiques (acides aminés, lactate, glycérol) lors d'un jeûne prolongé. Exemple : fabrication de glucose à partir de l'alanine musculaire après 24 h de jeûne.
\end{enumerate}
\end{corrige}

\begin{corrige}[Exercice 4 --- Calculs directs]
\begin{enumerate}
\item Bilan $=$ apports $-$ dépenses $= 12\,000 - 10\,500 = +\SI{1500}{kJ}$, soit $1500 / 4,18 \approx \SI{359}{kcal}$. Le bilan est \textbf{positif}.
\item Excédent sur 30 jours : $1500 \times 30 = \SI{45000}{kJ}$. Masse de lipides : $45\,000 / 38 \approx \SI{1184}{g} \approx \SI{1,2}{kg}$ de réserve lipidique théorique en un mois.
\item Consommation horaire d'\ce{O2} : $0,25 \times 60 = \SI{15}{L}$. Avec QR $= \ce{CO2}/\ce{O2} = 0,85$ : $\ce{CO2} = 15 \times 0,85 = \SI{12,75}{L}$.
\end{enumerate}
\textbf{Points de vigilance :} bien convertir les kJ en kcal en \emph{divisant} par 4,18 ; dans le calcul du QR, c'est le \ce{CO2} qui est au numérateur.
\end{corrige}

\courssub{Je m'entraîne}

\begin{corrige}[Exercice 5 --- Bilan énergétique d'un adolescent de 17 ans]
\begin{enumerate}
\item Apports totaux : $2\,400 + 4\,300 + 1\,600 + 3\,600 = \SI{11900}{kJ}$.
\item Dépenses totales : $7\,100 + 900 + 1\,800 + 1\,600 = \SI{11400}{kJ}$.
\item Bilan $= 11\,900 - 11\,400 = +\SI{500}{kJ/jour}$ : bilan \textbf{positif}. S'il se maintient, l'excédent mensuel ($500 \times 30 = \SI{15000}{kJ}$, soit $15\,000/38 \approx \SI{395}{g}$ de lipides par mois) conduit à une prise de masse grasse progressive, donc à un risque de surpoids à long terme.
\item Les deux postes les plus énergétiques sont le déjeuner (\SI{4300}{kJ}) et le dîner (\SI{3600}{kJ}). Modifications réalistes : réduire la portion de riz à 200 g et limiter l'huile et les arachides du ndolé ; remplacer la boisson sucrée de la collation par de l'eau et un fruit, et réduire la portion de bâton de manioc au dîner.
\end{enumerate}
\textbf{Points de vigilance :} ne pas oublier le métabolisme de base dans les dépenses ; conclure par le \emph{signe} du bilan et sa conséquence à long terme, pas seulement par le chiffre.
\end{corrige}

\begin{corrige}[Exercice 6 --- Test d'effort progressif et quotient respiratoire]
\begin{enumerate}
\item Repos : $0,24/0,30 = 0,80$. Effort modéré : $1,35/1,50 = 0,90$. Effort intense : $2,94/2,80 = 1,05$.
\item Repos (QR $= 0,80$) : mélange de lipides et de glucides, avec une part importante de lipides. Effort modéré (QR $= 0,90$) : mélange à dominante glucidique. Effort intense (QR $= 1,05$) : glucides quasi exclusifs, avec production d'acide lactique.
\item À l'effort intense, l'apport d'\ce{O2} ne suffit plus : la fermentation lactique produit de l'acide lactique, tamponné par les bicarbonates sanguins (\ce{H+} + \ce{HCO3-} -> \ce{CO2} + \ce{H2O}), ce qui libère du \ce{CO2} supplémentaire non issu de la respiration : le QR dépasse 1.
\item Allure de la courbe :
\end{enumerate}
\begin{popfigure}
\begin{center}
\begin{tikzpicture}
\begin{axis}[
width=0.75\linewidth, height=5.5cm,
xlabel={Intensité de l'effort (W)}, ylabel={QR},
xmin=0, xmax=250, ymin=0.6, ymax=1.2,
xtick={0,100,220}, ytick={0.7,0.8,0.9,1.0,1.1},
grid=both, grid style={popline!40},
]
\addplot[popcurve, mark=*, mark size=1.5pt] coordinates {(0,0.80) (100,0.90) (220,1.05)};
\addplot[dashed, popmuted, domain=0:250] {1};
\end{axis}
\end{tikzpicture}
\end{center}
\legende{Évolution du quotient respiratoire en fonction de l'intensité de l'effort. La ligne pointillée (QR $= 1$) marque le seuil d'oxydation exclusive des glucides ; son dépassement traduit la production d'acide lactique.}
\end{popfigure}
\textbf{Points de vigilance :} un QR $> 1$ ne signifie pas « plus de respiration » mais un \ce{CO2} d'origine non respiratoire (tamponnement) ; citer l'équation du tampon bicarbonate.
\end{corrige}

\begin{corrige}[Exercice 7 --- Interprétation d'une courbe de variation pondérale]
\begin{enumerate}
\item Courbe de la masse en fonction du temps :
\end{enumerate}
\begin{popfigure}
\begin{center}
\begin{tikzpicture}
\begin{axis}[
width=0.8\linewidth, height=6cm,
xlabel={Semaines}, ylabel={Masse (kg)},
xmin=0, xmax=11, ymin=57, ymax=65,
xtick={0,1,2,3,4,5,6,7,8,9,10,11},
grid=both, grid style={popline!40},
]
\addplot[popcurve, mark=*, mark size=1.5pt] coordinates {(0,58) (1,58.8) (2,59.6) (3,60.4) (4,61.2) (5,62) (6,63) (7,64) (8,62.8) (9,61.8) (10,60.9) (11,61)};
\end{axis}
\end{tikzpicture}
\end{center}
\legende{Variation de la masse corporelle sur 12 semaines : phase de prise de masse (semaines 0 à 7, bilan positif) puis phase de perte (semaines 7 à 11, bilan négatif).}
\end{popfigure}
\begin{enumerate}
\setcounter{enumi}{1}
\item Phase 1 (semaines 0 à 7) : prise de masse, bilan énergétique \textbf{positif}. Phase 2 (semaines 7 à 11) : perte de masse, bilan énergétique \textbf{négatif}.
\item Prise de masse : $64 - 58 = \SI{6}{kg}$. Perte de masse : $64 - 61 = \SI{3}{kg}$.
\item La prise de masse correspond au stockage de l'excédent sous forme de triglycérides dans les adipocytes du tissu adipeux. La perte n'est jamais purement lipidique car l'organisme mobilise aussi du glycogène (stocké avec de l'eau) et une partie des protéines musculaires via la néoglucogenèse.
\item Lipides perdus : $0,70 \times 3\,000 = \SI{2100}{g}$, soit $2\,100 \times 38 = \SI{79800}{kJ}$ sur 28 jours, d'où un déficit moyen de $79\,800 / 28 = \SI{2850}{kJ/jour}$.
\end{enumerate}
\end{corrige}

\begin{corrige}[Exercice 8 --- Programme de rééquilibrage]
\begin{enumerate}
\item IMC $= 69 / 1,60^2 = 69 / 2,56 \approx 27,0$ : Sandrine est en situation de \textbf{surpoids} (25 $<$ IMC $<$ 30).
\item Dépenses totales : $5\,900 + 2\,100 = \SI{8000}{kJ/jour}$. Bilan $= 10\,500 - 8\,000 = +\SI{2500}{kJ/jour}$ (positif).
\item Objectif : $10\,500 - 1\,500 = \SI{9000}{kJ/jour}$. Répartition : glucides $0,50$ à $0,55 \times 9\,000 = 4\,500$ à \SI{4950}{kJ} ; lipides $0,30$ à $0,35 \times 9\,000 = 2\,700$ à \SI{3150}{kJ} ; protéines $0,12$ à $0,15 \times 9\,000 = 1\,080$ à \SI{1350}{kJ}.
\item Activité d'endurance d'intensité modérée : marche rapide, footing léger, danse ou football, 30 à 45 min, 3 à 5 fois par semaine. L'intensité modérée favorise l'oxydation des lipides (QR proche de 0,7--0,8) et reste compatible avec les études.
\item Trois indicateurs : la masse corporelle (pesée hebdomadaire, à heure fixe), le tour de taille, et le journal des apports alimentaires ; on peut ajouter la durée d'activité physique hebdomadaire. Objectif raisonnable : perte de 0,5 kg par semaine maximum, sans régime drastique.
\end{enumerate}
\textbf{Points de vigilance :} l'IMC se calcule avec la taille \emph{au carré} en mètres ; un rééquilibrage ne doit jamais passer sous le métabolisme de base.
\end{corrige}

\courssub{Je me prépare au Probatoire}

\begin{corrige}[Exercice 9 --- Sujet type Probatoire : le jeûne de 48 heures]
\begin{enumerate}
\item La glycémie est la concentration de glucose dans le sang ; sa valeur normale à jeun est d'environ \SI{1}{g/L} (0,8 à \SI{1,1}{g/L}).
\item La glycémie diminue modérément (0,95 puis 0,82 puis \SI{0,78}{g/L}) mais reste proche de la normale : elle ne s'effondre pas malgré l'absence d'apport. C'est l'\textbf{homéostasie glycémique}, régulation qui maintient constante une grandeur du milieu intérieur.
\item L'insuline (hormone hypoglycémiante) chute de 12 à 4 unités, tandis que le glucagon (hormone hyperglycémiante) passe de 8 à 20 unités. Sous l'action du glucagon, le foie dégrade son glycogène (\textbf{glycogénolyse}) et libère du glucose dans le sang, ce qui compense l'absence d'apport pendant les 24 premières heures.
\item Le glycogène hépatique étant épuisé, deux voies prennent le relais : la \textbf{lipolyse} dans le tissu adipeux (libération d'acides gras et de glycérol, oxydés par les muscles et le foie) et la \textbf{néoglucogenèse} hépatique à partir du glycérol, du lactate et des acides aminés (alanine) issus de la dégradation des protéines musculaires.
\item \ce{2 C3H7NO2 + 2 H2O -> C6H12O6 + 2 NH3}. L'ammoniac (\ce{NH3}) formé, très toxique, est transformé en urée par le foie au cours du \textbf{cycle de l'urée} (uréogenèse), puis éliminé par les reins.
\item Les corps cétoniques (acétoacétate, bêta-hydroxybutyrate) sont produits par le foie à partir de l'oxydation massive des acides gras (cétogenèse). Ils servent de carburant de substitution au cerveau, ce qui \textbf{économise le glucose} et donc les protéines musculaires qui seraient sinon dégradées pour la néoglucogenèse.
\end{enumerate}
\textbf{Barème indicatif (sur 10) :} Q1 : 1 pt ; Q2 : 1,5 pt ; Q3 : 2 pts ; Q4 : 2 pts ; Q5 : 1,5 pt ; Q6 : 2 pts.
\textbf{Points de vigilance :} citer les deux hormones avec leur sens d'action (insuline hypoglycémiante, glucagon hyperglycémiant) ; respecter l'ordre chronologique des voies (glycogénolyse puis lipolyse et néoglucogenèse) ; vérifier l'équilibre de l'équation bilan (azote : 2 de chaque côté).
\end{corrige}

\begin{corrige}[Exercice 10 --- Sujet type Probatoire : effort sportif et métabolisme au Mont Cameroun]
\begin{enumerate}
\item Athlète A : repos $0,22/0,28 \approx 0,79$ ; ascension $3,42/3,60 = 0,95$. Athlète B : repos $0,24/0,30 = 0,80$ ; ascension $3,10/2,90 \approx 1,07$.
\item Pendant l'ascension, le QR de A (0,95) est inférieur à celui de B (1,07) : A oxyde davantage les lipides. Les réserves lipidiques (\SI{300000}{kJ}) étant près de 40 fois supérieures aux réserves glucidiques (\SI{8000}{kJ}), l'utilisation des lipides \textbf{ménage le glycogène} et retarde l'épuisement : avantage décisif en endurance.
\item Chez B, l'apport d'\ce{O2} est insuffisant : ses muscles recourent à la fermentation lactique \ce{C6H12O6 -> 2 C3H6O3}. L'acide lactique libère des ions \ce{H+}, tamponnés par les bicarbonates sanguins (\ce{H+} + \ce{HCO3-} -> \ce{CO2} + \ce{H2O}) : le \ce{CO2} supplémentaire, non issu de la respiration, fait passer le QR au-dessus de 1.
\item C'est la \textbf{dette d'oxygène} (consommation d'\ce{O2} élevée après l'effort). Elle sert à réoxyder l'acide lactique accumulé (au foie, où il est reconverti en glucose par le cycle de Cori), à reconstituer les réserves de glycogène et d'ATP, et à restaurer les réserves d'\ce{O2} de la myoglobine. Elle est plus importante chez B (1,40 contre \SI{0,90}{L/min}), qui a davantage fermenté.
\item Volume d'\ce{O2} : $3,60 \times 180 = \SI{648}{L}$. Énergie : $648 \times 20 = \SI{12960}{kJ}$ (environ \SI{3100}{kcal}).
\item Adaptations de l'entraînement : augmentation du nombre et de la taille des mitochondries musculaires et des enzymes oxydatives (meilleure utilisation des lipides) ; hypertrophie cardiaque avec augmentation du volume d'éjection systolique et meilleure vascularisation capillaire des muscles, d'où un apport d'\ce{O2} plus efficace (consommation maximale d'\ce{O2} plus élevée : 3,60 contre \SI{2,90}{L/min}).
\end{enumerate}
\textbf{Barème indicatif (sur 10) :} Q1 : 2 pts ; Q2 : 2 pts ; Q3 : 2 pts ; Q4 : 1,5 pt ; Q5 : 1,5 pt ; Q6 : 1 pt.
\textbf{Points de vigilance :} convertir les 3 h en 180 min avant le calcul ; le QR $> 1$ exige la mention du tampon bicarbonate ; la dette d'oxygène doit être reliée à l'élimination de l'acide lactique, pas seulement nommée.
\end{corrige}

\begin{corrige}[Exercice 11 --- Sujet type Probatoire : dénutrition et obésité, deux déséquilibres]
\begin{enumerate}
\item Groupe 1 : $6\,200 - 8\,400 = -\SI{2200}{kJ/jour}$ (bilan négatif). Groupe 2 : $12\,800 - 9\,600 = +\SI{3200}{kJ/jour}$ (bilan positif).
\item Groupe 1 (IMC $= 16,8 < 18,5$) : maigreur (dénutrition). Groupe 2 (IMC $= 27,5$, entre 25 et 30) : surpoids, proche de l'obésité, confirmé par une masse grasse de 32 \%.
\item Conséquences du bilan négatif chronique : fonte musculaire (protéines dégradées pour la néoglucogenèse), retard de croissance et troubles hormonaux (retard pubertaire, aménorrhée), baisse des défenses immunitaires (sensibilité accrue aux infections). Ordre d'utilisation des réserves : le glycogène, réserve rapidement mobilisable mais faible, est épuisé en 24 h ; les lipides, réserve considérable, sont ensuite mobilisés ; les protéines, qui assurent des fonctions structurales et fonctionnelles vitales (enzymes, muscles), ne sont dégradées qu'en dernier recours car leur perte compromet la survie.
\item L'excédent énergétique (glucides et lipides consommés au-delà des besoins) est transformé en triglycérides stockés dans les adipocytes du tissu adipeux, dont la masse augmente. Complications de l'obésité : diabète de type 2 (insulinorésistance), hypertension artérielle, maladies cardiovasculaires (athérosclérose, infarctus), auxquelles s'ajoutent les troubles articulaires.
\item Perte visée : $\SI{500}{g} \times \SI{38}{kJ/g} = \SI{19000}{kJ/semaine}$, soit $19\,000 / 7 \approx \SI{2714}{kJ/jour}$ de déficit. Ce déficit est inférieur à l'excédent actuel (\SI{3200}{kJ/jour}) : il suffit de combiner une restriction modérée des apports (vers \SI{10000}{kJ/jour}) et une augmentation des dépenses par une activité physique quotidienne (marche, sport, travaux champêtres) pour atteindre l'objectif sans régime drastique.
\item Exemple de texte attendu : « La dénutrition comme l'obésité résultent d'un même déséquilibre entre l'énergie apportée par les aliments et l'énergie dépensée par l'organisme. Quand les apports sont insuffisants, le corps puise dans ses réserves : maigreur, fatigue, infections. Quand ils sont excessifs, l'excédent est stocké en graisses : surpoids, diabète, hypertension. Pour rester en bonne santé : manger varié et équilibré en privilégiant les produits locaux (légumes, poissons, tubercules en quantités adaptées), limiter les boissons sucrées et les fritures, et pratiquer une activité physique régulière (marche, sport, travaux champêtres). »
\end{enumerate}
\textbf{Barème indicatif (sur 10) :} Q1 : 1 pt ; Q2 : 1 pt ; Q3 : 2 pts ; Q4 : 2 pts ; Q5 : 2 pts ; Q6 : 2 pts.
\textbf{Points de vigilance :} justifier l'ordre de mobilisation des réserves (disponibilité, quantité, rôle vital des protéines) ; dans le texte de sensibilisation, exiger la relation explicite apports/dépenses et trois mesures concrètes adaptées au contexte local.
\end{corrige}

\newpage

\coursec{Fiche bilan}

\begin{popfigure}
\centering
\begin{tikzpicture}[
    mindmap,
    grow cyclic,
    every node/.style={concept, execute at begin node=\hskip0pt},
    concept color=popblue,
    level 1/.style={level distance=5.5cm, sibling angle=72},
    level 2/.style={level distance=3.2cm, sibling angle=45},
    texte/.style={text width=2.8cm, align=center, font=\small\bfseries},
    branch1/.style={concept color=poporange, texte},
    branch2/.style={concept color=popgreen, texte},
    branch3/.style={concept color=poppurple, texte},
    branch4/.style={concept color=poppink, texte},
    branch5/.style={concept color=popgold, texte},
]
\node[texte, text width=3.2cm, font=\large\bfseries] {Déséquilibres\\ Énergétiques\\ Organisme}
    child[branch1] { node[align=center]{1. Production\\ Énergie Cellulaire}
        child { node[concept color=poporange!60, texte, text width=2.2cm] {Respiration cellulaire\\ (aérobie, 38 ATP)} }
        child { node[concept color=poporange!60, texte, text width=2.2cm] {Fermentation\\ (anaérobie, 2 ATP)} }
    }
    child[branch2] { node[align=center]{2. Bilan\\ Énergétique}
        child { node[concept color=popgreen!60, texte, text width=2.2cm] {Apports : Glucides,\\ Lipides, Protéines} }
        child { node[concept color=popgreen!60, texte, text width=2.2cm] {Dépenses : MB, Effort,\\ Thermogenèse, SDA} }
        child { node[concept color=popgreen!60, texte, text width=2.2cm] {$\Delta E = Apports - Dépenses$} }
    }
    child[branch3] { node[align=center]{3. Conséquences\\ Déséquilibres}
        child { node[concept color=poppurple!60, texte, text width=2.2cm] {Positif (+) : Surpoids,\\ Obésité, Diabète T2,\\ AVC, HTA} }
        child { node[concept color=poppurple!60, texte, text width=2.2cm] {Négatif (-) : Maigreur,\\ Dénutrition, Amaigrissement,\\ Affaiblissement immunitaire} }
    }
    child[branch4] { node[align=center]{4. Adaptations\\ Métaboliques}
        child { node[concept color=poppink!60, texte, text width=2.2cm] {Effort : Glycogenolyse,\\ Lipolyse, Cétogenèse} }
        child { node[concept color=poppink!60, texte, text width=2.2cm] {Jeûne : Néo-glucogenèse,\\ Économie protéique,\\ Corps cétoniques} }
    }
    child[branch5] { node[align=center]{5. Prévention\\ \& Correction}
        child { node[concept color=popgold!60, texte, text width=2.2cm] {Alimentation équilibrée\\ (Repas type camerounais)} }
        child { node[concept color=popgold!60, texte, text width=2.2cm] {Activité physique\\ régulière} }
        child { node[concept color=popgold!60, texte, text width=2.2cm] {Suivi médical,\\ Éducation nutritionnelle} }
    };
\end{tikzpicture}
\legende{Carte mentale récapitulative de la leçon ND05 : Des voies de production d'ATP aux mesures de prévention des déséquilibres énergétiques.}
\end{popfigure}

\begin{bilan}
\textbf{1. Définitions clés}
\begin{itemize}
    \item \cle{Respiration cellulaire} : Oxydation complète du glucose (ou acides gras/acides aminés) en présence d'O$_2$ $\rightarrow$ CO$_2$ + H$_2$O + \textbf{$\approx$ 38 ATP} (glycolyse + cycle de Krebs + chaîne respiratoire mitochondriale).
    \item \cle{Fermentation} : Dégradation incomplète du glucose sans O$_2$ $\rightarrow$ Lactate (muscle) ou Éthanol + CO$_2$ (levures) + \textbf{2 ATP} (glycolyse seulement). Régénère le NAD$^+$.
    \item \cle{Métabolisme de base (MB)} : Énergie minimale pour les fonctions vitales au repos, à jeun, thermoneutralité ($\approx$ 60--70\% dépenses totales). Facteurs : âge, sexe, masse maigre, thyroïde, température.
    \item \cle{Bilan énergétique} : $\Delta E = \text{Apports (kJ/j)} - \text{Dépenses (kJ/j)}$. Équilibré si $\Delta E \approx 0$.
    \item \cle{Valeurs énergétiques} : Glucides = 17 kJ/g (4 kcal/g) ; Lipides = 38 kJ/g (9 kcal/g) ; Protéines = 17 kJ/g (4 kcal/g) ; Alcool = 29 kJ/g (7 kcal/g).
\end{itemize}

\textbf{2. Formules et repères numériques à connaître}
\begin{center}
\begin{tabularx}{\linewidth}{|l|X|X|}\hline
\rowcolor{popblueL}
\textbf{Grandeur} & \textbf{Formule / Valeur} & \textbf{Interprétation / Piège} \\ \hline
Bilan journalier & $\Delta E = \sum(\text{Apports}_i) - (\text{MB} + \text{Activité} + \text{SDA} + \text{Thermorégulation})$ & SDA = Action dynamique spécifique ($\approx$ 10\% apports). Activité = MB $\times$ coefficient (1,3--2,4). \\ \hline
RQ (Quotient Respiratoire) & $RQ = \frac{V_{\ce{CO2}}}{V_{\ce{O2}}}$ & Glucides : 1,0 ; Lipides : 0,7 ; Protéines : 0,8 ; Mixte : 0,8--0,85. Indicateur du substrat oxydé. \\ \hline
Besoin énergétique (Homme 70 kg) & $\approx$ 10--11 MJ/j (2400--2600 kcal/j) & Femme 60 kg : $\approx$ 8--9 MJ/j (1900--2100 kcal/j). Varie fort avec l'activité (agriculteur vs bureau). \\ \hline
Seuil obésité (IMC) & $\text{IMC} = \frac{m}{t^2} \ge 30$ kg/m$^2$ & Surpoids : 25--29,9. IMC insuffisant seul (masse musculaire, répartition graisseuse). \\ \hline
Perte de poids théorique & 1 kg graisse $\approx$ 37 MJ (9000 kcal) & Déficit 500 kcal/j $\rightarrow$ $\approx$ 0,5 kg/semaine. Perte rapide = eau + muscle (dangereux). \\ \hline
\end{tabularx}
\end{center}

\textbf{3. Méthodes express}
\begin{methode}[Établir un bilan énergétique quantitatif]
    \begin{enumerate}
        \item Calculer les \textbf{apports} : Peser aliments $\times$ composition (table Ciqual/ANC) $\rightarrow$ kJ.
        \item Estimer le \textbf{Métabolisme de Base} : Formule de Black et al. (OMS 1985) ou Harris-Benedict révisée (Mifflin-St Jeor).
        \item Estimer les \textbf{dépenses d'activité} : MB $\times$ NAP (Niveau d'Activité Physique : 1,4--2,0).
        \item Ajouter \textbf{SDA} (10\% apports) et \textbf{thermorégulation} (négligeable en zone neutre).
        \item Comparer : $\Delta E > 0$ (stockage), $\Delta E < 0$ (déstockage), $\Delta E \approx 0$ (équilibre).
    \end{enumerate}
\end{methode}

\begin{methode}[Interpréter une courbe de consommation O$_2$ / production CO$_2$ à l'effort]
    \begin{enumerate}
        \item Identifier le \textbf{seuil aérobie} (RQ $\approx$ 0,85--0,90) : transition lipides $\rightarrow$ glucides.
        \item Identifier le \cle{seuil anaérobie} (point de rupture ventilation, RQ $> 1,0$) : apparition lactatémie, hyperventilation compensatrice.
        \item \textbf{VO$_{2\text{max}}$} : Palier de la consommation O$_2$ malgré l'augmentation charge $\rightarrow$ capacité aérobie max.
        \item \textbf{Dette en O$_2$} : Surface entre courbe O$_2$ réelle et théorique au démarrage/arrêt $\rightarrow$ remboursement post-effort (récupération).
    \end{enumerate}
\end{methode}

\textbf{4. Pièges fréquents (Probatoire)}
\begin{attention}[Erreurs à ne pas commettre]
    \begin{itemize}
        \item Confondre \textbf{Rendement ATP} (théorique max 38) et \textbf{Rendement réel} ($\approx$ 30--32 ATP : coût transport NADH cytosolique, fuites protons, thermogenèse).
        \item Oublier que les \textbf{protéines ne sont pas stockées} : en jeûne prolongé, néoglucogenèse à partir d'acides aminés musculaires $\rightarrow$ perte masse maigre.
        \item Croire que \textbf{lipides $\rightarrow$ glucose} : Acétyl-CoA (beta-oxydation) ne peut pas faire du glucose (pas de glyoxylate chez l'homme). Seuls glycérol (3C) et certains AA sont glucoformateurs.
        \item Négliger la \cle{thermogenèse adaptative} : En suralimentation, MB augmente (futile cycles, protéines découplantes UCP) ; en restriction, MB baisse (économie) $\rightarrow$ plateau de perte de poids.
        \item Unités : \textbf{kJ} (SI) vs \textbf{kcal} (ancien). 1 kcal = 4,184 kJ. Ne pas mélanger dans un calcul.
    \end{itemize}
\end{attention}
\end{bilan}

\begin{aretenir}[Checklist avant le Probatoire]
\begin{itemize}
    \item $\square$ Écrire l'équation bilan de la respiration cellulaire et de la fermentation lactique/alcoolique avec rendements ATP.
    \item $\square$ Définir Métabolisme de Base, SDA, NAP ; citer les formules de calcul (Black/OMS ou Mifflin).
    \item $\square$ Calculer un bilan énergétique journalier à partir de données brutes (repas + activité) et conclure (équilibre / + / -).
    \item $\square$ Expliquer les conséquences physiologiques d'un bilan \textbf{positif chronique} : hypertrophie/hyperplasie adipocytes, résistance à l'insuline, syndrome métabolique, risques cardiovasculaires.
    \item $\square$ Expliquer les conséquences d'un bilan \textbf{négatif chronique} : amaigrissement, perte masse musculaire, troubles immunitaires, aménorrhée, ostéoporose, marasme/kwashiorkor (enfant).
    \item $\square$ Décrire les adaptations métaboliques à l'\textbf{effort court/intense} (anaérobie lactique, dette O$_2$) vs \textbf{effort long/modéré} (aérobie, lipolyse, épargne glycogène).
    \item $\square$ Décrire les adaptations au \textbf{jeûne prolongé} : phases (post-prandial $\rightarrow$ glycogène $\rightarrow$ néoglucogenèse + cétogenèse $\rightarrow$ économie protéique par corps cétoniques).
    \item $\square$ Interpréter un RQ (Quotient Respiratoire) et lier au substrat oxydé (glucides vs lipides).
    \item $\square$ Proposer des mesures correctives concrètes et adaptées au contexte camerounais (alimentation locale diversifiée, réduction huile/huile rouge/sucre, activité physique, dépistage diabète/HTA au CSI).
    \item $\square$ Différencier \textbf{obésité androïde} (viscérale, risque métabolique fort) et \textbf{gynoïde} (fessier-cuisses, risque mécanique) ; tour de taille critique (H $>$ 102 cm, F $>$ 88 cm).
    \item $\square$ Argumenter pourquoi les régimes "miracles" très hypocaloriques ($<$ MB) échouent à long terme (baisse MB, effet yoyo, troubles du comportement alimentaire).
\end{itemize}
\end{aretenir}

\findecours

\end{document}
